\chapter{Building-Block Matrix}\label{app:blocks}

This appendix contains the complete offering matrix of the building blocks --- the \texttt{All*} lists of the code --- referenced by the state of the art
(Section~\ref{sec:sota-instances}); the inventory enabled per build configuration is named separately by each section. Each axis is an \emph{organ} of the living being: a concrete search algorithm (technically \texttt{SearchAlgorithm}, metaphorically a living being) is the fixed composition of its eighteen organ axes (T0--T17), each of which corresponds to an organ; the building-block variants listed per axis are the interchangeable forms of that organ, and their permutation is, in the metaphor canon of this thesis, the exchange of organs on the living being. The \emph{anatomy} of this living being is the wiring \emph{between} the organs: axis algorithms use the interfaces of other axes (the allocation axis, for instance, provides a shared allocation interface for the remaining organs). For each axis it lists the purpose, the \emph{static} sub-axes
(compile-time-selectable variant families) and the \emph{dynamic} sub-axes
(run-time parameters), as well as the concrete building-block variants with their source (paper identifier or
code/baseline). Beyond the eighteen composition slots the organ registry carries an optional organ meta-meta axis\footnote{Meta-meta in the sense of the M3 level of the four-level layering of the Meta Object Facility~\cite{omg2016mof}: the indirection of a description. Meta axes (M2) classify main and sub-axes, M1 is the conception or the run-time implementation, M0 the binary code. The assignment of M3 is fixed; the cut between M2 and M1 is a reading of this thesis.} \texttt{disk\_\allowbreak{}io}, which forms no nineteenth composition slot and never enters the \texttt{binary\_\allowbreak{}id} (Section~\ref{sec:app-blocks-disk-io}). The entries are taken from the code of the \texttt{comdare-cache-engine} and verified against
the axis-algorithm-paper-code map of the \texttt{ce} repository (file \texttt{docs/\allowbreak{}architecture/\allowbreak{}18\_\allowbreak{}achsen\_\allowbreak{}algorithmus\_\allowbreak{}paper\_\allowbreak{}code\_\allowbreak{}map.md}).

Beyond the eighteen organ slots this appendix documents the two remaining axis categories of the triad (cf.\ Section~\ref{sec:axis-triad}). The \emph{system realm} comprises exactly three main axes in canonical, code-guarded order --- target ISA (\texttt{target\_\allowbreak{}isa}), operating system (\texttt{operating\_\allowbreak{}system}) and extension hardware (\texttt{external\_\allowbreak{}utils}) --- plus the outer complex bracket \texttt{build\_\allowbreak{}target\_\allowbreak{}complex}, which is explicitly not a fourth main axis; the order is not an editorial decision but follows the single ordering source of the code. The \emph{measurement realm} comprises the measurement categories, the measurement tooling and the load framework as meta-meta axis at the end of the realm. In the anatomy metaphor, system and measurement axes are neither a further organ nor the blood (the fluid organ that connects the organ axes), but the \emph{habitat} of the animal (system axes) and the equipment of the \emph{laboratory} (measurement axes; Section~\ref{sec:anatomy-metaphor}); both are not part of the permutation tree of the \texttt{binary\_\allowbreak{}id} (\texttt{binary\_\allowbreak{}id=never}) and strictly neutral towards the binary identity; their sub-axes are nevertheless permuted by the planner (\texttt{stage=runtime}). The entries of both spheres are additionally verified against the three axis registries of the code (\texttt{cache\_\allowbreak{}engine\_\allowbreak{}axis\_\allowbreak{}registry.xml} for the organ realm, \texttt{system\_\allowbreak{}axis\_\allowbreak{}registry.xml} and \texttt{measurement\_\allowbreak{}axis\_\allowbreak{}registry.xml}). Three conserved offering catalogues outside the build-driving offering --- the SIMD catalogue Build-SIMD under the extension hardware as well as the page-type catalogue Build-PG and the platform hardware strategy Build-HW --- are marked as such and collected in Section~\ref{sec:app-blocks-konserviert} (Conserved designs) at the end of this appendix.

\setlength{\tabcolsep}{3pt}
\section[head={T0 Search Algorithm Axis (Search Structures and Methods)}]{T0 Search Algorithm Axis T0 (Search Structures and Search Methods)}

This axis is the \emph{search organ} (T0) of the living being \texttt{SearchAlgorithm} (cf.\ Section~\ref{sec:sota-instances}).

Systematic permutation of 22 concrete search algorithms and search structures (S01--S22): 17 base building blocks --- array-based, sorted vectors, original paper wrappers, search methods like k-ary/\allowbreak{}interpolation, and ordered structures like skip-list/\allowbreak{}B-tree --- plus five opt-in building blocks disabled by default, namely the four per-K k-ary wrappers K2/\allowbreak{}K4/\allowbreak{}K8/\allowbreak{}K16 (S18--S21) and the Swiss-table hash structure (S22), to measurably quantify the impact of search paradigm choice on performance characteristics within a uniform \texttt{std::map} interface\@. \emph{Registry state:} the enabled inventory of this build configuration comprises four building blocks (\texttt{k\_\allowbreak{}ary}, \texttt{interpolation}, \texttt{eytzinger}, \texttt{linear\_\allowbreak{}scan}); the catalogue S01--S22 is the offering. Sub-axes classify by access density (SA1 dense, SA2 sparse, SA3/\allowbreak{}SA4 multilevel), enable compliance proofs through paper mixins (original-code proof, supervisor requirement: \texttt{is\_\allowbreak{}original} marking per function/\allowbreak{}wrapper), and identify SIMD-capable and store-traversable variants.

\subsection*{Static sub-axes}

\begin{scriptsize}
\begin{longtable}{@{}>{\raggedright\arraybackslash}p{1.8cm} >{\raggedright\arraybackslash}p{3.2cm} >{\raggedright\arraybackslash}p{8.5cm}@{}}
\toprule
ID & Name & Description \\
\midrule
\endhead%
SA1 & Direct-Addressed High Density (Dense Access) & Array family: direct-addressed (O(1) lookup), high fill density >= 50\%. Includes \texttt{Array256\allowbreak{}Search\allowbreak{}Algo} (256-slot compact array, ART Node256 re-impl) and \texttt{Array65535\allowbreak{}Search\allowbreak{}Algo} (uint16-based, mid-density 25-50\% direct-addressed). \\
SA2 & Sparse/\allowbreak{}Ordered (Sparse Access) & Sorted vectors, ordered search methods, and tree/\allowbreak{}hash structures: low fill density, range-scan support, variable O(log n)--O(1) lookup. Includes \texttt{Vector\allowbreak{}U8\allowbreak{}U8\allowbreak{}Search\allowbreak{}Algo}, original paper wrappers (ART/\allowbreak{}HOT/\allowbreak{}START/\allowbreak{}Wormhole/\allowbreak{}Su\allowbreak{}RF), search methods (k-ary, interpolation, Eytzinger, hash, Swiss table), and structures (skip-list, BST, B-tree). \\
SA3 & Multi-Byte Discriminators (Multilevel Access) & Multi-byte search structures with indexed sub-discriminators (e.g., START with cost-DP multi-level). \texttt{Vector\allowbreak{}U16\allowbreak{}U16\allowbreak{}Search\allowbreak{}Algo} (u16-keys, sorted lower\_\allowbreak{}bound) and \texttt{Original\allowbreak{}Start\allowbreak{}Search\allowbreak{}Algo} (publication P05) belong here. \\
SA4 & Direct-Addressed Medium Density (Direct Multibyte) & Shares the direct-addressing access pattern (like SA1) with multi-byte granularity (uint16 size) and medium fill density 25--50\% (vs.\ SA1 >= 50\%). Example: prt-art Array65535. This sub-axis captures the density variant without breaking the SA1 abstraction. \\
\bottomrule
\end{longtable}
\end{scriptsize}

\subsection*{Building blocks}

\begin{scriptsize}
\begin{longtable}{@{}>{\raggedright\arraybackslash}p{3.5cm} >{\raggedright\arraybackslash}p{6.5cm} >{\raggedright\arraybackslash}p{3.5cm}@{}}
\toprule
Variant & Description & Source \\
\midrule
\endhead%
Array256\allowbreak{}Search\allowbreak{}Algo (S01) & ART Node256 direct-addressed re-implementation: \texttt{std::array<std::optional<u64>, 256>} with O(1) lookup/\allowbreak{}insert/\allowbreak{}erase. Cache-line-aligned, SIMD-capable, no heap allocation. Baseline for dense accesses. & Code/\allowbreak{}Baseline (Leis ICDE 2013 ART, publication P01) \\
Vector\allowbreak{}U8\allowbreak{}U8\allowbreak{}Search\allowbreak{}Algo (S02) & Sparse sorted key-value vector pair: \texttt{std::vector<u8> keys} + \texttt{std::vector<u64> values} with \texttt{std::lower\_\allowbreak{}bound}. O(log n) lookup, O(n) insert/\allowbreak{}erase. Concept-equivalent to HOT-k-constrained. SIMD-capable (bitmask scan). & Code/\allowbreak{}Baseline (Binna SIGMOD 2018 HOT, publication P02) \\
Vector\allowbreak{}U16\allowbreak{}U16\allowbreak{}Search\allowbreak{}Algo (S03) & Sparse sorted u16 vector pair: \texttt{std::vector<u16> keys} + \texttt{std::vector<u64> values}. Simplified START concept (multi-byte discriminators without cost-DP). O(log n) lookup, O(n) insert/\allowbreak{}erase. & Code/\allowbreak{}Baseline (Fent ICDEW 2020 START, publication P05) \\
Array65535\allowbreak{}Search\allowbreak{}Algo (S09) & Mid-density direct-addressed uint16-based (25--50\% fill density). Migrated from prt-art. O(1) lookup but with memory overhead vs.\ SA1 denser arrays. & Code/\allowbreak{}Baseline (migrated from prt-art) \\
Original\allowbreak{}Art\allowbreak{}Search\allowbreak{}Algo (S04) & ART paper binding (\texttt{unodb::db}): \texttt{is\_\allowbreak{}original\_\allowbreak{}module()=true} (4/\allowbreak{}4 functions original). Paper mixin via SHA256-validated tool code generation. Original-code proof (supervisor requirement): \texttt{get\_\allowbreak{}compiler()='gcc-9.5'}, all insert/\allowbreak{}lookup/\allowbreak{}erase/\allowbreak{}clear marked original. & P01 (Leis/\allowbreak{}Kemper/\allowbreak{}Neumann ICDE 2013, Apache-2.0 unodb) \\
Original\allowbreak{}Hot\allowbreak{}Search\allowbreak{}Algo (S05) & HOT paper binding (HOTRowex): \texttt{is\_\allowbreak{}original\_\allowbreak{}module()=false} (2/\allowbreak{}4 original)\@. insert/\allowbreak{}lookup original, but erase/\allowbreak{}clear gaps (Hot\allowbreak{}Rowex is append-only; re-impl fallback). Publication P02. & P02 (Binna et al. SIGMOD 2018, ISC) \\
Original\allowbreak{}Start\allowbreak{}Search\allowbreak{}Algo (S06) & START paper binding (sosd adapter): \texttt{is\_\allowbreak{}original\_\allowbreak{}module()=false} (2/\allowbreak{}4 original)\@. insert/\allowbreak{}lookup original, erase/\allowbreak{}clear gaps (no remove in the sosd adapter). Publication P05, multi-byte discriminators. & P05 (Fent/\allowbreak{}Jungmair/\allowbreak{}Kipf/\allowbreak{}Neumann ICDEW 2020, MIT) \\
Original\allowbreak{}Wormhole\allowbreak{}Search\allowbreak{}Algo (S07) & Wormhole paper binding: \texttt{is\_\allowbreak{}original\_\allowbreak{}module()=false} (3/\allowbreak{}4 original)\@. put/\allowbreak{}get/\allowbreak{}del original, clear gap. Hybrid hash+trie+B+ index. GPL-3.0 licensed (blocks original-code linking for the thesis). & P07 (Wu/\allowbreak{}Ni/\allowbreak{}Jiang Euro\allowbreak{}Sys 2019, GPL-3.0) \\
Original\allowbreak{}Surf\allowbreak{}Search\allowbreak{}Algo (S08) & Su\allowbreak{}RF paper binding (FST with LOUDS bitmap): \texttt{is\_\allowbreak{}original\_\allowbreak{}module()=false} (1/\allowbreak{}4 original, lookup only). Succinct range filter, multi-header library. Publication P10. & P10 (Zhang/\allowbreak{}Lim/\allowbreak{}Andersen SIGMOD 2018, Apache-2.0) \\
KAry\allowbreak{}Search\allowbreak{}Algo (S10) & k-ary search, a search \emph{method} (Schlegel/\allowbreak{}Gemulla/\allowbreak{}Lehner Da\allowbreak{}Mo\allowbreak{}N 2009): K-way partition instead of halving in binary search. K is iterable\_\allowbreak{}aspect (K $\in$ \{2,4,8,16\}). O($\lceil$log\_\allowbreak{}\{K+1\} n$\rceil$) iterations, parallel ILP comparisons. Array family but no faithful dedicated traversal organ (SA2 classification). C++23 re-impl, \texttt{is\_\allowbreak{}original=false}. & Schlegel et al.\ Da\allowbreak{}Mo\allowbreak{}N 2009~\cite{schlegel2009kary} \\
Interpolation\allowbreak{}Search\allowbreak{}Algo (S11) & Interpolation search, a search \emph{method} (Perl/\allowbreak{}Itai/\allowbreak{}Avni CACM 1978): distribution-aware position estimation instead of blind halving. O(log log N) avg with uniform distribution, O(N) worst-case. Not SIMD-vectorizable (data-dependent index). C++23 re-impl. & Perl/\allowbreak{}Itai/\allowbreak{}Avni CACM 1978~\cite{perl1978interpolation} \\
Eytzinger\allowbreak{}Search\allowbreak{}Algo (S12) & Eytzinger/\allowbreak{}BFS-layout branch-free search (Khuong/\allowbreak{}Morin JEA 2017): cache-layout search \emph{method} with branch-free kernel. Keys in BFS order of the implicit search tree. For large arrays, per the publication, the fastest general-purpose layout. C++23 re-impl (experiment harness without standard licence). & Khuong/\allowbreak{}Morin JEA 2017~\cite{khuong2017eytzinger} \\
Skip\allowbreak{}List\allowbreak{}Search\allowbreak{}Algo (S13) & Skip-list, a probabilistic ordered \emph{structure} (Pugh CACM 1990): O(log n) expected search/\allowbreak{}insert/\allowbreak{}erase without explicit rebalancing. Each node has 1..k\allowbreak{}Max\allowbreak{}Level forward pointers (coin-flip probability 0.5). Index-based (\texttt{std::vector<Node>} with uint32 indices, value-semantic). Textbook algorithm, C++23 re-impl. & Pugh CACM 1990~\cite{pugh1990skiplist} \\
Hash\allowbreak{}Search\allowbreak{}Algo (S14) & Open-addressing hash table, \emph{unordered} (Knuth TAOCP 3 \S6.4): Fibonacci hash, linear probing, tombstone-erase on resize. O(1) expected lookup/\allowbreak{}insert at load\_\allowbreak{}factor < 0.7. No range-scan support (distinct classification). Correct tombstone marking for probe-chain integrity. C++23 re-impl. & Knuth TAOCP 3 \S6.4~\cite{knuth1998taocp3} \\
Linear\allowbreak{}Scan\allowbreak{}Search\allowbreak{}Algo (S15) & Unsorted linear scan baseline (Leis ICDE 2013 ART Node4): O(n) lookup on an unsorted KV array. Simplest comparison-zero-point structure: insert O(1)-amortized, erase O(1) via swap-and-pop. Store-traversable (array family). No range-scan support. & Code/\allowbreak{}Baseline (Leis ICDE 2013 ART, publication P01, Node4) \\
Binary\allowbreak{}Search\allowbreak{}Tree\allowbreak{}Search\allowbreak{}Algo (S16) & Unbalanced binary search tree (Knuth TAOCP 3 \S6.2.2 Hibbard deletion): deterministic comparison-tree baseline (vs.\ skip-list probabilistic). Index-based (\texttt{std::vector<Node>} with uint32 indices, free-list). Hibbard erase (3 cases: leaf/\allowbreak{}1 child/\allowbreak{}2 children via in-order successor). O(n) worst-case with sorted input. & Knuth TAOCP 3 \S6.2.2~\cite{knuth1998taocp3} \\
BTree\allowbreak{}Search\allowbreak{}Algo (S17) & Balanced block-oriented multi-way B-tree (Bayer/\allowbreak{}Mc\allowbreak{}Creight Acta Inf. 1972 /\allowbreak{} CLRS Ch. 18): t=4, hence [3..7] keys per node, up to 8 children. Always balanced (O(log n) guaranteed). Index-based, \texttt{alignas(64)} nodes (cache-line alignment). Top-down splitting on insert, borrow/\allowbreak{}merge on erase. The standard structure for database indexes. & Bayer/\allowbreak{}Mc\allowbreak{}Creight Acta Inf.\ 1972 /\allowbreak{} CLRS~\cite{bayer1972btree} \\
\bottomrule
\end{longtable}
\end{scriptsize}

\subsection*{Building blocks (opt-in, disabled by default)}

The following five building blocks are registered in the catalogue and conformant, but disabled by default in the build configuration; the registration builds nothing, it makes the variants selectable.

\begin{scriptsize}
\begin{longtable}{@{}>{\raggedright\arraybackslash}p{3.5cm} >{\raggedright\arraybackslash}p{6.5cm} >{\raggedright\arraybackslash}p{3.5cm}@{}}
\toprule
Variant & Description & Source \\
\midrule
\endhead%
KAry\allowbreak{}Search\allowbreak{}Algo\allowbreak{}K2 (S18) & Compile-time K wrapper of the k-ary search with K=2 (binary-search baseline): one animal binary per K (\texttt{KAry\allowbreak{}Search\allowbreak{}Algo\allowbreak{}K<N>} $\rightarrow$ \texttt{KAry\allowbreak{}Traversal<K>}) with its own enable flag (default OFF) and its own name \texttt{k\_\allowbreak{}ary\_\allowbreak{}k2} (separation in the \texttt{binary\_\allowbreak{}id}). & Schlegel et al.\ Da\allowbreak{}Mo\allowbreak{}N 2009~\cite{schlegel2009kary} (Code/\allowbreak{}Baseline) \\
KAry\allowbreak{}Search\allowbreak{}Algo\allowbreak{}K4 (S19) & K=4 (paper default, 5-way partition); name \texttt{k\_\allowbreak{}ary\_\allowbreak{}k4}, otherwise as S18. & Schlegel et al.\ Da\allowbreak{}Mo\allowbreak{}N 2009~\cite{schlegel2009kary} (Code/\allowbreak{}Baseline) \\
KAry\allowbreak{}Search\allowbreak{}Algo\allowbreak{}K8 (S20) & K=8; name \texttt{k\_\allowbreak{}ary\_\allowbreak{}k8}, otherwise as S18. & Schlegel et al.\ Da\allowbreak{}Mo\allowbreak{}N 2009~\cite{schlegel2009kary} (Code/\allowbreak{}Baseline) \\
KAry\allowbreak{}Search\allowbreak{}Algo\allowbreak{}K16 (S21) & K=16; name \texttt{k\_\allowbreak{}ary\_\allowbreak{}k16}, otherwise as S18. & Schlegel et al.\ Da\allowbreak{}Mo\allowbreak{}N 2009~\cite{schlegel2009kary} (Code/\allowbreak{}Baseline) \\
Swiss\allowbreak{}Table\allowbreak{}Search\allowbreak{}Algo (S22) & CE-native re-implementation of the Swiss-table principle: open addressing with one control byte per slot (empty, deleted, 7-bit H2 fingerprint), groups of 16 slots, H1/\allowbreak{}H2 split and tombstones for erase; unordered, no range scan (SA2). Allocation exclusively through the allocator axis (four buffers in one allocation ledger). Disabled by default and, without organ backing, not to be read as a Swiss-table measurement; \texttt{is\_\allowbreak{}original}=false. & Google Abseil Swiss Tables~\cite{abseil_swisstable} (talk CppCon 2017); C++23 re-impl \\
\bottomrule
\end{longtable}
\end{scriptsize}

\emph{Contributing sources:} P01, P02, P05, P07, P10, k-ary (Schlegel et al.\ 2009), interpolation (Perl et al.\ 1978), Eytzinger (Khuong/\allowbreak{}Morin 2017), skip list (Pugh 1990), hashing/\allowbreak{}BST (Knuth TAo\allowbreak{}CP 3), B-tree (Bayer/\allowbreak{}Mc\allowbreak{}Creight 1972), Swiss Tables (Google/\allowbreak{}Abseil 2018)~\cite{leis2013art,binna2018hot,fent2020start,wu2019wormhole,zhang2018surf,schlegel2009kary,perl1978interpolation,khuong2017eytzinger,pugh1990skiplist,knuth1998taocp3,bayer1972btree,abseil_swisstable}.

\section{T1 Cache Traversal (Fanout Lookup)}

This axis is the \emph{cache-traversal organ} (T1) of the living being \texttt{SearchAlgorithm} (cf.\ Section~\ref{sec:sota-instances}).

Axis \texttt{axis\_\allowbreak{}03b} defines three distinct strategies for key resolution in small fanout arrays (typically B+ tree or trie inner nodes). It spans a spectrum from linear scan (O(N), cache-friendly, small), via multiplicative hashing with open addressing (O(1) amortized, Knuth-Fibonacci), to sorted binary search (O(log N), ordered traversal, B+ tree standard). All three are textbook baselines without original code linking, implemented as pure C++23 re-implementations.

\subsection*{Static sub-axes}
\begin{scriptsize}
\begin{longtable}{@{}>{\raggedright\arraybackslash}p{1.8cm} >{\raggedright\arraybackslash}p{3.2cm} >{\raggedright\arraybackslash}p{8.5cm}@{}}
\toprule
ID & Name & Description \\
\midrule
\endhead%
CT1 & Linear Access & Linear-scan-based lookup methods on small, unsorted fanout arrays. Characterized by O(N) resolution with simple comparisons, typical for nodes with fanout $<$ 32 in B+ trees. Sub-axis for \texttt{Linear\allowbreak{}Fanout} and \texttt{Binary\allowbreak{}Search\allowbreak{}Fanout} (both: \texttt{axis\_\allowbreak{}tag} = \texttt{linear\_\allowbreak{}access\_\allowbreak{}tag}). \\
CT2 & Hash-Based Access & Hash lookup with open addressing and linear probing. O(1) amortized resolution, characterized by Fibonacci multiplier (11400714819323198485 $\approx$ $2^{64}/\allowbreak{}\varphi$), power-of-2 bucket count, dynamic rehashing when load factor $>$ 0.7. Only \texttt{Hash\allowbreak{}Lookup} wrapper. Sub-axis: \texttt{axis\_\allowbreak{}tag} = \texttt{hash\_\allowbreak{}access\_\allowbreak{}tag}. \\
\bottomrule
\end{longtable}
\end{scriptsize}

\subsection*{Building blocks}
\begin{scriptsize}
\begin{longtable}{@{}>{\raggedright\arraybackslash}p{2.5cm} >{\raggedright\arraybackslash}p{6.5cm} >{\raggedright\arraybackslash}p{4.5cm}@{}}
\toprule
Variant & Description & Source \\
\midrule
\endhead%
\texttt{Linear\allowbreak{}Fanout} & Unsorted fanout traversal via \texttt{std::find\_\allowbreak{}if} and linear scan. O(N) lookup, no hash overhead, cache-friendly for small N ($<$ 32). Classic B-tree baseline (Bayer/\allowbreak{}Mc\allowbreak{}Creight 1972). Belongs to sub-axis CT1. & Code/\allowbreak{}Baseline --- Bayer/\allowbreak{}Mc\allowbreak{}Creight 1972 (Acta Informatica) \\
\texttt{Hash\allowbreak{}Lookup} & Fibonacci-hash open addressing with linear probing (Knuth TAOCP Vol. 3, §6.4). O(1) amortized, dynamic rehashing (doubles capacity when load $\geq$ 0.7), iterable runtime capacity (\texttt{k\allowbreak{}Iterable\allowbreak{}Initial\allowbreak{}Capacities}: \{8, 16, 64, 256, 1024\}). Standalone C++23 re-implementation, no original code binding. Belongs to sub-axis CT2. & Code/\allowbreak{}Baseline --- Knuth 1998 (The Art of Computer Programming Vol.3) \\
\texttt{Binary\allowbreak{}Search\allowbreak{}Fanout} & Sorted fanout traversal via \texttt{std::lower\_\allowbreak{}bound}. O(log N) lookup with sorted ordering, preserves key sequencing for range-capable routing layers. Classic B+ tree inner-node search (Bayer/\allowbreak{}Mc\allowbreak{}Creight 1972), third distinct strategy of the axis. C++23 re-implementation, \texttt{is\_\allowbreak{}original}=false. Belongs to sub-axis CT1. & Code/\allowbreak{}Baseline --- Bayer/\allowbreak{}Mc\allowbreak{}Creight 1972 (Acta Informatica) \\
\bottomrule
\end{longtable}
\end{scriptsize}

\emph{Contributing sources:} Bayer/\allowbreak{}Mc\allowbreak{}Creight 1972, Knuth TAo\allowbreak{}CP Vol.\,3~\cite{bayer1972btree,knuth1998taocp3}.

\emph{Related literature:} P17, P21, P32~\cite{bender2005coboblivious,chen2001prefetch,schmidt2025hbm}.

\section{T2 Slot-to-Offset Mapping (Traversal Sublayer)}

This axis is the \emph{mapping organ} (T2) of the living being \texttt{SearchAlgorithm} (cf.\ Section~\ref{sec:sota-instances}).

Mediation between logical slot indices and physical memory offsets in trie inner nodes and memory pools. The axis provides two fundamentally different mapping strategies: direct absolute offsets (linear packing) or pool-relative offsets with automatic rebase mechanism supporting memory reallocation and persistent data structures.

\subsection*{Static sub-axes}
\begin{scriptsize}
\begin{longtable}{@{}>{\raggedright\arraybackslash}p{1.8cm} >{\raggedright\arraybackslash}p{3.2cm} >{\raggedright\arraybackslash}p{8.5cm}@{}}
\toprule
ID & Name & Description \\
\midrule
\endhead%
MP1 & Direct-Access Mapping (MP1) & Direct slot-to-absolute-offset mapping without indirection. Wrapper concept for classic array-of-pointers layout as in B-tree inner nodes (Bayer/\allowbreak{}Mc\allowbreak{}Creight 1972). Simple cache locality, O(N) lookup via \texttt{std::find\_\allowbreak{}if}. \\
MP2 & Pool-Relative Mapping (MP2) & Slot-to-pool-relative-offset mapping using based-pointer arena pattern. All offsets stored relative to dynamic \texttt{pool\_\allowbreak{}base\_\allowbreak{}address}. Enables O(1) rebase on memory reallocation --- core property for persistent data structures (Driscoll et al. 1989) and Custom\allowbreak{}Aligned\allowbreak{}Structure pattern in prt-art. \\
\bottomrule
\end{longtable}
\end{scriptsize}

\subsection*{Building blocks}
\begin{scriptsize}
\begin{longtable}{@{}>{\raggedright\arraybackslash}p{2.6cm} >{\raggedright\arraybackslash}p{8.5cm} >{\raggedright\arraybackslash}p{2.4cm}@{}}
\toprule
Variant & Description & Source \\
\midrule
\endhead%
Direct\allowbreak{}Placement & Wrapper for direct slot-to-absolute-offset mapping (linear packing). \texttt{std::vector<pair<slot\_\allowbreak{}index\_\allowbreak{}type, offset\_\allowbreak{}type>>}; \texttt{resolve\_\allowbreak{}offset}/\allowbreak{}\texttt{register\_\allowbreak{}slot} via linear \texttt{std::find\_\allowbreak{}if}. Usage: prt-art \texttt{INode::placement\_\allowbreak{}page\_\allowbreak{}} + mapping table. Storage: 16-bit slot index + \texttt{size\_\allowbreak{}t} offset per entry. No pool base dependency. & CE-Baseline \\
Pool\allowbreak{}Relative & Wrapper for pool-relative-offset mapping with O(1) rebase capability. \texttt{std::vector<pair<slot\_\allowbreak{}index\_\allowbreak{}type, relative\_\allowbreak{}offset>>} + mutable \texttt{pool\_\allowbreak{}base\_\allowbreak{}address}. On pool reallocation: \texttt{rebase(new\_\allowbreak{}pool\_\allowbreak{}base)} updates only \texttt{pool\_\allowbreak{}base\_\allowbreak{}} --- all relative offsets remain valid. Used with prt-art Custom\allowbreak{}Aligned\allowbreak{}Structure + versioning scenarios (persistent data structures per Driscoll/\allowbreak{}Sarnak/\allowbreak{}Sleator/\allowbreak{}Tarjan 1989). & CE-Baseline \\
\bottomrule
\end{longtable}
\end{scriptsize}

\emph{Contributing sources:} --- (CE baseline, own wrappers).

\emph{Related literature:} Bayer/\allowbreak{}Mc\allowbreak{}Creight 1972, Driscoll et al.\ 1989, P03, P07~\cite{bayer1972btree,driscoll1989persistent,mao2012masstree,wu2019wormhole}.

\section{T3 Path Compression}

This axis is the \emph{path-compression organ} (T3) of the living being \texttt{SearchAlgorithm} (cf.\ Section~\ref{sec:sota-instances}).

Defines the compression granularity and strategy for key paths in trie structures. Controls how inner nodes store and navigate redundant byte sequences or significant bits of common prefixes to optimize memory and cache efficiency.

\subsection*{Static sub-axes}

\begin{scriptsize}
\begin{longtable}{@{}>{\raggedright\arraybackslash}p{1.8cm} >{\raggedright\arraybackslash}p{3.2cm} >{\raggedright\arraybackslash}p{8.5cm}@{}}
\toprule
ID & Name & Description \\
\midrule
\endhead%
PC1 & Compression Granularity & Defines the minimal unit of path compression: bit-wise (Patricia), byte-wise (ART), or no compression (raw path). Determines trie height and optimal code blocks for traversal. Code token: \texttt{granularity\_\allowbreak{}tag}. \\
PC2 & Skip Strategy & Defines the navigation pattern at inner nodes: linear over 8 bytes (Byte\allowbreak{}Wise), single-bit comparison (Patricia), or direct jump without intermediate steps. Influences branch density and prefetch behaviour. Code token: \texttt{skip\_\allowbreak{}strategy\_\allowbreak{}tag}. \\
PC3 & Decode Complexity & Classifies the runtime complexity when extracting prefix information: O(1) for direct extraction, O(log n) for structure-based search, O(n) for sequential check. Important for operations like \texttt{common\_\allowbreak{}prefix\_\allowbreak{}len()} and key comparisons. Code token: \texttt{decode\_\allowbreak{}complexity\_\allowbreak{}tag}. \\
\bottomrule
\end{longtable}
\end{scriptsize}

\subsection*{Building blocks}

\begin{scriptsize}
\begin{longtable}{@{}>{\raggedright\arraybackslash}p{3.0cm} >{\raggedright\arraybackslash}p{8.5cm} >{\raggedright\arraybackslash}p{2.0cm}@{}}
\toprule
Variant & Description & Source \\
\midrule
\endhead%
\texttt{Path\allowbreak{}Compression\allowbreak{}None} & No path compression; raw key paths are stored completely and uncompressed in nodes. Serves as measurement baseline with \texttt{compression\_\allowbreak{}ratio=1.0}. Used as control case to quantify memory and height savings of compressing strategies. & Code/\allowbreak{}Baseline \\
\texttt{Byte\allowbreak{}Wise\allowbreak{}Path\allowbreak{}Compression} & Byte-by-byte compression (ART-style); stores up to 7 prefix bytes per node in a compact \texttt{u64} word with length tag in high byte. Typical 2x reduction in trie height (\texttt{compression\_\allowbreak{}ratio=0.5}). Origin: Leis et al., ART ICDE 2013 (P01). Employs goldstandard operations \texttt{common\_\allowbreak{}prefix\_\allowbreak{}len()} and \texttt{cut()} for navigation optimization. & P01 \\
\texttt{Patricia\allowbreak{}Path\allowbreak{}Compression} & Single-bit-split path compression (Patricia trie); navigates bit-wise instead of byte-wise, reduces trie height by typically factor 3 (\texttt{compression\_\allowbreak{}ratio=0.3}). Origin: Morrison PATRICIA JACM 1968; modernly implemented in HOT (Binna SIGMOD 2018, P02) and Wormhole (Wu Euro\allowbreak{}Sys 2019, P07). F15-operative primitives: \texttt{key\_\allowbreak{}split\_\allowbreak{}bit()} for MSB-first bit extraction, \texttt{path\_\allowbreak{}descend\_\allowbreak{}scan()} for goldstandard-compliant \texttt{\_\allowbreak{}scan} signature with key-dependent variable depth. & P02 \\
\bottomrule
\end{longtable}
\end{scriptsize}

\emph{Contributing sources:} P01, P02, P07, Morrison 1968~\cite{leis2013art,binna2018hot,wu2019wormhole,morrison1968patricia}.

\emph{Related literature:} P04, P05, P09~\cite{boffa2024coco,fent2020start,jacobson1989louds}.

\section{T4 Node Types (Node-Type Axis)}

This axis is the \emph{node-type organ} (T4) of the living being \texttt{SearchAlgorithm} (cf.\ Section~\ref{sec:sota-instances}).

The axis \texttt{axis\_\allowbreak{}04\_\allowbreak{}node\_\allowbreak{}type} defines the four node formats of the Adaptive Radix Tree (ART) as exchangeable variants: Node4, Node16, Node48, and Node256. Each node type offers a different balance between memory, cache efficiency, and lookup speed. This axis forms the foundation for adaptive tree reorganization and is central to the design of search algorithms.

\subsection*{Static sub-axes}

\begin{scriptsize}
\begin{longtable}{@{}>{\raggedright\arraybackslash}p{1.8cm} >{\raggedright\arraybackslash}p{3.2cm} >{\raggedright\arraybackslash}p{8.5cm}@{}}
\toprule
ID & Name & Description \\
\midrule
\endhead%
NT1 & Capacity Class (NT1) & Classification by fanout size and slot count: small (Node4 with 4 slots, linear search), medium (Node16 with 16 slots, SIMD binary search) to maximal (Node256 with 256 slots, direct-addressed). \\
NT2 & Access Pattern (NT2) & Method of key lookup within a node: linear scan (Node4), SIMD binary search (Node16), indirect indexing (Node48), or direct addressing (Node256). \\
NT3 & Compactness (NT3) & Memory layout and density: sparse and cache-aligned (Node4), dense with indirection (Node48), or maximally direct (Node256), with impacts on memory usage and cache-line misses. \\
\bottomrule
\end{longtable}
\end{scriptsize}

\subsection*{Building blocks}

The table lists three kinds of entries from the same axis directory; only the first block is the value set of the axis.

\begin{scriptsize}
\begin{longtable}{@{}>{\raggedright\arraybackslash}p{3.0cm} >{\raggedright\arraybackslash}p{8.5cm} >{\raggedright\arraybackslash}p{2.0cm}@{}}
\toprule
Variant & Description & Source \\
\midrule
\endhead%
\multicolumn{3}{@{}l}{\emph{Axis variants (\texttt{All\allowbreak{}Node\allowbreak{}Types}, four)}} \\
\texttt{Node4\allowbreak{}Node\allowbreak{}Type} & Small ART node with 4 slots and linear search. Ideal for sparsely populated upper tree levels, fits completely in a cache line (64 bytes). Implemented as CE re-implementation of the ART specification with \texttt{node\_\allowbreak{}find\_\allowbreak{}scan()} as linear scan over sorted keys. & P01 \\
\texttt{Node16\allowbreak{}Node\allowbreak{}Type} & Medium ART node with 16 slots, optimized for SIMD-based binary search (SSE2 PCMPESTRM /\allowbreak{} AVX2 VPCMPESTRM). Balance between memory and search speed; largest capacity without indirection. CE re-implementation of the Leis ICDE 2013 specification. & P01 \\
\texttt{Node48\allowbreak{}Node\allowbreak{}Type} & Large ART node with 48 slots, using indirection: 256-byte Child\allowbreak{}Index (byte→slot mapping, 0=empty) plus 48-slot Child\allowbreak{}Array. Trade-off between memory (vs Node256 direct) and cache misses (vs Node4 linear). Double-indirection access pattern (index + child). & P01 \\
\texttt{Node256\allowbreak{}Node\allowbreak{}Type} & Maximal ART node with 256 slots, directly addressed without index. O(1) lookup (no scan, no index touch). High memory overhead (256 pointers), but optimal search speed for dense nodes. Default for lower tree levels. & P01 \\
\midrule
\multicolumn{3}{@{}l}{\emph{Paper node formats (reconstruction targets of this thesis, no CE baseline)}} \\
Internal/\allowbreak{}Border & Node-type pair of Masstree (trie of B\textsuperscript{+} trees): internal nodes navigate the 8-byte key slices of the B\textsuperscript{+} level, border nodes form the leaf level with the values or the transition to the next trie layer. Paper node format (reconstruction target of this thesis, no CE baseline). & P03 (Mao et al.\ 2012)~\cite{mao2012masstree} \\
Decision/\allowbreak{}Span & Node-type pair of the B\textsuperscript{2}-tree (B\textsuperscript{+} with embedded binary inner nodes, cache-aware): decision nodes carry the binary branching decision within the page, span nodes group contiguous key ranges. Paper node format (reconstruction target of this thesis, no CE baseline). & P06 (Schmeisser et al.\ 2022)~\cite{schmeisser2022b2tree} \\
LOUDS-Dense/\allowbreak{}Sparse & Node-type pair of the succinct range filter Su\allowbreak{}RF\@: LOUDS-Dense encodes the upper, densely populated trie levels bitmap-based with direct addressing, LOUDS-Sparse the lower, sparsely populated levels byte-sequentially succinct. Paper node format (reconstruction target of this thesis, no CE baseline). & P10 (Zhang et al.\ 2018)~\cite{zhang2018surf} \\
\midrule
\multicolumn{3}{@{}l}{\emph{Storage organs of the composition (storage organs, no \texttt{node\_\allowbreak{}type} values)}} \\
\texttt{Node\allowbreak{}Type\allowbreak{}Slot\allowbreak{}Store} & Storage organ: node-type-driven bounded slot storage. Spans \texttt{std::array} with capacity \texttt{N::max\_\allowbreak{}capacity()} over a shared uint64 key. Fulfills the Storage\allowbreak{}Organ concept and is drop-in substrate for Composed\allowbreak{}Search. Baseline variant without layout/\allowbreak{}allocator integration. & Code/\allowbreak{}Baseline \\
\texttt{Composed\allowbreak{}Store} & 3-axis storage organ: node\_\allowbreak{}type(N) $\oplus$ layout(L) $\oplus$ allocator(A). Unbounded vector over allocator A\@; layout flows in as static cache\_\allowbreak{}line\_\allowbreak{}size and via optional \texttt{organ\_\allowbreak{}scan\_\allowbreak{}field\_\allowbreak{}sum()}. Measurable (requirement F15). & Code/\allowbreak{}Baseline \\
\texttt{Node\allowbreak{}Chunked\allowbreak{}Store} & 3-axis storage organ with runtime impact of the node\_\allowbreak{}type axis: slots lie in chunks of capacity \texttt{N::max\_\allowbreak{}capacity()}. Chunk count = number of node-sized allocations; thus the node\_\allowbreak{}type axis is measurable (Node4 many chunks, Node256 few). Proof of delegation. & Code/\allowbreak{}Baseline \\
\texttt{Layout\allowbreak{}Aware\allowbreak{}Chunked\allowbreak{}Store} & Layout-honoring variant of Node\allowbreak{}Chunked\allowbreak{}Store: stores records at layout-driven stride (\texttt{eff\_\allowbreak{}stride}). Each record occupies \texttt{eff\_\allowbreak{}stride} bytes (key 0..8, value 8..16, rest padding). Makes the memory\_\allowbreak{}layout axis measurable, OOB-safe for scan methods. & Code/\allowbreak{}Baseline \\
\texttt{Observable\allowbreak{}Node\allowbreak{}Type} & Observable\allowbreak{}Axis wrapper around a node\_\allowbreak{}type strategy. Passes through static API (max\_\allowbreak{}capacity/\allowbreak{}name/\allowbreak{}topic\_\allowbreak{}tag) + offers \texttt{observe\_\allowbreak{}node\_\allowbreak{}find()} as instance driver. Tracks find\_\allowbreak{}count, keys\_\allowbreak{}stored, queries\_\allowbreak{}run, last\_\allowbreak{}checksum. Format-divergent lookup latency measurable (out-of-process path). & Code/\allowbreak{}Baseline \\
\bottomrule
\end{longtable}
\end{scriptsize}

\emph{Contributing sources:} P01~\cite{leis2013art}, P03~\cite{mao2012masstree}, P06~\cite{schmeisser2022b2tree}, P10~\cite{zhang2018surf}.

\section{T5 Memory Layout}

This axis is the \emph{memory-layout organ} (T5) of the living being \texttt{SearchAlgorithm} (cf.\ Section~\ref{sec:sota-instances}).

This axis classifies the spatial organization of records in memory, directly impacting cache efficiency, SIMD vectorization potential, and false-sharing characteristics. It captures four orthogonal design dimensions, three of which exist as sub-axis tags in the code (HM1--HM3): (1) alignment strategies (cache-line padding vs.\ strict packing), (2) data organization (row-oriented Ao\allowbreak{}S vs.\ column-oriented So\allowbreak{}A vs.\ hybrid Ao\allowbreak{}So\allowbreak{}A), (3) packing density (dense vs.\ succinct bit-packed); the fourth, the stride pattern (HM4, sequential vs.\ interleaved), is a planned dimension without a carrying layout building block, for which the code carries no tag (a permutable stride value would break the \texttt{All\allowbreak{}Layouts} list of the 320 catalogue, which is fixed at five layout building blocks). Layout choice is a primary driver of cache locality and SIMD capability during index traversal.

\subsection*{Static sub-axes}

\begin{scriptsize}
\begin{longtable}{@{}>{\raggedright\arraybackslash}p{1.8cm} >{\raggedright\arraybackslash}p{3.2cm} >{\raggedright\arraybackslash}p{8.5cm}@{}}
\toprule
ID & Name & Description \\
\midrule
\endhead%
HM1 & Alignment Strategy & Choice between cache-line-padded layout (false-sharing avoidance) and strict packing (memory density). Determines stride overhead during sequential traversal. \\
HM2 & Data Organization & Primary memory ordering: Array-of-Structs (Ao\allowbreak{}S, row-oriented), Struct-of-Arrays (So\allowbreak{}A, column-oriented), or hybrid Ao\allowbreak{}So\allowbreak{}A (block-mixed). Critical for cache-line utilization during field-selective scans. \\
HM3 & Packing Density & Storage density level: from standard records (8--64 bytes) to compact bit-packing (e.g., LOUDS bitmaps with \textasciitilde{}1 bit per slot). Affects decode latency and L1/\allowbreak{}L2 utilization. \\
HM4 & Stride Pattern (planned, no code tag) & Addressing pattern during traversal: sequential strides (consecutive records), interleaved (e.g., to separate hot/\allowbreak{}cold data), or cache-aware layouts (e.g., Eytzinger BFS). Optimizes prefetch efficiency. In the code without a carrying layout building block and without a tag: a permutable stride value would break the \texttt{All\allowbreak{}Layouts} list of the 320 catalogue, which is fixed at five layout building blocks. \\
\bottomrule
\end{longtable}
\end{scriptsize}

\subsection*{Dynamic sub-axes}

\begin{scriptsize}
\begin{longtable}{@{}>{\raggedright\arraybackslash}p{1.8cm} >{\raggedright\arraybackslash}p{3.2cm} >{\raggedright\arraybackslash}p{8.5cm}@{}}
\toprule
ID & Parameter & Description \\
\midrule
\endhead%
Cache\allowbreak{}Line\allowbreak{}Size & Cache-line size (line-size dimension of the \texttt{cacheline} sub-axis) & Planner parameter for the target platform cache line: 32 bytes (embedded), 64 bytes (standard x86-64/\allowbreak{}ARM64), 128 bytes (among others some Power ISAs) or 256 bytes (for example s390x). The value space of the \texttt{cacheline} sub-axis comprises 4 line sizes $\times$ 3 alignment values $\times$ 5 prefetch-hint values = 60 configurations; the choice is applied per build as the NTTP \texttt{Cache\allowbreak{}Line\allowbreak{}Config} and affects alignment/\allowbreak{}padding overhead for all aligned layouts. \\
\bottomrule
\end{longtable}
\end{scriptsize}

\subsection*{Building blocks}

\begin{scriptsize}
\begin{longtable}{@{}>{\raggedright\arraybackslash}p{3.0cm} >{\raggedright\arraybackslash}p{7.0cm} >{\raggedright\arraybackslash}p{3.5cm}@{}}
\toprule
Variant & Description & Source \\
\midrule
\endhead%
\texttt{Ao\allowbreak{}SStrict\allowbreak{}Memory\allowbreak{}Layout} & Array-of-Structs, tightly packed without cache-line alignment. Advantage: minimal memory footprint (record\_\allowbreak{}size strides, dense). Disadvantage: false-sharing under concurrent writes, poor cache-hit ratio for narrow field scans. Typical for Wormhole index. & Code/\allowbreak{}Baseline \\
\texttt{Cache\allowbreak{}Line\allowbreak{}Aligned\allowbreak{}Memory\allowbreak{}Layout} & Ao\allowbreak{}S with 64-byte per-record alignment. Avoids false-sharing via padding (stride = round\_\allowbreak{}up(record\_\allowbreak{}size, 64)). Standard for ART/\allowbreak{}HOT/\allowbreak{}Masstree/\allowbreak{}START\@. Build constant: at record\_\allowbreak{}size=48 (measurement build), 16 bytes of padding arise per record; a measurably higher latency than with aos\_\allowbreak{}strict is the expected cache effect (measurement evidence in the corpus: n/a, the complete measurement series is not available). & Code/\allowbreak{}Baseline \\
\texttt{So\allowbreak{}AMemory\allowbreak{}Layout} & Struct-of-Arrays: all fields stored column-wise sequentially. Advantage: contiguous field scans (e.g., keys-only without values) maximize cache-line utilization (16$\times$ fewer cache-lines than Ao\allowbreak{}S for single-field scan). Ideal for SIMD vectorization and OLAP indexes. Benefit described in Abadi et al. SIGMOD 2008. & Abadi et al.\ 2008~\cite{abadi2008columnrow} \\
\texttt{Packed\allowbreak{}Bitmap\allowbreak{}Memory\allowbreak{}Layout} & Bit-packed succinct layout (LOUDS /\allowbreak{} Su\allowbreak{}RF FST). Stores n entries in \textasciitilde{}n*lg(sigma) bits instead of bytes. Extreme compactness (32$\times$ fewer cache-lines than Ao\allowbreak{}S), but higher decode cost per lookup. Foundation: Jacobson LOUDS FOCS 1989 (P09). Modern impl: Su\allowbreak{}RF Apache-2.0. & P09 (Jacobson LOUDS FOCS 1989) \\
\texttt{Ao\allowbreak{}So\allowbreak{}AMemory\allowbreak{}Layout} & Array-of-Structures-of-Arrays: hybrid layout with block width W=8 (AVX2 lanes). Within block: So\allowbreak{}A (fields contiguous). Across blocks: Ao\allowbreak{}S (blocks strided). Unites SIMD vectorization (So\allowbreak{}A within) with spatial locality (Ao\allowbreak{}S across). Standard in HPC/\allowbreak{}SIMD engines. Convention; no single originating paper, but an idiom of the SPMD compiler ispc~\cite{pharr2012ispc}. & Code/\allowbreak{}Baseline (HPC/\allowbreak{}SIMD convention; ispc~\cite{pharr2012ispc}) \\
\bottomrule
\end{longtable}
\end{scriptsize}

\emph{Contributing sources:} P09, Abadi et al.\ 2008, ispc (Pharr/\allowbreak{}Mark 2012)~\cite{jacobson1989louds,abadi2008columnrow,pharr2012ispc}.

\emph{Related literature:} P06, P11, P12, P32~\cite{schmeisser2022b2tree,rao1999css,rao2000csb,schmidt2025hbm}.

\section{T6 Allocator Axis (T6)}

This axis is the \emph{allocation organ} (T6) of the living being \texttt{SearchAlgorithm} (cf.\ Section~\ref{sec:sota-instances}).

This axis provides the memory management strategy for the cache engine. It defines 26 concrete allocator wrappers for various algorithms and techniques (from classical baselines like Buddy and Slab through modern research allocators like mimalloc, jemalloc, snmalloc to specialized variants for NUMA, PIM and formally verified systems as well as a memory-type virtualization following the VAMPIR design, A24). All wrappers satisfy standardized concepts and are orthogonally structurable by 7 compile-time variant families (freelist topology, size-class schema, thread locality, synchronization, allocation policy, reclamation, fragmentation strategy).

\subsection*{Static sub-axes}

\begin{scriptsize}
\begin{longtable}{@{}>{\raggedright\arraybackslash}p{1.8cm} >{\raggedright\arraybackslash}p{3.2cm} >{\raggedright\arraybackslash}p{8.5cm}@{}}
\toprule
ID & Name & Description \\
\midrule
\endhead%
AA1 & Free-List Topology & How are free memory blocks organized? Examples: per-page sharding (mimalloc A04), per-thread caching (snmalloc A07), buddy trees (A19), single-global free-list (dlmalloc A20), spans (scalloc A08), hybrid strategies (hmalloc A15). \\
AA2 & Size-Class Schema & How are block sizes classified? Examples: slab per object size (A02), dlmalloc bins (A20), jemalloc size-classes (A05), scalloc spans (A08), buddy power-of-2 (A19), formally verified bins (Star\allowbreak{}Malloc A13), std::pmr pools (A22). \\
AA3 & Thread Locality & How thread/\allowbreak{}core-local is allocation? Examples: tcmalloc thread-local caches (A06), rpmalloc thread-heaps (A10), snmalloc message-passing (A07), tcmalloc-warehouse hyperscale (A14), Hoard per-heap locks (A01). \\
AA4 & Synchronization & What concurrency mechanism? Examples: lock-free CAS (Michael A03, LRMalloc A11), per-heap locks (Hoard A01), single-threaded (Exgen A18), wait-free (Crystalline A17). \\
AA5 & Allocation Policy & How are allocation requests routed? Examples: NUMA-origin awareness (NUMAlloc A09), cache-set awareness (CAMA A12), NUCA awareness (TCMalloc-Warehouse A14), PIM awareness (PIM-malloc A16), polymorphic memory resource (std::pmr A22), NFP memory-type virtualization (Vampir\allowbreak{}Nfp A24). \\
AA6 & Reclamation & How are deleted blocks reclaimed? Examples: magazine cache (Slab A02), page-heap release (TCMalloc A06), lock-free reclamation (Crystalline A17), deferred-free (mimalloc A04), Vmem magazine extension (A23). \\
AA7 & Fragmentation Strategy & How is fragmentation prevented? Examples: dlmalloc coalescing (A20), object colouring (Slab A02), buddy splitting (A19), page-local sharding (mimalloc A04), span pooling (scalloc A08). \\
\bottomrule
\end{longtable}
\end{scriptsize}

\subsection*{Building blocks}

\begin{scriptsize}
\begin{longtable}{@{}>{\raggedright\arraybackslash}p{4.0cm} >{\raggedright\arraybackslash}p{6.5cm} >{\raggedright\arraybackslash}p{3.0cm}@{}}
\toprule
Variant & Description & Source \\
\midrule
\endhead%
\texttt{Hoard\allowbreak{}Allocator} (A01) & Per-heap free-lists with SMP scalability. Berger et al. (ASPLOS 2000). Synchronization concept via per-heap locks + local heap caches per thread. & P— (OSS, Apache-2.0) \\
\texttt{Slab\allowbreak{}Allocator} (A02) & Object-caching kernel allocator with magazine-cache strategy. Bonwick (USENIX 1994, 2001). Classic textbook paradigm without standalone OSS code. & P— (Baseline/\allowbreak{}Textbook) \\
\texttt{Michael\allowbreak{}Lock\allowbreak{}Free\allowbreak{}Allocator} (A03) & Lock-free CAS allocator with hazard pointers. Maged Michael (PLDI 2004). Re-implementation with IBM patent background. & P— (LGPL Re-Impl) \\
\texttt{Mimalloc\allowbreak{}Allocator} (A04) & Free-list sharding with deferred-free and page-local management. Leijen/\allowbreak{}Zorn/\allowbreak{}de Moura (MSR-TR-2019-18, APLAS 2019). Original code linking (is\_\allowbreak{}original=2/\allowbreak{}2). & P— (OSS, MIT, is\_\allowbreak{}original) \\
\texttt{Jemalloc\allowbreak{}Allocator} (A05) & Arena-based size-classes with per-CPU arenas and tcache layer. Evans (BSDCan 2006). Standard allocator of Free\allowbreak{}BSD and Facebook. Original code linking (is\_\allowbreak{}original=2/\allowbreak{}2). & P— (OSS, BSD-2-Clause, is\_\allowbreak{}original) \\
\texttt{TCMalloc\allowbreak{}Allocator} (A06) & Thread-caching malloc with central page-heap. Ghemawat/\allowbreak{}Menage (Google ca. 2005) + TEMERAIRE (OSDI 2021). Google production standard. & P— (OSS, Apache-2.0) \\
\texttt{Snmalloc\allowbreak{}Allocator} (A07) & Message-passing allocator with asymmetric heap operations. Paul Liétar et al. (ISMM 2019). Original code linking (is\_\allowbreak{}original=2/\allowbreak{}2). & P— (OSS, MIT, is\_\allowbreak{}original) \\
\texttt{Scalloc\allowbreak{}Allocator} (A08) & Span-based free-list with virtual spans for multi-core scalability. Aigner/\allowbreak{}Kirsch/\allowbreak{}Lippautz/\allowbreak{}Sokolova (OOPSLA 2015). Low fragmentation via span pooling. & P— (OSS, BSD-3-Clause) \\
\texttt{NUMAlloc\allowbreak{}Allocator} (A09) & NUMA-origin awareness with faster memory allocation. UTSASRG (ISMM 2023). Optimized for multi-socket architectures. & P— (OSS, licence not determined) \\
\texttt{RPMalloc\allowbreak{}Allocator} (A10) & Per-thread span caching without dependencies on other allocator libraries. Mattias Jansson (2017). Original code linking (is\_\allowbreak{}original=2/\allowbreak{}2). & P— (OSS, Public-Domain/\allowbreak{}MIT, is\_\allowbreak{}original) \\
\texttt{LRMalloc\allowbreak{}Allocator} (A11) & Lock-free allocator with hazard pointers for multi-core scalability. Leite/\allowbreak{}Rocha (VECPAR 2018, LNCS 11333). Original code linking (is\_\allowbreak{}original=2/\allowbreak{}2). & P— (OSS, MIT, is\_\allowbreak{}original) \\
\texttt{CAMAAllocator} (A12) & Cache-Aware Memory Allocator with predictable behaviour. Herter/\allowbreak{}Backes/\allowbreak{}Haupenthal/\allowbreak{}Reineke (ECRTS 2011, Saarland). No standalone OSS implementation. & P— (Baseline/\allowbreak{}Research) \\
\texttt{Star\allowbreak{}Malloc\allowbreak{}Allocator} (A13) & Formally verified hardened memory allocator with F*/\allowbreak{}Steel. Inria-Prosecco: Reitz/\allowbreak{}Fromherz/\allowbreak{}Protzenko et al. (OOPSLA 2024). Mathematically verified allocator. & P— (OSS, Apache-2.0) \\
\texttt{TCMalloc\allowbreak{}Warehouse\allowbreak{}Allocator} (A14) & TCMalloc hyperscale extension with Huge\allowbreak{}Page support for warehouse-scale computing. TEMERAIRE (Hunter OSDI 2021) + warehouse characterization (Zhou ASPLOS 2024). NUCA-aware TCMalloc variant. & P— (OSS, Apache-2.0) \\
\texttt{HMalloc\allowbreak{}Allocator} (A15) & Hybrid lock-free allocator with scalable memory management. Li/\allowbreak{}Yao/\allowbreak{}Tang/\allowbreak{}Lin (IEEE ICPADS 2019). No standalone OSS implementation. & P— (Research Baseline) \\
\texttt{PIMMalloc\allowbreak{}Allocator} (A16) & Processing-in-memory allocator for heterogeneous hardware with fast local memory. VIA-Research (HPCA 2026, ar\allowbreak{}Xiv:2505.13002). Specialized for PIM hardware. & P— (OSS, MIT) \\
\texttt{Crystalline\allowbreak{}Reclamation} (A17) & Wait-free safe-memory-reclamation family (Nikolaev/\allowbreak{}Ravindran DISC 2021 / PLDI 2024). Reclamation policy, not a standalone allocator. & P— (Research) \\
\texttt{Exgen\allowbreak{}Allocator} (A18) & Single-threaded specialized allocator with optimizations for single-threaded workloads. UT Austin (IEEE CAL 2025, ar\allowbreak{}Xiv:2510.10219)\@. \enquote{Old is Gold} principle. & P— (Research, licence not determined) \\
\texttt{Buddy\allowbreak{}Allocator} (A19) & Buddy system with power-of-2 splitting for simple fragmentation control. Knuth (TAo\allowbreak{}CP Vol.1, 1968). Classic textbook paradigm. & P— (Baseline/\allowbreak{}Textbook) \\
\texttt{Dlmalloc\allowbreak{}Allocator} (A20) & Doug Lea malloc with bins-based classification. Lea (1987--2012). Original code linking (is\_\allowbreak{}original=2/\allowbreak{}2). & P— (OSS, Public-Domain/\allowbreak{}CC0, is\_\allowbreak{}original) \\
\texttt{Pt\allowbreak{}Malloc2\allowbreak{}Allocator} (A21) & ptmalloc2 (glibc malloc) with explicit wrapper control. Wolfram Gloger /\allowbreak{} Ulrich Drepper (1990s-2000s). Licence-blocked (LGPL-2.1+). & P— (OSS, LGPL-2.1+, blocks linking) \\
\texttt{Std\allowbreak{}Malloc} (A22a) & Standard libc malloc (proof of concept) — portable fallback allocator. Baseline for all standard library implementations. & P— (Standard-Library) \\
\texttt{Pmr\allowbreak{}Resource\allowbreak{}Allocator} (A22b) & std::pmr::memory\_\allowbreak{}resource (polymorphic memory resource). Halpern (N3916, C++17). Abstract interface without behaviourally distinct default implementation. & P— (Standard-Library, C++17) \\
\texttt{Pool\allowbreak{}Resource\allowbreak{}Allocator} (A22c) & std::pmr::unsynchronized\_\allowbreak{}pool\_\allowbreak{}resource with own size-class pool. Halpern (N3916, C++17). Behaviourally distinct without vendor linking (F15-operative). & P— (Standard-Library, C++17, F15-operative) \\
\texttt{Vmem\allowbreak{}Magazines\allowbreak{}Allocator} (A23) & Vmem + magazines as extension of the slab allocator for multi-core scaling. Bonwick (USENIX ATC 2001). No standalone OSS implementation. & P— (Baseline/\allowbreak{}Textbook) \\
\texttt{Vampir\allowbreak{}Nfp\allowbreak{}Allocator} (A24) & Allocator front following the VAMPIR design (virtualization of non-functional memory properties): allocations pass through an NFP-labelled memory tier (target tier DRAM), the NFP descriptor (\texttt{bandwidth\_\allowbreak{}class}, \texttt{latency\_\allowbreak{}class}, \texttt{capacity}) is exposed. Structural reconstruction of the buildable memory-type virtualization part only: poster and review provide neither code nor a cost formula; energy is deliberately not modelled (no counter, otherwise an empty metric); compensation, transparent OS migration and global pipeline scheduling are out of scope. \texttt{is\_\allowbreak{}original}=false, disabled by default. & P33 (VAMPIR poster SOSP 2023~\cite{vampir2023sosp}, TU Dresden; reconstruction, is\_\allowbreak{}original=false) \\
\bottomrule
\end{longtable}
\end{scriptsize}

\emph{Registry state:} the generated organ registry carries three enabled building blocks for the allocator axis --- \texttt{std\_\allowbreak{}malloc}, \texttt{pmr\_\allowbreak{}resource} and \texttt{pool\_\allowbreak{}resource}; the catalogue A01--A24 is retained as the offering catalogue, the remaining wrappers can be switched on per vendor availability and enable flag.

\emph{Contributing reclamation sources:} P29, P30~\cite{mckenney2001rcu,michael2004hazard}; for the full allocator corpus A01--A24 of the code (the literature corpus A01--A23 of Section~\ref{sec:sota-instances} plus the allocator front A24 reconstructed from the SOTA profile P33 without an allocator profile of its own) see section \enquote{A-Korpus}.

\section{T7 Prefetch Strategy}

This axis is the \emph{prefetch organ} (T7) of the living being \texttt{SearchAlgorithm} (cf.\ Section~\ref{sec:sota-instances}).

The prefetch axis determines the strategy for anticipatory loading of nodes and addresses along the expected search path in a trie data structure. It supports four concrete strategies with different trigger mechanisms and heuristics: no prefetch (measurement baseline), distance-estimated software prefetch based on node density and cache latency (ART standard), explicit hardware prefetch instructions (Wormhole standard), and path-oriented prefetch with bundle granularity (prt-art, this thesis).

\subsection*{Static sub-axes}

\begin{scriptsize}
\begin{longtable}{@{}>{\raggedright\arraybackslash}p{1.8cm} >{\raggedright\arraybackslash}p{3.2cm} >{\raggedright\arraybackslash}p{8.5cm}@{}}
\toprule
ID & Name & Description \\
\midrule
\endhead%
PF1 & Trigger Mechanism & Type of control over the prefetch timing: whether no mechanism (baseline), through compiler hints via distance estimator, explicit hardware prefetch instructions, or path-oriented prediction \\
PF2 & Distance Heuristic & Calculation method for prefetch distance (how far ahead to load): linear, adaptive based on node density and tier latency, or no heuristic \\
PF3 & Granularity & Prefetch size per trigger: cache-line (64 bytes, individual), page (4\allowbreak{}K, OS window), or bundle (multiple cache-lines, trie-path-oriented) \\
\bottomrule
\end{longtable}
\end{scriptsize}

\subsection*{Building blocks}

\begin{scriptsize}
\begin{longtable}{@{}>{\raggedright\arraybackslash}p{3.0cm} >{\raggedright\arraybackslash}p{7.0cm} >{\raggedright\arraybackslash}p{3.5cm}@{}}
\toprule
Variant & Description & Source \\
\midrule
\endhead%
\texttt{None\allowbreak{}Prefetch} & No prefetch strategy (baseline). Serves as comparison reference for measurement series. Trigger: none; heuristic: none; granularity: N/\allowbreak{}A. \texttt{is\_\allowbreak{}active()} returns false. & Code/\allowbreak{}Baseline \\
\texttt{Distance\allowbreak{}Estimator\allowbreak{}Prefetch} & Distance-estimated software prefetch (ART standard per Leis ICDE 2013). Estimates cache distance to next node based on node density [\%] and tier latency [cycles] with heuristic: estimate = clamp(1 + sparseness*8 + latency\_\allowbreak{}factor, [1,16]). Trigger: compiler hints via \texttt{\_\allowbreak{}\_\allowbreak{}builtin\_\allowbreak{}prefetch}; heuristic: adaptive (density + latency); granularity: cache-line. Applies ART pattern (unodb/\allowbreak{}libart). & P01 (ART/\allowbreak{}unodb, Leis ICDE 2013) \\
\texttt{Hardware\allowbreak{}Prefetch} & Explicit hardware prefetch instructions (Wormhole standard per Wu Euro\allowbreak{}Sys 2019). Uses PREFETCHT0/\allowbreak{}T1/\allowbreak{}T2/\allowbreak{}NTA on hash-anchor-chain without software heuristic; CPU estimates distance autonomously. Trigger: explicit hardware instruction; heuristic: hardware-managed; granularity: hardware-determined. License GPL-3.0 blocks \texttt{is\_\allowbreak{}original} linking. & P07 (Wormhole, Wu Euro\allowbreak{}Sys 2019) \\
\texttt{Path\allowbreak{}Oriented\allowbreak{}Prefetch} & Path-oriented prefetch heuristic (prt-art, this thesis). Pre-loads all nodes along the predicted trie path before the search starts. Tracker follows the search path and recommends the next address via linear step extrapolation. Hot-path hints: interprets the first 8 bytes of a key as uint64 address. Trigger: path-oriented prediction; heuristic: linear (step extrapolation); granularity: bundle (multiple cache-lines). Related literature: Chen/\allowbreak{}Gibbons/\allowbreak{}Mowry SIGMOD 2001 (p\allowbreak{}B+-Trees Prefetching) + Mowry 1994. & P21 (related: Chen/\allowbreak{}Gibbons/\allowbreak{}Mowry SIGMOD 2001); this thesis \\
\bottomrule
\end{longtable}
\end{scriptsize}

\emph{Contributing sources:} P01, P07, P21~\cite{leis2013art,wu2019wormhole,chen2001prefetch}.

\emph{Related literature:} P23, P25, P26, P27~\cite{khan2010adaptive,mahling2025hotpath,zhang2024pathprefix,zhang2025hierarchical}.

\section[head={T8 Concurrency Synchronization Axis (Concurrency Axis)}]{T8 Concurrency Synchronization Axis (Concurrency Axis T8)}

This axis is the \emph{concurrency organ} (T8) of the living being \texttt{SearchAlgorithm} (cf.\ Section~\ref{sec:sota-instances}).

This axis controls synchronization patterns and safe memory reclamation for concurrent access to tree/\allowbreak{}index structures. The axis provides compile-time selectable synchronization patterns (None/\allowbreak{}Blocking/\allowbreak{}Reader-Writer/\allowbreak{}Optimistic/\allowbreak{}Lock-Free/\allowbreak{}Wait-Free) as well as orthogonal reclamation schemes (RCU/\allowbreak{}Hazard-Pointer) for safe memory deallocation in lock-free structures.

\subsection*{Static sub-axes}

\begin{scriptsize}
\begin{longtable}{@{}>{\raggedright\arraybackslash}p{1.8cm} >{\raggedright\arraybackslash}p{3.2cm} >{\raggedright\arraybackslash}p{8.5cm}@{}}
\toprule
ID & Name & Description \\
\midrule
\endhead%
CC1 & Synchronization Pattern (CC1) & Compile-time variant family of synchronization patterns: None (single-threaded baseline), Blocking (pessimistic mutex), Reader\allowbreak{}Writer (shared\_\allowbreak{}mutex), Optimistic (lock-free version validation), Lock\allowbreak{}Free (CAS-loops), Wait\allowbreak{}Free (bounded steps). Each family has distinct runtime characteristics (0-overhead to mutex-overhead). \\
CC2 & Reclamation Scheme (CC2) & Orthogonal compile-time variant family for safe memory deallocation in lock-free structures: RCU (deferred grace-period via userspace-rcu), Hazard-Pointer (per-thread in-use marks). Composable with lock-free patterns, optional for Blocking/\allowbreak{}Optimistic. \\
\bottomrule
\end{longtable}
\end{scriptsize}

\subsection*{Building blocks}

\begin{scriptsize}
\begin{longtable}{@{}>{\raggedright\arraybackslash}p{3.0cm} >{\raggedright\arraybackslash}p{8.0cm} >{\raggedright\arraybackslash}p{2.5cm}@{}}
\toprule
Variant & Description & Source \\
\midrule
\endhead%
\texttt{None\allowbreak{}Concurrency} & No synchronization, single-threaded baseline. Both acquire() and release() are no-ops (volatile dummy). Comparison zero point for concurrency-overhead measurement (F15). & Code/\allowbreak{}Baseline \\
\texttt{Blocking\allowbreak{}Concurrency} & Pessimistic mutex locking with std::mutex. Global coarse-grained lock. Serializes all accesses. Upper bound of the locking overhead. Publication: Dijkstra CACM 1965. & Code/\allowbreak{}Baseline \\
\texttt{Reader\allowbreak{}Writer\allowbreak{}Concurrency} & Single-writer/\allowbreak{}multi-reader via std::shared\_\allowbreak{}mutex. Lock-sharing for read-heavy workloads. Middleweight between pessimistic and optimistic patterns. Publication: Courtois et al. CACM 1971. & Code/\allowbreak{}Baseline \\
\texttt{Olc\allowbreak{}Optimistic\allowbreak{}Concurrency} & Optimistic Lock-Coupling (ART-Sync): version counter per node, lock-free readers with retry-on-conflict. Cheap under low contention. Publication: Leis Da\allowbreak{}Mo\allowbreak{}N 2016 /\allowbreak{} DEBULL 2019 (P08, is\_\allowbreak{}original=true, Apache-2.0, unodb). & P08 \\
\texttt{Lock\allowbreak{}Free\allowbreak{}Concurrency} & Lock-free via Compare-And-Swap (std::atomic CAS-loops). System-wide progress. Generic CAS pattern without blocking. Publication: Michael/\allowbreak{}Scott PODC 1996 (reference paper). & Code/\allowbreak{}Baseline \\
\texttt{Wait\allowbreak{}Free\allowbreak{}Concurrency} & Wait-free: each thread completes in a bounded number of steps (strongest guarantee). Without retry loop like Lock\allowbreak{}Free. Realized via sequence lock /\allowbreak{} fetch-add. Publication: Herlihy TOPLAS 1991 (theory reference). & Code/\allowbreak{}Baseline \\
\texttt{Rcu\allowbreak{}Concurrency} & Read-Copy-Update (Mc\allowbreak{}Kenney OLS 2001): lock-free readers via per-thread nesting counter (relaxed order, \emph{no} memory barrier on the read side). Reclamation scheme orthogonal to pattern. Local re-impl (P29 userspace-rcu LGPL-2.1, is\_\allowbreak{}original=false). & P29 \\
\texttt{Hazard\allowbreak{}Pointer\allowbreak{}Concurrency} & Hazard Pointers (Michael TPDS 2004, P30): per-thread marked in-use pointers prevent premature deallocation (ABA-safe). Reclamation scheme. Local re-impl (P30 NO LICENSE, is\_\allowbreak{}original=false). SEQ\_\allowbreak{}CST store on acquire(). & P30 \\
\texttt{Olc\allowbreak{}Reserved\allowbreak{}Blocks\allowbreak{}Concurrency} & OLC + cache-line-reserved value blocks (multi-writer cache-coherence-storm avoidance). Specialization of the Optimistic pattern with additional atomic block-ID reservation (fetch\_\allowbreak{}add)\@. prt-art option (\texttt{is\_\allowbreak{}original}=false). & Code/\allowbreak{}Baseline \\
\bottomrule
\end{longtable}
\end{scriptsize}

\emph{Contributing sources:} Dijkstra 1965, Courtois et al.\ 1971, P08, Michael/\allowbreak{}Scott 1996, Herlihy 1991, P29, P30~\cite{dijkstra1965concurrent,courtois1971readers,leis2016olc,michael1996queue,herlihy1991waitfree,mckenney2001rcu,michael2004hazard}.

\section{T9 Serialization (Encoding and Compression)}

This axis is the \emph{serialization organ} (T9) of the living being \texttt{SearchAlgorithm} (cf.\ Section~\ref{sec:sota-instances}).

Axis \texttt{axis\_\allowbreak{}10} (organ slot T9) determines how data is encoded and compressed within the cache-engine system: from raw-binary storage layouts via variable-length encoding to succinct bit-packing. The choice of serialization strategy directly impacts CPU cost (delta encoding, bit manipulation, LEB128 varint), memory overhead and IO-bandwidth utilization --- critical for LSM-trees and read-heavy workloads with approximate filters (Su\allowbreak{}RF/\allowbreak{}LOUDS).

\subsection*{Static sub-axes}
\begin{scriptsize}
\begin{longtable}{@{}>{\raggedright\arraybackslash}p{1.8cm} >{\raggedright\arraybackslash}p{3.2cm} >{\raggedright\arraybackslash}p{8.5cm}@{}}
\toprule
ID & Name & Description \\
\midrule
\endhead%
SR1 & Byte-Order and Word Format & How are bytes arranged within a record? Raw (memcpy raw-layout), Varlen (LEB128-varint with signalling bits), or Packed (bit-by-bit packing)? \\
SR2 & Density (Memory Utilization) & Memory overhead per record: Sparse (complete records), Dense (with padding), or Succinct (information-theoretically dense, n*H + o(n) bits)? \\
SR3 & Compression and Preprocessing & Block compression (LZ4/\allowbreak{}Snappy, LZ77-based) or delta-encoding (zigzag-transform for negative deltas)? Or no compression? \\
\bottomrule
\end{longtable}
\end{scriptsize}

\subsection*{Building blocks}
\begin{scriptsize}
\begin{longtable}{@{}>{\raggedright\arraybackslash}p{3.0cm} >{\raggedright\arraybackslash}p{8.5cm} >{\raggedright\arraybackslash}p{2.0cm}@{}}
\toprule
Variant & Description & Source \\
\midrule
\endhead%
\texttt{Raw\allowbreak{}Binary\allowbreak{}Serialization} & Raw byte storage layout (memcpy). Baseline without transformation. Minimal CPU cost. Typical for in-memory workloads without compression. Variant: \texttt{serialize\_\allowbreak{}scan()} reads 32-bit fields + order-sensitive checksum. & Code/\allowbreak{}Baseline \\
\texttt{Var\allowbreak{}Len\allowbreak{}Serialization} & Variable-length encoding with signalling bits (LEB128-Var\allowbreak{}Int). Standard for ART (Leis ICDE 2013 P01). Small values = few bytes. Signalling bits mark byte-length. Data-dependent loop (1-5 bytes). FNV-1a-mix of varint bytes. & P01 \\
\texttt{Succinct\allowbreak{}Serialization} & Bit-by-bit packing with minimal bit-width (LOUDS/\allowbreak{}Su\allowbreak{}RF). Achieves information-theoretic minimum n*H + o(n) bits. Read-only after build. Optimal for approx-filters + read-heavy. Highest per-element CPU cost (bit manipulation). Variant: \texttt{serialize\_\allowbreak{}scan()} with 64-bit accumulator. & P10,P09 \\
\texttt{Compressed\allowbreak{}Serialization} & Block compression (LZ4/\allowbreak{}Snappy) layered on raw-binary. Trade-off: CPU cost vs IO/\allowbreak{}memory bandwidth. Variant: delta-encoding against predecessor + zigzag-transform (negative delta $\rightarrow$ small unsigned). Typical for LSM-trees. Medium CPU cost. & Code/\allowbreak{}LZ77 (Ziv/\allowbreak{}Lempel 1977)~\cite{ziv1977lz77} \\
\bottomrule
\end{longtable}
\end{scriptsize}

\emph{Contributing sources:} P01, P09, P10, Ziv/\allowbreak{}Lempel 1977~\cite{leis2013art,jacobson1989louds,zhang2018surf,ziv1977lz77}.

\section{T10 Value Handle (Value-Handle)}

This axis is the \emph{value-handle organ} (T10) of the living being \texttt{SearchAlgorithm} (cf.\ Section~\ref{sec:sota-instances}).

The value-handle axis defines how data values (Values) are stored and accessed in the cache index: inline within node slots, externally in a pool, via immutable-shared references (RCU-style), versioned pointers (MVCC), or chained external references (multi-value). This axis characterizes the storage-layout strategy for values and thus determines access latency (direct vs.\ indirect) and concurrency semantics (snapshot vs.\ versioning).

\subsection*{Static sub-axes}
\begin{scriptsize}
\begin{longtable}{@{}>{\raggedright\arraybackslash}p{1.8cm} >{\raggedright\arraybackslash}p{3.2cm} >{\raggedright\arraybackslash}p{8.5cm}@{}}
\toprule
ID & Name & Description \\
\midrule
\endhead%
VH1 & Storage Location & Compile-time variants for the physical placement of the value: inline in the node slot (combines key/\allowbreak{}value in one slot without indirection) or external in a separate pool (the node holds offset/\allowbreak{}pointer only). \\
VH2 & Ownership Semantics & Compile-time variants for ownership and sharing strategy: owned (exclusive per key), shared (immutable reference-counted via RCU-style), weak (borrowing without increment). \\
VH3 & Versioning & Compile-time variants for concurrency control: no versioning (simple ownership), version tags in pointer (MVCC for reader snapshot isolation). \\
\bottomrule
\end{longtable}
\end{scriptsize}

\subsection*{Building blocks}
\begin{scriptsize}
\begin{longtable}{@{}>{\raggedright\arraybackslash}p{3.2cm} >{\raggedright\arraybackslash}p{8.0cm} >{\raggedright\arraybackslash}p{2.3cm}@{}}
\toprule
Variant & Description & Source \\
\midrule
\endhead%
\texttt{Inline\allowbreak{}Value\allowbreak{}Handle} & Value embedded directly in the node slot (no indirection). Standard for small values like uint64\_\allowbreak{}t. Minimal access latency (1 direct read), but node memory is consumed by larger values. Source: ART (Leis ICDE 2013, P01). & P01 \\
\texttt{External\allowbreak{}Pool\allowbreak{}Value\allowbreak{}Handle} & Value stored externally in a pool; the node holds the pool offset only. Enables variable value sizes and more compact nodes, but costs 1 additional cache miss per lookup (pointer chasing). Source: Wormhole (Wu Euro\allowbreak{}Sys 2019, P07). & P07 \\
\texttt{Immutable\allowbreak{}Shared\allowbreak{}Ref\allowbreak{}Value\allowbreak{}Handle} & Value via reference-counted pointer (RCU-style snapshot sharing). Optimal for reader concurrency: readers see stable snapshots, writers create new immutable snapshots. Requires Ref\allowbreak{}Count acquire/\allowbreak{}release per access. Source: RCU (Mc\allowbreak{}Kenney OLS 2001, P29). & P29 \\
\texttt{Versioned\allowbreak{}Pointer\allowbreak{}Value\allowbreak{}Handle} & Pointer with version tag in the upper byte (MVCC). Enables multi-version concurrency control: readers get a snapshot version for visibility control, writers create a new version. Tombstone detection via version bit. Source: Masstree (Mao Euro\allowbreak{}Sys 2012, P03) and SMART (OSDI 2023). & P03 \\
\texttt{Chain\allowbreak{}Ref\allowbreak{}Value\allowbreak{}Handle} & Multi-value chained external reference (linked list). The node holds the offset to the chain head in the pool; each chain node holds (value\_\allowbreak{}offset, next\_\allowbreak{}offset). Most expensive variant with 2 dependent dereferences per value access. Source: prt-art (local re-implementation). & Code/\allowbreak{}Baseline \\
\bottomrule
\end{longtable}
\end{scriptsize}

\emph{Contributing sources:} P01, P03, P07, P29, SMART (OSDI 2023)~\cite{leis2013art,mao2012masstree,wu2019wormhole,mckenney2001rcu,luo2023smart}.

\section{T11 Index Organization}

This axis is the \emph{index-organization organ} (T11) of the living being \texttt{SearchAlgorithm} (cf.\ Section~\ref{sec:sota-instances}).

Classification of data storage and index structuring across three orthogonal dimensions: storage order (heap/\allowbreak{}clustered/\allowbreak{}unordered), index count (none/\allowbreak{}single/\allowbreak{}multiple) and data embedding (separate storage vs.\ embedded in leaf pages). This axis models fundamental DBMS design decisions that decisively influence cache behaviour, access patterns and I/\allowbreak{}O characteristics based on the alignment of index order with storage order. The three sub-axes carry the identifiers IO1--IO3 (index organization); the I/\allowbreak{}O dispatch axis (T12) uses the same identifiers for different subjects.

\subsection*{Static sub-axes}

\begin{scriptsize}
\begin{longtable}{@{}>{\raggedright\arraybackslash}p{1.8cm} >{\raggedright\arraybackslash}p{3.2cm} >{\raggedright\arraybackslash}p{8.5cm}@{}}
\toprule
ID & Name & Description \\
\midrule
\endhead%
IO1 & Storage Order & Arrangement of records in storage: unordered (insertion order in heap), sorted by primary key (clustered), or with pointer indirection (Non\allowbreak{}Clustered). Defines the sequential or random access patterns during range scans. Code token: \texttt{storage\_\allowbreak{}order\_\allowbreak{}tag}. \\
IO2 & Index Multiplicity & Number of supported indexes per table: none (heap baseline), single index (clustered, My\allowbreak{}SQL Inno\allowbreak{}DB primary key), or multiple secondary indexes (Non\allowbreak{}Clustered, SQL Server NONCLUSTERED INDEX). Code token: \texttt{index\_\allowbreak{}count\_\allowbreak{}tag}. \\
IO3 & Data Embedding & Storage placement of the data relative to the index: separate in the row store (Clustered/\allowbreak{}Non\allowbreak{}Clustered with pointer hop), or directly embedded in index leaf pages (IOT/\allowbreak{}trie-style like ART/\allowbreak{}HOT). Code token: \texttt{data\_\allowbreak{}embedding\_\allowbreak{}tag}. \\
\bottomrule
\end{longtable}
\end{scriptsize}

\subsection*{Building blocks}

\begin{scriptsize}
\begin{longtable}{@{}>{\raggedright\arraybackslash}p{3.5cm} >{\raggedright\arraybackslash}p{6.0cm} >{\raggedright\arraybackslash}p{4.0cm}@{}}
\toprule
Variant & Description & Source \\
\midrule
\endhead%
\texttt{Clustered\allowbreak{}Index\allowbreak{}Organization} & Storage order corresponds to index order (1 primary key). Optimized for range scans along the primary key with sequential disk reads (cache-friendly). Standard: My\allowbreak{}SQL Inno\allowbreak{}DB, Postgre\allowbreak{}SQL CLUSTER, SQL Server CLUSTERED INDEX\@. Access pattern: sequential forward iteration without predicate evaluation, O(n) aggregation scan. & Bayer \& Mc\allowbreak{}Creight, Acta Informatica 1972 (10.1007/\allowbreak{}BF00288683) \\
\texttt{Heap\allowbreak{}Index\allowbreak{}Organization} & Storage order = insertion order, \emph{no} index (baseline). Lookup via O(n) full scan with predicate evaluation per record (data-dependent branch, no early exit). Standard: Postgre\allowbreak{}SQL heap tables, HBase. Access pattern: unordered sequential scan with comparison load and branch misprediction. & Garcia-Molina, Ullman \& Widom, Database Systems (Textbook Baseline) \\
\texttt{Iot\allowbreak{}Index\allowbreak{}Organization} & Data directly embedded in index leaf pages (Oracle IOT, SQL Server CLUSTERED INDEX style). Saves the pointer indirection vs.\ Non\allowbreak{}Clustered. Equivalent to B+tree leaves with embedded data or trie nodes (ART/\allowbreak{}HOT/\allowbreak{}Wormhole). Access pattern: sequential like Clustered, but with additional load from the data in the leaf (two fields per record). & Srinivasan et al., VLDB 2000 (Oracle8i IOT)~\cite{srinivasan2000oracle8iiot} \\
\texttt{Non\allowbreak{}Clustered\allowbreak{}Index\allowbreak{}Organization} & Index order $\neq$ storage order, multiple secondary indexes per table. Lookup via index + pointer hop to the data row (random stride, cache misses). Standard: SQL Server NONCLUSTERED INDEX, Postgre\allowbreak{}SQL CREATE INDEX\@. Access pattern: index navigation followed by random storage hops (LCG-determined, actual cache-miss latency from TLB/\allowbreak{}L2 invalidation), then a residual-predicate comparison per fetched record (data-dependent branch on the fetched data row after the pointer hop). & Comer, ACM Computing Surveys 1979 (The Ubiquitous B-Tree, 10.1145/\allowbreak{}356770.356776) \\
\bottomrule
\end{longtable}
\end{scriptsize}

\emph{Contributing sources:} Bayer \& Mc\allowbreak{}Creight 1972, Comer 1979, Srinivasan et al.\ 2000, Garcia-Molina, Ullman \& Widom (Textbook)~\cite{bayer1972btree,comer1979btree,srinivasan2000oracle8iiot,garciamolina2009dbsystems}.


\section{T12 I/\allowbreak{}O Dispatch Axis}

This axis is the \emph{I/O-dispatch organ} (T12) of the living being \texttt{SearchAlgorithm} (cf.\ Section~\ref{sec:sota-instances}).

Modelling and selection of I/\allowbreak{}O strategies for persistence and memory access in the index: RAM-only storage (baseline), buffered via the OS page cache (standard), direct I/\allowbreak{}O with the \texttt{O\_\allowbreak{}DIRECT} flag (NVMe optimization), or memory-mapped file I/\allowbreak{}O (persistent memory). This axis decouples the search and tree structures from the persistence implementation and permits A/\allowbreak{}B measurements of the dispatch overhead profiles under identical workloads. The three sub-axes carry the identifiers IO1--IO3 (I/\allowbreak{}O); the index-organization axis (T11) uses the same identifiers for different subjects.

\subsection*{Static sub-axes}

\begin{scriptsize}
\begin{longtable}{@{}>{\raggedright\arraybackslash}p{1.8cm} >{\raggedright\arraybackslash}p{3.2cm} >{\raggedright\arraybackslash}p{8.5cm}@{}}
\toprule
ID & Name & Description \\
\midrule
\endhead%
IO1 & Persistence Model & Classifies the storage type: RAM-only (no persistence, pure measurement baseline), or disk/\allowbreak{}persistent memory (file-backed indexes with various caching strategies). \\
IO2 & Caching Strategy & Characterizes OS buffer usage: Buffered (OS page cache, read-ahead + write-back, standard), Direct I/\allowbreak{}O (\texttt{O\_\allowbreak{}DIRECT}, page-cache bypass, 512\allowbreak{}B/\allowbreak{}4\allowbreak{}KB alignment), or Memory-Mapped (file-backed \texttt{mmap} with automatic OS paging and page-fault-driven access). \\
IO3 & Write Durability & Defines the consistency guarantees: no \texttt{fsync} (highest performance, no safeguard), \texttt{fsync} after batch (compromise, batch commits), or atomic writes (highest safety, lowest performance). In this implementation primarily present as a conceptual sub-axis; operational realization occurs through the caching strategy. \\
\bottomrule
\end{longtable}
\end{scriptsize}

\subsection*{Building blocks}

\begin{scriptsize}
\begin{longtable}{@{}>{\raggedright\arraybackslash}p{2.5cm} >{\raggedright\arraybackslash}p{8.5cm} >{\raggedright\arraybackslash}p{2.5cm}@{}}
\toprule
Variant & Description & Source \\
\midrule
\endhead%
\texttt{In\allowbreak{}Memory\allowbreak{}Only} & Purely RAM-based index without persistence. Direct, sequential record access without dispatch overhead. Fastest path and comparison zero point of the axis. Used for the measurement baseline and performance referencing under conditions without I/\allowbreak{}O latency. & Code/\allowbreak{}Baseline \\
\texttt{Buffered\allowbreak{}Io} & Standard \texttt{read()}/\allowbreak{}\texttt{write()} via the OS page cache with automatic read-ahead and write-back. Default for general-purpose DBMS\@. Trade-off: convenient, but double buffering (app cache + OS cache leads to higher memory usage). Reference publication: Stonebraker CACM 1981 (concept/\allowbreak{}position paper Operating System Support for Database Management, DOI 10.1145/\allowbreak{}358699.358703). Classification: engineering pattern, no algorithm code of the reference. & P-Stonebraker-1981 \\
\texttt{Direct\allowbreak{}Io} & Bypass of the OS page cache via the \texttt{O\_\allowbreak{}DIRECT} flag (\texttt{open(2)}), NVMe-SSD-optimal. 512\allowbreak{}B/\allowbreak{}4\allowbreak{}KB alignment required. Enables DBMS-owned cache strategies (Rocks\allowbreak{}DB, My\allowbreak{}SQL). Higher CPU cost, but no cache pollution from background activity. No dedicated publication as source (Linux API since 2.4.10); treated as a deliberate engineering baseline. & Code/\allowbreak{}Baseline \\
\texttt{Mmap\allowbreak{}Io} & Memory-mapped file I/\allowbreak{}O (\texttt{mmap}) with persistent memory or read-heavy workloads. The OS pages automatically, no copy between user and kernel space. Trade-off: page-fault latency spikes vs.\ reuse-friendliness. Reference publication: Crotty et al. CIDR 2022 (anti-pattern baseline Are You Sure You Want to Use MMAP in Your DBMS?, db.cs.cmu.edu/\allowbreak{}\ldots/\allowbreak{}cidr2022-p13-crotty.pdf). Deliberately cited as anti-pattern reference to document limits and overhead. OSS reference MIT-licensed. & P-Crotty-2022 \\
\bottomrule
\end{longtable}
\end{scriptsize}

\emph{Contributing sources:} Stonebraker-1981, Crotty-2022~\cite{stonebraker1981os,crotty2022mmap}.

\section{T13 Migration Strategy Axis}

This axis is the \emph{migration organ} (T13) of the living being \texttt{SearchAlgorithm} (cf.\ Section~\ref{sec:sota-instances}).

The migration strategy axis (T13, \texttt{axis\_\allowbreak{}migration}) encapsulates decision logic for data placement across multi-tier storage hierarchies (RAM/\allowbreak{}SSD/\allowbreak{}HDD, hot/\allowbreak{}cold tiers). It models when and how frequently accessed data is promoted to faster storage or rarely accessed data is demoted, without performing an actual tier migration (decision-logic measurement only). Supports four concrete strategies: static (baseline), heat/\allowbreak{}cold separation, multi-tier hierarchy, ML-driven adaptive.

\subsection*{Static sub-axes}

\begin{scriptsize}
\begin{longtable}{@{}>{\raggedright\arraybackslash}p{1.8cm} >{\raggedright\arraybackslash}p{3.2cm} >{\raggedright\arraybackslash}p{8.5cm}@{}}
\toprule
ID & Name & Description \\
\midrule
\endhead%
MG1 & Trigger Mechanism (MG1) & Determines when a migration decision is triggered: none (static), threshold-based (Threshold), observation-driven (Observation). Compile-time variant per strategy. \\
MG2 & Migration Direction (MG2) & Determines the direction of data movement: none (static), upward (RAM→SSD), downward (SSD→RAM), bidirectional. Compile-time variant. \\
MG3 & Migration Granularity (MG3) & Specifies the scaling unit for migration decisions: none (static), Page, Block, Object (single item); the size of the page and block unit is not fixed by the code (pure tag types \texttt{granularity\_\allowbreak{}tag}; the migration building blocks measure the decision logic without a real block transfer). Compile-time category. \\
\bottomrule
\end{longtable}
\end{scriptsize}

\subsection*{Dynamic sub-axes}

The three quantities are design parameters of the axis; the code carries them as compile-time constants inside the decision scans (\texttt{w\_\allowbreak{}freq} = 3 and \texttt{w\_\allowbreak{}rec} = 1 in \texttt{axis\_\allowbreak{}migration\_\allowbreak{}adaptive.hpp}, \texttt{k\allowbreak{}Tiers} = 3 in \texttt{axis\_\allowbreak{}migration\_\allowbreak{}tier\_\allowbreak{}based.hpp}; the hot/\allowbreak{}cold scan decides by recency comparison without a threshold), a run-time or user parameter is implemented for none of the three, and the generated registry carries no run-time sub-axis for T13.

\begin{scriptsize}
\begin{longtable}{@{}>{\raggedright\arraybackslash}p{1.8cm} >{\raggedright\arraybackslash}p{3.2cm} >{\raggedright\arraybackslash}p{8.5cm}@{}}
\toprule
ID & Parameter & Description \\
\midrule
\endhead%
D1 & Decision Threshold & Run-time parameter (plan; no code parameter): hot/\allowbreak{}cold classification boundary (e.g., recency counter limit for LRU-based promotions/\allowbreak{}demotions), user-configurable; the hot/\allowbreak{}cold scan in the code decides per record by recency comparison with the last seen field value, without a boundary. \\
D2 & Feature Scoring Weights & Run-time parameter (plan; a constant in the code) for Adaptive\allowbreak{}Migration: weights (w\_\allowbreak{}freq, w\_\allowbreak{}rec) for frequency and recency features in the ML linear combination, fixed in the code as \texttt{w\_\allowbreak{}freq} = 3 and \texttt{w\_\allowbreak{}rec} = 1. Self-adaptation of the weights foreseen but not implemented; implemented is a running score accumulator with fixed weights. \\
D3 & Tier Configuration & Run-time (plan; a constant in the code): number of storage tiers (e.g., k\allowbreak{}Tiers=3 for RAM/\allowbreak{}SSD/\allowbreak{}HDD in Tier\allowbreak{}Based\allowbreak{}Migration; in the code the compile-time constant \texttt{k\allowbreak{}Tiers} = 3), capacities per tier and promotion/\allowbreak{}demotion thresholds (plan; without a code object). \\
\bottomrule
\end{longtable}
\end{scriptsize}

\subsection*{Building blocks}

\begin{scriptsize}
\begin{longtable}{@{}>{\raggedright\arraybackslash}p{3.0cm} >{\raggedright\arraybackslash}p{7.0cm} >{\raggedright\arraybackslash}p{3.5cm}@{}}
\toprule
Variant & Description & Source \\
\midrule
\endhead%
\texttt{No\allowbreak{}Migration} & Static placement without migration (baseline). Data is positioned once, no movement between tiers. Benchmark zero point for comparison with active strategies. & Code/\allowbreak{}Baseline \\
\texttt{Hot\allowbreak{}Cold\allowbreak{}Migration} & Heat/\allowbreak{}cold separation via LRU + probabilistic recency counting. Frequently accessed data (hot) is promoted to faster cache, rarely accessed data (cold) is demoted to slow storage. Standard for LRU-K and LFU cache systems. & Anti-Caching (PVLDB 2013)~\cite{debrabant2013anticaching} \\
\texttt{Tier\allowbreak{}Based\allowbreak{}Migration} & Multi-level migration across the tiers RAM→SSD→HDD (\texttt{k\allowbreak{}Tiers} = 3 in the code; the paper-code map of the \texttt{ce} repository names the tier sequence RAM→NVM→SSD→HDD, an NVM tier is not implemented in the code) with modulo-based tier mapping. Each data item is mapped modulo-wise to its target tier by a field value. Standard for Rocks\allowbreak{}DB, LSM-trees and tiered cache systems with configurable thresholds; the building block itself has no threshold parameter (table of the dynamic sub-axes). & Managing NVM in DBMSs (SIGMOD 2018)~\cite{vanrenen2018nvm} \\
\texttt{Adaptive\allowbreak{}Migration} & ML-driven adaptive migration with online learning (Cachelib/\allowbreak{}Le\allowbreak{}Ca\allowbreak{}R-style). Weighted linear combination of frequency and recency features computes an adaptive score; a running learning accumulator models online learning. Higher overhead vs.\ optimal hit rate. & ML-based Le\allowbreak{}Ca\allowbreak{}R (USENIX Hot\allowbreak{}Storage 2018)~\cite{vietri2018lecar} \\
\bottomrule
\end{longtable}
\end{scriptsize}

\emph{Contributing sources:} Anti-Caching (PVLDB 2013), LeCaR (Hot\allowbreak{}Storage 2018), Managing NVM (SIGMOD 2018)~\cite{debrabant2013anticaching,vietri2018lecar,vanrenen2018nvm}.

\section{T14 Filter Axis}

This axis is the \emph{filter organ} (T14) of the living being \texttt{SearchAlgorithm} (cf.\ Section~\ref{sec:sota-instances}).

The filter axis T14 offers selectable probabilistic membership-query filters (point and range queries; all four building blocks answer approximately, their construction is deterministic). It covers classic Bloom filters with k hash functions, the deletable cuckoo filter with bucket hashing and fingerprints, the succinct range filter Su\allowbreak{}RF with LOUDS-trie navigation, and the space-saving xor filter with constant 3-way-XOR latency. Purpose: make the trade-offs between space, query latency, mutability and false-positive rate directly measurable in the cache-engine experiment. As an exact comparison zero point the pass-through variant \texttt{None\allowbreak{}Filter} is additionally registered (E14), disabled by default in the build configuration.

\subsection*{Static sub-axes}

\begin{scriptsize}
\begin{longtable}{@{}>{\raggedright\arraybackslash}p{1.8cm} >{\raggedright\arraybackslash}p{3.2cm} >{\raggedright\arraybackslash}p{8.5cm}@{}}
\toprule
ID & Name & Description \\
\midrule
\endhead%
FT1 & Query Type & Distinction between point queries (membership: key present?) and range queries (range: keys in [a,b]?). Point filters are Bloom/\allowbreak{}Cuckoo/\allowbreak{}Xor; only Range\allowbreak{}Surf\allowbreak{}Filter supports range semantics. \\
FT2 & Mutability & Division into immutable filters (Bloom, Range\allowbreak{}Surf, Xor immutable after build) vs.\ mutable (Cuckoo with delete operation). Determines wrapper design and caching possibilities during runtime. \\
FT3 & Error Profile & Characterization of the error-tolerant semantics: false-positive-only (Bloom, Cuckoo, Xor --- filter answers \enquote{yes} incorrectly) vs.\ lossless (theoretically deterministic answer). Bloom: approx.\ k*n hash tests; Xor: exactly 3 XOR operations (branch-free). \\
\bottomrule
\end{longtable}
\end{scriptsize}

\subsection*{Building blocks}

\begin{scriptsize}
\begin{longtable}{@{}>{\raggedright\arraybackslash}p{2.6cm} >{\raggedright\arraybackslash}p{7.4cm} >{\raggedright\arraybackslash}p{3.5cm}@{}}
\toprule
Variant & Description & Source \\
\midrule
\endhead%
\texttt{Bloom\allowbreak{}Filter} & Classical Bloom filter per Bloom 1970: m-bit bitmap, k independent hash functions (double hashing per Kirsch/\allowbreak{}Mitzenmacher 2006). Point queries only; false positives possible, no deletes. Implementation with an actual persistent m-bit bitmap (8~KiB, m=65536, k=4 double hashing), which is probed in the out-of-process measurement path over the actually inserted keys; the pseudo-bitmap probe declared as such (k=4 hash tests per query, characteristic early-exit latency on mismatch) remains only in the in-process segment-timing scan. & Bloom 1970 (CACM 13:422-426) \\
\texttt{Cuckoo\allowbreak{}Filter} & Cuckoo filter per Fan et al. Co\allowbreak{}NEXT 2014: bucketed cuckoo hashing with 1-byte fingerprints. Supports deletes and lookup counting (advantage vs Bloom). Probe characteristic: 2 bucket lookups + fingerprint comparison (partial-key cuckoo per Co\allowbreak{}NEXT 2014 \S3: i2 = i1 XOR hash(fp)). Higher insert cost, but flexible mutability. & Fan et al. Co\allowbreak{}NEXT 2014 (DOI:10.1145/\allowbreak{}2674005.2674994) \\
\texttt{Range\allowbreak{}Surf\allowbreak{}Filter} & Su\allowbreak{}RF (Succinct Range Filter) per Zhang et al. SIGMOD 2018: LOUDS-bitmap trie for range queries. Only filter with range semantics (lower\_\allowbreak{}bound/\allowbreak{}upper\_\allowbreak{}bound at byte-prefix level). Probe characteristic: multi-level trie descent with succinct-LOUDS navigation (up to k\allowbreak{}Max\allowbreak{}Depth=6), not exactly 1 hash test, hence latency variable per prefix length. Practical in Rocks\allowbreak{}DB SST optimizations. & P10 (Zhang et al. SIGMOD 2018) \\
\texttt{Xor\allowbreak{}Filter} & Xor filter per Graf+Lemire 2020 (JEA 25): XOR-based filter with approx. 9 bits/\allowbreak{}key (vs. 10 bits/\allowbreak{}key Bloom at the same FP rate). Immutable after build. Probe characteristic: exactly 3 independent table accesses (3 disjoint table thirds h0, h1, h2) + 3-way XOR, branch-free, constant latency without early exit. Safe against optimization-away, smallest memory footprint of all filters. & Graf+Lemire 2020 (JEA 25:Article 1.5, DOI:10.1145/\allowbreak{}3376122) \\
\texttt{None\allowbreak{}Filter} & Pass-through variant without filtering (E14): every probe passes unconditionally (input equals output), 0~bit structure memory and 0 hash tests per probe --- the comparison zero point for the space/time trade-off of the four filters (concept 5, end-append, golden-320-neutral). Registered in the catalogue and conformant, but disabled by default in the build configuration (option \texttt{COMDARE\_\allowbreak{}AXIS\_\allowbreak{}FILTER\_\allowbreak{}ENABLE\_\allowbreak{}NONE}, default OFF), so that the enabled inventory of the axis comprises four building blocks. & Code/\allowbreak{}Baseline \\
\bottomrule
\end{longtable}
\end{scriptsize}

\emph{Contributing sources:} Bloom 1970, Fan et al. Co\allowbreak{}NEXT 2014, Zhang et al. SIGMOD 2018, Graf+Lemire JEA 2020~\cite{bloom1970,fan2014cuckoo,zhang2018surf,graf2020xor}.

\section{T15 Buffering Q1 (Operations Buffer)}

This axis is the \emph{buffering organ (Q1)} (T15) of the living being \texttt{SearchAlgorithm} (cf.\ Section~\ref{sec:sota-instances}).

The axis \texttt{axis\_\allowbreak{}q1\_\allowbreak{}queuing} realizes various buffering strategies for operations in the cache-engine context. It encompasses textbook ADTs (FIFO, LIFO, Append-Only, Priority-Heap), specialized versioning techniques (Delta-Chain, Copy-on-Write, Epoch, Tombstone) and modern lock-free queues (SPSC/\allowbreak{}MPMC following Lamport and Vyukov). The strategies are selected at compile time via the axis registry (the \texttt{mp\_list} of the axis registry, filtered by enable flags) per permutation binary (no runtime switch) and are classifiable by access pattern (sequential, ordered, cyclic, versioned, batched, lock-free).

\subsection*{Static sub-axes}

\begin{scriptsize}
\begin{longtable}{@{}>{\raggedright\arraybackslash}p{1.8cm} >{\raggedright\arraybackslash}p{3.2cm} >{\raggedright\arraybackslash}p{8.5cm}@{}}
\toprule
ID & Name & Description \\
\midrule
\endhead%
QS1 & Sequential Access & FIFO/\allowbreak{}LIFO/\allowbreak{}Append-Only: access at head or tail, typical for write-coalescing + Mem\allowbreak{}Table buffering \\
QS2 & Ordered Access & Skip-List/\allowbreak{}Priority-Heap: sorted drain, typical for LSM compaction and hot-key eviction \\
QS3 & Cyclic Access & Bounded ring buffer: fixed capacity, Disruptor pattern, wrap-around with overflow handling \\
QS4 & Versioned Access & Delta-Chain/\allowbreak{}Tombstone/\allowbreak{}Co\allowbreak{}W/\allowbreak{}Epoch: multi-version concurrency control, snapshot isolation, reclamation windows \\
QS5 & Batched Operations & Batched-Insert-Buffer: accumulation of operations into sub-batches, SIMD vectorization, cache-locality optimization \\
QS6 & Lock-Free Concurrent & SPSC/\allowbreak{}MPMC\@: CAS loop or wait-free Lamport model, lock-freedom without mutex \\
\bottomrule
\end{longtable}
\end{scriptsize}

\subsection*{Building blocks}

\begin{scriptsize}
\begin{longtable}{@{}>{\raggedright\arraybackslash}p{3.0cm} >{\raggedright\arraybackslash}p{8.0cm} >{\raggedright\arraybackslash}p{2.5cm}@{}}
\toprule
Variant & Description & Source \\
\midrule
\endhead%
\texttt{No\allowbreak{}Buffer} & Passthrough/\allowbreak{}identity (baseline): \texttt{put()} is a no-op, \texttt{get()} always returns \texttt{std::nullopt}. Useful as comparison point and in classic B+/\allowbreak{}ART scenarios with direct leaf access. Family Q01, QS1 \texttt{sequential\_\allowbreak{}access}. & Code/\allowbreak{}Baseline \\
\texttt{Append\allowbreak{}Only\allowbreak{}Buffer} & Append-only linear buffer (LSM Mem\allowbreak{}Table, Bw-Tree delta): monotonically growing \texttt{std::vector} without deletion in the middle. Optimized for pure append + bulk drain. Family Q02, QS1 \texttt{sequential\_\allowbreak{}access}. Literature: Levandoski/\allowbreak{}Lomet/\allowbreak{}Sengupta Bw-Tree ICDE 2013, O'Neil LSM 1996. & Code/\allowbreak{}Baseline \\
\texttt{FIFOQueue\allowbreak{}Buffer} & FIFO queue (\texttt{std::deque}): first-in-first-out via \texttt{deque::push\_\allowbreak{}back} /\allowbreak{} \texttt{deque::pop\_\allowbreak{}front}. Unbounded, can grow efficiently in the middle. Textbook ADT, typical for write-coalescing. Family Q03, QS1 \texttt{sequential\_\allowbreak{}access}. & Code/\allowbreak{}Baseline \\
\texttt{LIFOStack\allowbreak{}Buffer} & LIFO stack (\texttt{std::vector}): last-in-first-out via \texttt{vector::push\_\allowbreak{}back} /\allowbreak{} \texttt{vector::pop\_\allowbreak{}back}. Unbounded, optimized for hot-path reuse and version tombstones. Family Q04, QS1 \texttt{sequential\_\allowbreak{}access}. & Code/\allowbreak{}Baseline \\
\texttt{Bounded\allowbreak{}Ring\allowbreak{}Buffer} & Bounded ring buffer (Disruptor pattern): fixed capacity, cyclic access with wrap-around. Overflow: drop\_\allowbreak{}oldest (default). \texttt{iterable\_\allowbreak{}aspect\_\allowbreak{}t}: capacity as runtime loop. Family Q05, QS3 \texttt{cyclic\_\allowbreak{}access}. Literature: LMAX Disruptor 2011. & LMAX 2011 \\
\texttt{Priority\allowbreak{}Heap\allowbreak{}Buffer} & Priority heap (\texttt{std::priority\_\allowbreak{}queue}): max-heap, \texttt{get()} returns the highest priority first. First strategy with \texttt{supports\_\allowbreak{}priority\_\allowbreak{}ordering()}=true. Unbounded, textbook ADT (Williams 1964). Application: LRU approximation, hot-key promotion, eviction queues. Family Q06, QS2 \texttt{ordered\_\allowbreak{}access}. & Code/\allowbreak{}Baseline \\
\texttt{Delta\allowbreak{}Chain\allowbreak{}Buffer} & Delta-chain buffer (Bw-Tree, Levandoski 2013): first versioned strategy (\texttt{is\_\allowbreak{}versioned()}=true). Each \texttt{put()} appends a delta record with ascending version ID to the top of the chain. \texttt{get()} LIFO drain with version ID\@. Append-versioning paradigm. Family Q07, QS4 \texttt{versioned\_\allowbreak{}access}. & Bw-Tree (Levandoski et al.\ ICDE 2013)~\cite{levandoski2013bwtree} \\
\texttt{Skiplist\allowbreak{}Buffer} & Skip-list buffer (sorted buffer, Rocks\allowbreak{}DB-style): \texttt{put()} inserts in sorted order via \texttt{std::set}. \texttt{get()} ASCENDING drain (minimum first), typical for LSM-compact output. Difference from \texttt{Priority\allowbreak{}Heap\allowbreak{}Buffer}: skiplist=ascending, heap=descending. Literature: Pugh CACM 1990. Family Q08, QS2 \texttt{ordered\_\allowbreak{}access}. & Pugh 1990, Rocks\allowbreak{}DB \\
\texttt{Tombstone\allowbreak{}Buffer} & Tombstone-marker buffer (LSM + MVCC): \texttt{put()} creates a live entry, deleted entries remain as tombstones (\texttt{std::nullopt}) until the compaction run. \texttt{get()} skips tombstones FIFO-like, \texttt{is\_\allowbreak{}versioned()}=true (erase-versioning, unlike delta-chain append-versioning). \texttt{size()} counts live slots only. Family Q09, QS4 \texttt{versioned\_\allowbreak{}access}. Literature: O'Neil LSM-tree 1996. & O'Neil 1996 \\
\texttt{Copy\allowbreak{}On\allowbreak{}Write\allowbreak{}Buffer} & Copy-on-write snapshot buffer (persistent data structures): \texttt{put()}/\allowbreak{}\texttt{get()} increments the snapshot version, earlier snapshots remain logically preserved via \texttt{shared\_\allowbreak{}ptr}. Snapshot-versioning paradigm (unlike delta-chain append, tombstone erase). \texttt{is\_\allowbreak{}versioned()}=true. Family Q10, QS4 \texttt{versioned\_\allowbreak{}access}. Literature: Driscoll/\allowbreak{}Sarnak/\allowbreak{}Sleator/\allowbreak{}Tarjan JCSS 1989. & Driscoll et al.\ 1989 \\
\texttt{Epoch\allowbreak{}Buffer} & Epoch-based buffer /\allowbreak{} QSBR (quiescent-state-based reclamation): entries accumulated in the current epoch. After the \texttt{epoch\_\allowbreak{}threshold} transition => the earlier epoch can be reclaimed. \texttt{iterable\_\allowbreak{}aspect\_\allowbreak{}t}: \texttt{epoch\_\allowbreak{}threshold} as runtime parameter. \texttt{is\_\allowbreak{}versioned()}=true (reclamation-window versioning). Family Q11, QS4 \texttt{versioned\_\allowbreak{}access}. Literature: Mc\allowbreak{}Kenney RCU OLS 2001, Masstree Euro\allowbreak{}Sys 2012, SMART OSDI 2023. & P29 (RCU OLS 2001), Masstree 2012 \\
\texttt{Batched\allowbreak{}Insert\allowbreak{}Buffer} & Batched-insert buffer (OLAP indexing, ART bulk insert): \texttt{put()} buffers into sub-batches of size \texttt{batch\_\allowbreak{}size\_\allowbreak{}}, complete batches are handed over (bulk drain). Goal: amortized insert cost through bulk operations, SIMD vectorization, cache locality. First QS5 strategy. Family Q12, QS5 \texttt{batched\_\allowbreak{}access}. Literature: Leis ART ICDE 2013 §6. & P01 (Leis ART ICDE 2013) \\
\texttt{Lock\allowbreak{}Free\allowbreak{}SPSCBuffer} & Lock-free SPSC ring queue (Lamport 1977, 1983): single-producer/\allowbreak{}single-consumer. \texttt{put()} wait-free (1 atomic.store + 1 atomic.load), \texttt{get()} wait-free. No CAS needed thanks to disjoint modified sets. \texttt{Progress\allowbreak{}Guarantee::Lock\allowbreak{}Free} (even wait-free). \texttt{iterable\_\allowbreak{}aspect\_\allowbreak{}t}: capacity. Overflow: dropping. Mandatory contract: exactly 1 producer + 1 consumer. Family Q13a, QS6 \texttt{lock\_\allowbreak{}free\_\allowbreak{}access}. & Lamport 1977, 1983 \\
\texttt{Lock\allowbreak{}Free\allowbreak{}MPMCBuffer} & Lock-free bounded MPMC queue (Vyukov 2010--11): multi-producer/\allowbreak{}multi-consumer via per-cell atomic sequence number and CAS loop. \texttt{Cell.sequence} == pos (producer writes), then pos+1 (consumer reads). Retries on contention (wait-free per op \emph{not} guaranteed). Capacity \emph{must} be a power of 2. \texttt{iterable\_\allowbreak{}aspect\_\allowbreak{}t}. Family Q13b, QS6 \texttt{lock\_\allowbreak{}free\_\allowbreak{}access}. \texttt{is\_\allowbreak{}original\_\allowbreak{}eligible}=true (Vyukov BSD-2). & Vyukov 1024cores 2010--11, Michael/\allowbreak{}Scott PODC 1996 \\
\texttt{Original\allowbreak{}Lock\allowbreak{}Free\allowbreak{}Mpmc\allowbreak{}Concurrent\allowbreak{}Queue} & \texttt{moodycamel::Concurrent\allowbreak{}Queue} (Cameron Desrochers, block-based MPMC): proof wrapper for original code with original-code linking (2/\allowbreak{}6 methods original: put, get; 4/\allowbreak{}6 gaps: emplace, peek\_\allowbreak{}front, peek\_\allowbreak{}back, clear are re-impl). \texttt{is\_\allowbreak{}original\_\allowbreak{}module()}=false because of the gap fillings. Header-only, BSD-2/\allowbreak{}BSL-1.0 dual. Analogous to \texttt{Lock\allowbreak{}Free\allowbreak{}MPMCBuffer}, but with an actual submodule binding. Family Q15, QS6 \texttt{lock\_\allowbreak{}free\_\allowbreak{}access}. \texttt{is\_\allowbreak{}original\_\allowbreak{}eligible}=true (moodycamel BSD-2). & moodycamel 2014, BSD-2/\allowbreak{}BSL-1.0 \\
\bottomrule
\end{longtable}
\end{scriptsize}

The family numbers run Q01--Q13b and Q15; Q14 is not assigned in the catalogue, the two lock-free queues share family Q13 (Q13a SPSC, Q13b MPMC), the original wrapper carries Q15.

\emph{Contributing sources:} LMAX 2011, P01 (Leis ART ICDE 2013), Bw-Tree (Levandoski et al.\ ICDE 2013), P29 (RCU OLS 2001), Lamport 1977, 1983, Vyukov 1024cores 2010--11, Michael/\allowbreak{}Scott PODC 1996, Pugh CACM 1990, O'Neil LSM-tree 1996, Driscoll et al.\ 1989, Masstree Euro\allowbreak{}Sys 2012, SMART OSDI 2023, Williams 1964, moodycamel 2014~\cite{thompson2011disruptor,leis2013art,levandoski2013bwtree,mckenney2001rcu,lamport_spsc,lamport1983specifying,vyukov_mpmc,michael1996queue,pugh1990skiplist,oneil1996lsm,driscoll1989persistent,mao2012masstree,luo2023smart,williams1964heapsort,desrochers_concurrentqueue}.

\section{T16 Buffering Q2 (Flush Policy)}

This axis is the \emph{buffering organ (Q2)} (T16) of the living being \texttt{SearchAlgorithm} (cf.\ Section~\ref{sec:sota-instances}).

Axis Q2 defines automatic flush strategies for write buffers. It decides \emph{when} a filled buffer propagates its output to the next stage --- after each operation (Eager), never (Lazy), on threshold excess (Watermark), after a time window (Timed) or adaptively by workload frequency (EWMA-controlled LSM style). This is central to the memory/\allowbreak{}throughput trade-off and the batching efficiency in cache engines.

\subsection*{Static sub-axes}
\begin{scriptsize}
\begin{longtable}{@{}>{\raggedright\arraybackslash}p{1.8cm} >{\raggedright\arraybackslash}p{3.2cm} >{\raggedright\arraybackslash}p{8.5cm}@{}}
\toprule
ID & Name & Description \\
\midrule
\endhead%
FS1 & Event-Triggered & Flush is triggered by explicit events (put/\allowbreak{}get operation or eviction). Encompasses Eager\allowbreak{}Flush (after every operation) and Lazy\allowbreak{}Flush (only on eviction). \\
FS2 & Threshold-Triggered & Flush is triggered when the buffer fill exceeds a threshold (default 75\%). Reduces unnecessary flushes under sparse traffic. \\
FS3 & Time-Triggered & Flush is triggered periodically after a time window (micro-batching), independent of the buffer fill. Classic Hadoop/\allowbreak{}Spark/\allowbreak{}Kafka pattern for bounded-latency guarantees. \\
FS4 & Adaptive-Triggered & Flush threshold is learned adaptively from the workload frequency via EWMA\@. High write rate $\rightarrow$ flush earlier (60\% threshold); low rate $\rightarrow$ flush later (95\% threshold). First self-adaptive Q2 strategy. \\
\bottomrule
\end{longtable}
\end{scriptsize}

\subsection*{Building blocks}
\begin{scriptsize}
\begin{longtable}{@{}>{\raggedright\arraybackslash}p{2.6cm} >{\raggedright\arraybackslash}p{6.5cm} >{\raggedright\arraybackslash}p{4.4cm}@{}}
\toprule
Variant & Description & Source \\
\midrule
\endhead%
\texttt{Eager\allowbreak{}Flush} & Pessimistic write-through policy: flushes completely (Full\allowbreak{}Flush) after \emph{every} \texttt{put()} operation. Lowest latency, but maximum flush frequency and highest throughput overhead --- classic synchronous-write pattern. & Textbook Write-Through Policy /\allowbreak{} Code Baseline \\
\texttt{Lazy\allowbreak{}Flush} & Optimistic write-back policy: \emph{never} flushes automatically, only on explicit eviction from the buffer. Maximum batching, minimum flush frequency --- classic deferred-write pattern. \texttt{should\_\allowbreak{}flush()} \emph{always} returns No\allowbreak{}Flush. & Textbook Write-Back Policy /\allowbreak{} Code Baseline \\
\texttt{Watermark\allowbreak{}Flush} & Threshold-controlled flush policy (default 75\% capacity). Iterable aspect with 5 variants (50\%, 65\%, 75\%, 85\%, 95\%) for runtime permutation. Related to the LSM-tree memtable trigger (O'Neil 1996). & LSM-Tree O'Neil 1996 (related) /\allowbreak{} Code Implementation \\
\texttt{Timed\allowbreak{}Flush} & Time-window-based micro-batching policy (iterable aspect: 10/\allowbreak{}100/\allowbreak{}1000/\allowbreak{}10000 ms). Inspired by Kafka (Net\allowbreak{}DB 2011) and Spark D-Streams (SOSP 2013) for bounded-latency guarantees. Independent of the buffer fill. & Kafka Net\allowbreak{}DB 2011 /\allowbreak{} Spark SOSP 2013 /\allowbreak{} Code Implementation \\
\texttt{Adaptive\allowbreak{}Lsm\allowbreak{}Flush} & Self-adaptive watermark policy with EWMA over the burst rate. Threshold interpolates between 60\% (high write rate) and 95\% (low rate). Conceptually Rocks\allowbreak{}DB Dynamic Leveled Compaction + O'Neil LSM 1996; the EWMA logic is CE-specific. First Q2 strategy with \texttt{is\_\allowbreak{}adaptive=true}. & Rocks\allowbreak{}DB Dynamic Leveled Compaction /\allowbreak{} O'Neil LSM-Tree 1996 /\allowbreak{} Code Implementation \\
\bottomrule
\end{longtable}
\end{scriptsize}

\emph{Contributing sources:} O'Neil et al.\ 1996 (LSM-Tree), Kafka (Net\allowbreak{}DB 2011), Spark D-Streams (SOSP 2013)~\cite{oneil1996lsm,kreps2011kafka,zaharia2013dstreams}.

\section{T17 Persistence Target (Persistence-Target)}

This axis is the \emph{persistence-target organ} (T17) of the living being \texttt{SearchAlgorithm} (cf.\ Section~\ref{sec:sota-instances}); it is the eighteenth organ slot (T17).

The axis makes the storage target of the structure explicit as an organ slot of its own instead of implicitly carrying an in-memory assumption. In the golden reference catalogue it is pinned to its in-memory representative: seventeen axes with two expressions each and \texttt{persistence\_\allowbreak{}target} with one yield $2^{17} = 131072$ reference identities per system permutation (simplifying, superseded assumption of two configurations per axis; regression detector, not an upper bound); every \texttt{binary\_\allowbreak{}id} nevertheless carries the eighteenth segment \texttt{persistence\_\allowbreak{}target=\allowbreak{}persistence\_\allowbreak{}memory\_\allowbreak{}only}. The cardinality of a complete measurement series is different from that: it is not a static number but is counted per profile from the Director walk, as an upper bound from the configurations selected per axis in the XML experiment, known only at execution time (Section~\ref{sec:mess-chain}).

\subsection*{Building blocks}
\begin{scriptsize}
\begin{longtable}{@{}>{\raggedright\arraybackslash}p{3.0cm} >{\raggedright\arraybackslash}p{8.5cm} >{\raggedright\arraybackslash}p{2.0cm}@{}}
\toprule
Variant & Description & Source \\
\midrule
\endhead%
\texttt{Memory\allowbreak{}Only\allowbreak{}Target} & In-memory storage as the default target; the only enabled representative of the canonical measurement catalogue. No device write-back path (\texttt{has\_\allowbreak{}device\_\allowbreak{}writeback\_\allowbreak{}path()=false}). & Code/\allowbreak{}Baseline \\
\texttt{Disk\allowbreak{}Writeback\allowbreak{}Target} & Disk write-back target; fully present as a building block (type, registry, sub-axes PT1--PT2), but switched off per option in the canonical catalogue, so that the axis remains single-element there. It is at the same time the carrier of the organ meta-meta axis \texttt{disk\_\allowbreak{}io} (Section~\ref{sec:app-blocks-disk-io}), whose stamp suffix attaches only to compositions with this target. & Code/\allowbreak{}Baseline \\
\bottomrule
\end{longtable}
\end{scriptsize}

\emph{Contributing sources:} --- (CE baseline, own building blocks).

\section{Organ Meta-Meta Axis: Disk I/O (\texttt{disk\_io})}\label{sec:app-blocks-disk-io}

This quantity is the only organ meta-meta axis of the organ realm (meta-meta in the sense of the M3 level, cf.\ the footnote at the beginning of this appendix). It is no nineteenth composition slot: the composition stays at eighteen organ axes, and the axis never enters the \texttt{binary\_\allowbreak{}id} (\texttt{binary\_\allowbreak{}id=never}, \texttt{stage=ct}, \texttt{baustein\_\allowbreak{}count=1}); the organ registry carries it as category \texttt{organ\_\allowbreak{}meta\_\allowbreak{}meta} beside the eighteen slots.

The axis carries the stamping identity of the member \texttt{disk\_\allowbreak{}io}: its axis code version (\texttt{1.0.0.c}) attaches as a bracket appendage to the end of the organ line (member [3] of the animal fingerprint), and exclusively for compositions with the persistence target \texttt{Disk\allowbreak{}Writeback\allowbreak{}Target} (carrier axis \texttt{persistence\_\allowbreak{}target}, carrier value \texttt{persistence\_\allowbreak{}disk\_\allowbreak{}writeback}); compositions with \texttt{Memory\allowbreak{}Only\allowbreak{}Target} remain byte-identical, their suffix is empty and produces no empty bracket pair. The form follows the system model: the organ line of an I/O carrier ends in \texttt{;[disk\_\allowbreak{}io=code@1.0.0.c]}. The binding strategy type $\rightarrow$ meta-meta lives in the ABI layer, not in the axis header; carriers of the selection are the eighteen existing slot aliases of the composition. For the channel machinery, \texttt{disk\_\allowbreak{}io} is declared as a family channel of its own; the channel itself is declared as a concept but not implemented.
Thus the stamp mechanics of the axis are implemented, the channel construction is not; also not implemented is the profile connection: the profile parser does not read the registry block of the axis, and a reference to it in organ position is rejected by the profile validation, so that the axis is not selectable via experiment profile and is switched off; in the golden reference catalogue its suffix remains empty (Section~\ref{sec:axis-triad}).

\subsection*{Dynamic sub-axes}

\begin{scriptsize}
\begin{longtable}{@{}>{\raggedright\arraybackslash}p{1.8cm} >{\raggedright\arraybackslash}p{3.2cm} >{\raggedright\arraybackslash}p{8.5cm}@{}}
\toprule
ID & Parameter & Description \\
\midrule
\endhead%
\texttt{disk\_\allowbreak{}io} & Provisioning stage of the disk I/O & \texttt{stage=runtime}, \texttt{value\_\allowbreak{}type=token}, \texttt{option\_\allowbreak{}count=1}. Only option \texttt{staging}: \texttt{Disk\allowbreak{}Writeback\allowbreak{}Target} measures exclusively the write-back preparation (no device write path, \texttt{has\_\allowbreak{}device\_\allowbreak{}writeback\_\allowbreak{}path()=false}); a value \texttt{device} follows only with the actual device path --- the registry contains only occupied options. The family spans its sub-axis of the same name (sub-axis label = axis label). \\
\bottomrule
\end{longtable}
\end{scriptsize}

\subsection*{Building blocks}

\begin{scriptsize}
\begin{longtable}{@{}>{\raggedright\arraybackslash}p{3.0cm} >{\raggedright\arraybackslash}p{8.5cm} >{\raggedright\arraybackslash}p{2.0cm}@{}}
\toprule
Variant & Description & Source \\
\midrule
\endhead%
\texttt{Disk\allowbreak{}Io\allowbreak{}Organ\allowbreak{}Meta\allowbreak{}Meta} & Building block \texttt{disk\_\allowbreak{}io}, enabled (\texttt{enabled=true}); satisfies the concept \texttt{Organ\allowbreak{}Meta\allowbreak{}Meta\allowbreak{}Axis} (category directory \texttt{organ\_\allowbreak{}axes/\allowbreak{}organ\_\allowbreak{}meta\_\allowbreak{}meta/}, root model is the SIMD family of the system hub); axis kind \texttt{organ\_\allowbreak{}meta\_\allowbreak{}meta}, explicitly no organ main axis (\texttt{static\_\allowbreak{}assert}). & Code/\allowbreak{}Baseline \\
\bottomrule
\end{longtable}
\end{scriptsize}

\emph{Contributing sources:} --- (CE baseline, own building block).

\section{System Sphere: Target ISA (\texttt{target\_isa})}

This quantity is the system main axis \texttt{target\_isa} --- itself a complex axis (cf.\ Sections~\ref{sec:axis-triad} and~\ref{sec:complex-axis}).

The axis specifies the target CPU instruction set architecture as a compile-time constant and thereby enables cross-platform optimizations for four dominant server architectures: x86-64 (Intel/\allowbreak{}AMD), AArch64 (ARM), Power-ISA (IBM) and RISC-V (open source). Each variant defines a hardware-specific SIMD baseline (SSE2, NEON, VSX, scalar) and thus characterizes the vector acceleration potential for cache-engine operations such as \texttt{simd\_\allowbreak{}field\_\allowbreak{}sum} (accumulation of 32-bit word fields).

\subsection*{Static sub-axes}

\begin{scriptsize}
\begin{longtable}{@{}>{\raggedright\arraybackslash}p{1.8cm} >{\raggedright\arraybackslash}p{3.2cm} >{\raggedright\arraybackslash}p{8.5cm}@{}}
\toprule
ID & Name & Description \\
\midrule
\endhead%
IS1 & Word Size & CPU data-path width: 32-bit, 64-bit or 128-bit. In the cache engine standardized to 64-bit for all four ISA profiles. \\
IS2 & CPU Vendor/\allowbreak{}Family & Instruction set family: \texttt{intel\_\allowbreak{}amd} (x86-64), \texttt{arm} (AArch64), \texttt{ibm} (Power-ISA), \texttt{open} (RISC-V). Determines SIMD extension compatibility: SSE2/\allowbreak{}AVX2/\allowbreak{}AVX-512 (x86-64), NEON/\allowbreak{}SVE/\allowbreak{}SVE2 (AArch64), VSX (Power), RVV (RISC-V). \\
IS3 & ZIH Cluster Affinity & Production machines of the ZIH of TU Dresden~\cite{zih_hpc_compendium_hardware}: Barnard (Intel Xeon Platinum 8470, Sapphire Rapids, x86-64), Capella (AMD EPYC 9334, x86-64), the Grace Hopper Superchip test system (72 Arm Neoverse V2 cores, AArch64; test system with limited availability~\cite{zih_hpc_compendium_gh200}), N/\allowbreak{}A (RISC-V and Power: not among the current HPC resources of the ZIH compendium). Enables hardware-specific experiments with production cluster profiles. \\
\bottomrule
\end{longtable}
\end{scriptsize}

Beyond these three design dimensions, the build-driving system registry carries five sub-axes on the complex wrapper: two compile-time ones (\texttt{target\_\allowbreak{}isa\_\allowbreak{}complex} and \texttt{scheduling}, each \texttt{stage=ct} with \texttt{kind=fixed\_\allowbreak{}enum\_\allowbreak{}tuple}) and three dynamic ones (\texttt{numa\_\allowbreak{}node}, \texttt{page} and \texttt{core\_\allowbreak{}class}). The first bundles the named recombinations of the target machines (\texttt{complex\_\allowbreak{}count=2}) with three fixed members each.

\begin{scriptsize}
\begin{longtable}{@{}>{\raggedright\arraybackslash}p{3.0cm} >{\raggedright\arraybackslash}p{3.6cm} >{\raggedright\arraybackslash}p{6.9cm}@{}}
\toprule
Recombination & Member & Declaration \\
\midrule
\endhead%
\texttt{prod1\_\allowbreak{}zen5} (\texttt{isa=x86\_\allowbreak{}64}) & \texttt{ram\_\allowbreak{}frequency\_\allowbreak{}mhz} & \texttt{declared=false} --- no declared figure. \\
\texttt{prod1\_\allowbreak{}zen5} & \texttt{cas\_\allowbreak{}latency\_\allowbreak{}cl} & \texttt{declared=false} --- no declared figure. \\
\texttt{prod1\_\allowbreak{}zen5} & \texttt{cpu\_\allowbreak{}fabrication} & \texttt{declared=true}, \texttt{value=Authentic\allowbreak{}AMD/\allowbreak{}26/\allowbreak{}68/\allowbreak{}0} (vendor/\allowbreak{}family/\allowbreak{}model/\allowbreak{}stepping). \\
\texttt{prod2\_\allowbreak{}alder\_\allowbreak{}lake} (\texttt{isa=x86\_\allowbreak{}64}) & \texttt{ram\_\allowbreak{}frequency\_\allowbreak{}mhz} & \texttt{declared=true}, \texttt{value=4800}, \texttt{declaration\_\allowbreak{}source=\allowbreak{}dmidecode 2026-07}. \\
\texttt{prod2\_\allowbreak{}alder\_\allowbreak{}lake} & \texttt{cas\_\allowbreak{}latency\_\allowbreak{}cl} & \texttt{declared=false} --- no declared figure. \\
\texttt{prod2\_\allowbreak{}alder\_\allowbreak{}lake} & \texttt{cpu\_\allowbreak{}fabrication} & \texttt{declared=true}, \texttt{value=Genuine\allowbreak{}Intel/\allowbreak{}6/\allowbreak{}151/\allowbreak{}2}. \\
\bottomrule
\end{longtable}
\end{scriptsize}

\emph{Note:} the values originate from \texttt{k\allowbreak{}Declared\allowbreak{}Machines}, which in turn mirrors the \texttt{<machines>} block of the user XML\@. \texttt{declaration\_\allowbreak{}source} preserves the fixed origin of a declared figure as a note, not as a provenance level; run-time levels never appear in this XML\@.

The second compile-time sub-axis \texttt{scheduling} carries five fixed members. It hangs as a \texttt{sub\_\allowbreak{}axis} on the \texttt{target\_\allowbreak{}isa} complex wrapper and is, by the design decision of this thesis, at the same time a permutation axis of the planner --- sub-axis and permutation axis are not mutually exclusive. In the build offering it is \texttt{stage=ct} with a fixed five-tuple; the permutation through the planner is fixed as a design but not implemented.

\begin{scriptsize}
\begin{longtable}{@{}>{\raggedright\arraybackslash}p{4.0cm} >{\raggedright\arraybackslash}p{2.4cm} >{\raggedright\arraybackslash}p{7.1cm}@{}}
\toprule
Member & Registry value & Description \\
\midrule
\endhead%
\texttt{worker\_\allowbreak{}pool\_\allowbreak{}layout} & \texttt{value\_\allowbreak{}ordinal=0} & Distribution of the worker pool across the cores. \\
\texttt{simd\_\allowbreak{}worker\_\allowbreak{}count\_\allowbreak{}limit} & \texttt{value=2} & Upper bound of concurrent SIMD workers. \\
\texttt{hetero\_\allowbreak{}core\_\allowbreak{}dispatch} & \texttt{value\_\allowbreak{}ordinal=0} & Dispatch across heterogeneous cores (P and E cores). \\
\texttt{co\_\allowbreak{}routine\_\allowbreak{}strategy} & \texttt{value\_\allowbreak{}ordinal=0} & Coroutine strategy of job processing. \\
\texttt{batch\_\allowbreak{}granularity} & \texttt{value\_\allowbreak{}ordinal=0} & Granularity of batch processing. \\
\bottomrule
\end{longtable}
\end{scriptsize}

\emph{Note:} the ordinals are raw enum values without a compile-time label template; the identifier strings are compile-coupled to the accessor calls of \texttt{concepts::\allowbreak{}scheduling\_\allowbreak{}strategy.hpp}.

\subsection*{Dynamic sub-axes}

\begin{scriptsize}
\begin{longtable}{@{}>{\raggedright\arraybackslash}p{1.8cm} >{\raggedright\arraybackslash}p{3.2cm} >{\raggedright\arraybackslash}p{8.5cm}@{}}
\toprule
ID & Parameter & Description \\
\midrule
\endhead%
\texttt{numa\_\allowbreak{}node} & NUMA node & \texttt{stage=runtime}, \texttt{value\_\allowbreak{}type=token}, \texttt{option\_\allowbreak{}source=machine\_\allowbreak{}resolved}. Mirror of the organ member \texttt{alloc\_\allowbreak{}hw.numa\_\allowbreak{}node}; the admissible values are the node identifiers of the concrete machine and are not knowable at build time --- hence a declared option source instead of an option catalogue. \\
\texttt{page} & Page-size release & \texttt{stage=runtime}, \texttt{value\_\allowbreak{}type=token}, \texttt{option\_\allowbreak{}source=machine\_\allowbreak{}resolved}. The release of page sizes (4k/\allowbreak{}2m/\allowbreak{}1g) as a system capability; enforcement lives on the organ side in the \texttt{Alloc\allowbreak{}Page\allowbreak{}Hint} NTTP, this axis is the system side of the release. \\
\texttt{core\_\allowbreak{}class} & Core class & \texttt{stage=runtime}, \texttt{value\_\allowbreak{}type=token}, \texttt{option\_\allowbreak{}source=machine\_\allowbreak{}resolved}. The class of the executing core (P or E core on hybrid processors) as the execution locality of the measurement --- a sibling of \texttt{numa\_\allowbreak{}node} and \texttt{page}, which capture the memory locality; distinct from \texttt{scheduling}, which contains the policy fixed at compile time: \texttt{core\_\allowbreak{}class} says whether and which core classes the machine carries. Admissible values are machine-resolved; the separate measurement on P and E cores is thus named in the catalogue as the system side of core pinning, the measurement itself is performed by the measurement system (Section~\ref{sec:measurement-system}). \\
\bottomrule
\end{longtable}
\end{scriptsize}

\emph{Name collision (guarded at compile time):} the sub-axis is called \texttt{page} and is not \texttt{page\_\allowbreak{}type}. \texttt{axis\_\allowbreak{}01\_\allowbreak{}page\_\allowbreak{}type} is the tree-node \texttt{Page\allowbreak{}Kind} of the organ world and an entirely different axis; the \texttt{static\_\allowbreak{}assert} in \texttt{target\_\allowbreak{}isa\_\allowbreak{}sub\_\allowbreak{}axes.hpp} records this. The building-block catalogue of the page type is in the conserved offering catalogue \enquote{Build-PG}.

\subsection*{Building blocks}

\begin{scriptsize}
\begin{longtable}{@{}>{\raggedright\arraybackslash}p{2.4cm} >{\raggedright\arraybackslash}p{7.6cm} >{\raggedright\arraybackslash}p{3.5cm}@{}}
\toprule
Variant & Description & Source \\
\midrule
\endhead%
\texttt{Amd64\allowbreak{}Isa} & x86-64 /\allowbreak{} Intel 64 ISA (AMD/\allowbreak{}Intel server+desktop). Dominant server platform with SSE2 SIMD baseline (4 uint32 lanes per 128-bit register). Implements SSE2 vectorization via \texttt{\_\allowbreak{}mm\_\allowbreak{}loadu\_\allowbreak{}si128} + \texttt{\_\allowbreak{}mm\_\allowbreak{}add\_\allowbreak{}epi32} on native x86-64 hosts; scalar fallback on non-x86 build hosts. ZIH presence~\cite{zih_hpc_compendium_hardware}: Barnard (Intel Xeon Platinum 8470, Sapphire Rapids), Capella (AMD EPYC 9334), Alpha Centauri (AMD EPYC 7352), Romeo (AMD EPYC 7702). & Vendor spec (AMD 2000 /\allowbreak{} Intel ongoing) \\
\texttt{Aarch64\allowbreak{}Isa} & ARMv8-A 64-bit ISA (ARM Holdings). Growing server relevance (Apple M-Series, AWS Graviton, Ampere Altra). SIMD baseline: NEON (128-bit, 4 uint32 lanes). Cache-engine implementation: 4-way unrolled scalar fallback (models NEON width without ARM intrinsics on x86 build hosts); native NEON vectorization on ARMv8 hosts. ZIH presence: Grace Hopper Superchip test system (72 Arm Neoverse V2 cores; test system with limited availability~\cite{zih_hpc_compendium_gh200}). & Vendor spec (ARM Ltd., since 2011/\allowbreak{}2013) \\
\texttt{Power\allowbreak{}Pc\allowbreak{}Isa} & Power ISA (POWER9/\allowbreak{}10, ppc64le little-endian). IBM Power Systems (AC922, IC922, S1022). SIMD baseline: VSX (Vector-Scalar Extension, 128-bit, 4 uint32 lanes). Cache-engine implementation: 4-way unrolled scalar fallback (models VSX width without Power intrinsics on x86 build hosts). Relevant for HPC workloads with NVLink GPU connectivity (AC922). ZIH presence: carried in the ZIH compendium only as a decommissioned archive system (Power9), not among the current HPC resources. & Vendor/\allowbreak{}standards spec (IBM/\allowbreak{}Open\allowbreak{}POWER, v3.0 2017, v3.1 2020) \\
\texttt{Risc\allowbreak{}VIsa} & RISC-V RV64\allowbreak{}GC open-source ISA (RV64 General + Compressed). Growing relevance (Si\allowbreak{}Five, Alibaba Xuantie, Andes, ETH Zurich cores). SIMD baseline: \emph{none} (RVV is optional, not part of the RV64\allowbreak{}GC baseline). Cache-engine implementation: deliberately plain scalar path (not unrolled) -- models the missing vector unit of the base ISA and thus produces a real, slower runtime versus x86-64 and ARM\@. ZIH presence: not listed in the ZIH compendium. & Standards spec (RISC-V Intl.; UCB TR since 2011) \\
\bottomrule
\end{longtable}
\end{scriptsize}

\emph{Registry state:} the build-driving system registry carries two enabled building blocks --- \texttt{x86\_\allowbreak{}64} (native, without triple and march flags) and \texttt{aarch64} (cross compilation with \texttt{triple=\allowbreak{}-target aarch64-linux-gnu} and \texttt{march=\allowbreak{}-march=armv8-a}). The four-ISA catalogue listed above is retained as an \emph{offering} catalogue: Power-ISA and RISC-V are described as building blocks but are not enabled in the build offering.

\emph{Contributing sources:} ---

\section{System Sphere: Operating System (\texttt{operating\_system})}

This quantity is the second system main axis and thus part of the \emph{habitat} and of no organ (cf.\ Section~\ref{sec:axis-triad}); in the canonical, code-guarded order it occupies position two of three.

The main axis carries the \emph{class} identity only: exactly three operating-system concepts (\texttt{baustein\_\allowbreak{}count=3}), together with \texttt{stage=ct} and \texttt{binary\_\allowbreak{}id=never}. The \emph{instance} --- distribution, release, kernel and patch level --- resides entirely in the three sub-axes; entering a distribution as a fourth expression of the main axis would miss that cut. The existing machines (Windows 11 and Windows Server 2022, the Linux Docker matrix, two macOS machines) fall entirely into the three concepts.

\subsection*{Dynamic sub-axes}

\begin{scriptsize}
\begin{longtable}{@{}>{\raggedright\arraybackslash}p{1.8cm} >{\raggedright\arraybackslash}p{3.2cm} >{\raggedright\arraybackslash}p{8.5cm}@{}}
\toprule
ID & Name & Description \\
\midrule
\endhead%
\texttt{os\_\allowbreak{}version} & Distribution/\allowbreak{}release identity & The release identity of the operating-system instance within its family: on Linux \texttt{ID} and \texttt{VERSION\_\allowbreak{}ID} from \texttt{/etc/\allowbreak{}os-release}, on Windows the product line, on macOS the product version. \\
\texttt{kernel} & Kernel release & The kernel release of the instance: on Linux and macOS the \texttt{uname} release, on Windows the NT kernel triple. \\
\texttt{build} & Build and patch level & The build and patch level of the instance \emph{including} the update state (design decision of this thesis: the update state is a constituent of the build identifier). A separate \texttt{update\_\allowbreak{}zustand} is therefore explicitly \emph{not} a fourth sub-axis but a constituent of this one. \\
\bottomrule
\end{longtable}
\end{scriptsize}

\emph{Note:} all three are \texttt{stage=\allowbreak{}runtime} with \texttt{value\_\allowbreak{}type=\allowbreak{}token} and \texttt{option\_\allowbreak{}source=\allowbreak{}machine\_\allowbreak{}resolved}: they have no compile-time option catalogue because their admissible values depend on the concrete machine and are not knowable at build time. The run-time acquisition is implemented in probe files of their own (\texttt{operating\_\allowbreak{}system\_\allowbreak{}probe\_\allowbreak{}\{linux,\allowbreak{}macos,\allowbreak{}windows\}.hpp}); the attachment to the registry declares only the options. The three names are part of the generated registry and thus fixed.

\subsection*{Building blocks}

\begin{scriptsize}
\begin{longtable}{@{}>{\raggedright\arraybackslash}p{2.4cm} >{\raggedright\arraybackslash}p{7.6cm} >{\raggedright\arraybackslash}p{3.5cm}@{}}
\toprule
Variant & Description & Source \\
\midrule
\endhead%
\texttt{linux} & Linux family (\texttt{Linux\allowbreak{}Operating\allowbreak{}System}, \texttt{posix=true}). Covers the entire Docker distribution matrix; the individual distribution is a matter for the sub-axis \texttt{os\_\allowbreak{}version}, not for this expression. & Code/\allowbreak{}Baseline \\
\texttt{windows} & Windows family (\texttt{Windows\allowbreak{}Operating\allowbreak{}System}, \texttt{posix=false}). Covers desktop and server product lines (Windows 11, Windows Server 2022). & Code/\allowbreak{}Baseline \\
\texttt{macos} & macOS family (\texttt{Macos\allowbreak{}Operating\allowbreak{}System}, \texttt{posix=true}). Covers the existing Apple Silicon and Intel machines. & Code/\allowbreak{}Baseline \\
\bottomrule
\end{longtable}
\end{scriptsize}

\emph{Contributing sources:} --- (CE baseline, own building blocks).

\emph{Stamp neutrality:} the three sub-axes are run-time sub-axes and therefore never appear in the binary stamp. Their attachment changes neither the stamp lines nor the build-version suffix, and it must not raise the axis code level of \texttt{operating\_\allowbreak{}system} (check against a version jump): such a jump would shift the SHA512 of all new builds against the existing stock and would force exactly the full rebuild of all legacy binaries that the design decision of this thesis rules out.

\section{System Sphere: Extension Hardware (\texttt{external\_utils})}

This quantity is the third system main axis; it is exclusively the collecting hub of the hardware meta-meta axes --- SIMD and vector extensions, GPU, FPGA and NPU as well as external accelerators (cf.\ Section~\ref{sec:axis-triad}).

The hub carries \emph{only} system meta-metas; the load framework explicitly does not belong to it. Because the meta-meta axes receive no stamp line of their own but appear as a bracket appendage at the end of the respective realm line, this hub closes the system realm: the appendage is folded type-driven from the member type list of the hub (\texttt{External\allowbreak{}Utils\allowbreak{}Hub} $\rightarrow$ \texttt{Simd\allowbreak{}External\allowbreak{}Utils\allowbreak{}Family}), the level resides in the bracket depth and not in the name. In the build-driving offering the axis carries \texttt{baustein\_\allowbreak{}count=1}, \texttt{stage=ct} and \texttt{binary\_\allowbreak{}id=never}.

\subsection*{Dynamic sub-axes}

\begin{scriptsize}
\begin{longtable}{@{}>{\raggedright\arraybackslash}p{1.8cm} >{\raggedright\arraybackslash}p{3.2cm} >{\raggedright\arraybackslash}p{8.5cm}@{}}
\toprule
ID & Name & Description \\
\midrule
\endhead%
\texttt{simd} & Vector extension level & \texttt{stage=runtime}, \texttt{value\_\allowbreak{}type=token}, \texttt{option\_\allowbreak{}count=3}. The concept identifier is at the same time the sub-axis label (notation \texttt{external\_\allowbreak{}utils.simd=\dots}); the per-dialect reflected compiler flags come from the option, not from the concept. \\
\bottomrule
\end{longtable}
\end{scriptsize}

The three options of the sub-axis with their flags per compiler dialect:

\begin{scriptsize}
\begin{longtable}{@{}>{\raggedright\arraybackslash}p{3.0cm} >{\raggedright\arraybackslash}p{5.2cm} >{\raggedright\arraybackslash}p{5.3cm}@{}}
\toprule
Option & Flag (\texttt{gpp} and \texttt{clang}) & Flag (\texttt{msvc}) \\
\midrule
\endhead%
\texttt{no\_\allowbreak{}extension} & --- (empty) & --- (empty) \\
\texttt{avx2} & \texttt{-mavx2} & \texttt{/arch:AVX2} \\
\texttt{avx512} & \texttt{-mavx512f} & \texttt{/arch:AVX512} \\
\bottomrule
\end{longtable}
\end{scriptsize}

\emph{Note:} \texttt{simd} materializes as a flag of the compilation seam and therefore appears in the build-version suffix (cf.\ Section~\ref{sec:stamp-model}) --- never as a run-time type switch and never in the \texttt{binary\_\allowbreak{}id}.

The sub-axis is accompanied by a fine-grained feature catalogue \texttt{simd\_\allowbreak{}feature\_\allowbreak{}catalog} with \texttt{count=23} features, tiered by \texttt{tier}: fourteen features of tier \texttt{avx512}, four of tier \texttt{avx256}, three companion features (\texttt{companion}) and two scalar ones (\texttt{scalar}). The flags are identical throughout for \texttt{gpp} and \texttt{clang}.

\begin{scriptsize}
\begin{longtable}{@{}>{\raggedright\arraybackslash}p{4.4cm} >{\raggedright\arraybackslash}p{5.4cm} >{\raggedright\arraybackslash}p{3.7cm}@{}}
\toprule
Feature (\texttt{cpuinfo}) & Flag (\texttt{gpp} and \texttt{clang}) & Tier (\texttt{tier}) \\
\midrule
\endhead%
\texttt{avx512f} & \texttt{-mavx512f} & \texttt{avx512} \\
\texttt{avx512cd} & \texttt{-mavx512cd} & \texttt{avx512} \\
\texttt{avx512vl} & \texttt{-mavx512vl} & \texttt{avx512} \\
\texttt{avx512dq} & \texttt{-mavx512dq} & \texttt{avx512} \\
\texttt{avx512bw} & \texttt{-mavx512bw} & \texttt{avx512} \\
\texttt{avx512ifma} & \texttt{-mavx512ifma} & \texttt{avx512} \\
\texttt{avx512vbmi} & \texttt{-mavx512vbmi} & \texttt{avx512} \\
\texttt{avx512\_\allowbreak{}vbmi2} & \texttt{-mavx512vbmi2} & \texttt{avx512} \\
\texttt{avx512\_\allowbreak{}vnni} & \texttt{-mavx512vnni} & \texttt{avx512} \\
\texttt{avx512\_\allowbreak{}bitalg} & \texttt{-mavx512bitalg} & \texttt{avx512} \\
\texttt{avx512\_\allowbreak{}vpopcntdq} & \texttt{-mavx512vpopcntdq} & \texttt{avx512} \\
\texttt{avx512\_\allowbreak{}vp2intersect} & \texttt{-mavx512vp2intersect} & \texttt{avx512} \\
\texttt{avx512\_\allowbreak{}bf16} & \texttt{-mavx512bf16} & \texttt{avx512} \\
\texttt{avx512\_\allowbreak{}fp16} & \texttt{-mavx512fp16} & \texttt{avx512} \\
\texttt{avx2} & \texttt{-mavx2} & \texttt{avx256} \\
\texttt{fma} & \texttt{-mfma} & \texttt{avx256} \\
\texttt{f16c} & \texttt{-mf16c} & \texttt{avx256} \\
\texttt{avx\_\allowbreak{}vnni} & \texttt{-mavxvnni} & \texttt{avx256} \\
\texttt{gfni} & \texttt{-mgfni} & \texttt{companion} \\
\texttt{vaes} & \texttt{-mvaes} & \texttt{companion} \\
\texttt{vpclmulqdq} & \texttt{-mvpclmulqdq} & \texttt{companion} \\
\texttt{bmi2} & \texttt{-mbmi2} & \texttt{scalar} \\
\texttt{popcnt} & \texttt{-mpopcnt} & \texttt{scalar} \\
\bottomrule
\end{longtable}
\end{scriptsize}

\emph{Machine signatures (\texttt{count=3}):} the registry records for each target machine the subset of the features actually present --- \texttt{prod1\_\allowbreak{}zen5} (\texttt{host\_\allowbreak{}isa=x86\_\allowbreak{}64}) with 22 features, \texttt{prod2\_\allowbreak{}alder\_\allowbreak{}lake} with 9 and \texttt{odroid\_\allowbreak{}gracemont} with 9. Only the Zen 5 machine has the \texttt{avx512} tier; the other two end at \texttt{avx256} together with the companion and scalar features.

\subsection*{Building blocks}

\begin{scriptsize}
\begin{longtable}{@{}>{\raggedright\arraybackslash}p{2.4cm} >{\raggedright\arraybackslash}p{7.6cm} >{\raggedright\arraybackslash}p{3.5cm}@{}}
\toprule
Variant & Description & Source \\
\midrule
\endhead%
\texttt{simd} & Vector extension family (\texttt{Simd\allowbreak{}External\allowbreak{}Utils\allowbreak{}Family}, \texttt{family\_\allowbreak{}id=simd}, enabled). The only enabled member of the hub; the families for GPU, FPGA and NPU are foreseen in the design but are not part of the build-driving offering. & Code/\allowbreak{}Baseline \\
\bottomrule
\end{longtable}
\end{scriptsize}

\emph{Contributing sources:} --- (CE baseline, own building blocks).

\section[head={System Sphere: Complex Bracket and Build Toolchain}]{System Sphere: Outer Complex Bracket (\texttt{build\_target\_complex}) and Build Toolchain (\texttt{build\_toolchain})}

This quantity is the command-bracket wrapper over the fixed recombination of the three system main axes (\texttt{member\_\allowbreak{}count=3}: \texttt{target\_\allowbreak{}isa}, \texttt{operating\_\allowbreak{}system}, \texttt{external\_\allowbreak{}utils}) and thus closes the system realm block (cf.\ Section~\ref{sec:complex-axis}).

As a bracket it is explicitly \emph{not} a fourth main axis: the canonical system axis order stays at three. Three static exclusion checks (\texttt{static\_\allowbreak{}assert}) record that \texttt{compiler}, \texttt{scheduling} and \texttt{load\_\allowbreak{}framework} are no system main axes. The load framework is not a \emph{member} of this bracket either: it belongs to the measurement realm, whereas the system meta-metas are carried under the hub \texttt{external\_\allowbreak{}utils}. The axis carries \texttt{stage=ct} and \texttt{binary\_\allowbreak{}id=never}.

\subsection*{Static sub-axes}

The sub-axis group \texttt{build\_\allowbreak{}toolchain} (\texttt{sub\_\allowbreak{}axis\_\allowbreak{}count=3}) hangs on the bracket. The registry names the three members flat; the nesting underneath shows the actual parent chain.

\begin{scriptsize}
\begin{longtable}{@{}>{\raggedright\arraybackslash}p{2.2cm} >{\raggedright\arraybackslash}p{3.4cm} >{\raggedright\arraybackslash}p{7.9cm}@{}}
\toprule
Member & Parent reference & Description \\
\midrule
\endhead%
\texttt{compiler} & \texttt{build\_\allowbreak{}target\_\allowbreak{}complex} & \texttt{stage=ct}, \texttt{baustein\_\allowbreak{}count=2}. The compiler choice itself; it is a member of the sub-axis group \texttt{build\_\allowbreak{}toolchain} of the outer bracket. \\
\texttt{opt\_\allowbreak{}level} & \texttt{compiler} & \texttt{stage=runtime}, \texttt{value\_\allowbreak{}type=token}, \texttt{option\_\allowbreak{}count=5}. Optimization level per compiler dialect. \\
\texttt{atomic128} & \texttt{compiler} & \texttt{stage=runtime}, \texttt{value\_\allowbreak{}type=token}, \texttt{option\_\allowbreak{}count=2}. Release of the 128-bit compare-and-swap instruction. \\
\bottomrule
\end{longtable}
\end{scriptsize}

\emph{Note:} the group hangs on \texttt{build\_\allowbreak{}target\_\allowbreak{}complex} and not on \texttt{target\_\allowbreak{}isa} because it describes the \emph{build} arising from the recombination of all three main axes and not a property of the target ISA alone.

\subsection*{Building blocks}

\begin{scriptsize}
\begin{longtable}{@{}>{\raggedright\arraybackslash}p{2.4cm} >{\raggedright\arraybackslash}p{7.6cm} >{\raggedright\arraybackslash}p{3.5cm}@{}}
\toprule
Variant & Description & Source \\
\midrule
\endhead%
\texttt{gcc} & GNU compiler (\texttt{Gcc\allowbreak{}Compiler\allowbreak{}Axis}, driver \texttt{g++-16}, \texttt{supports\_\allowbreak{}fno\_\allowbreak{}gnu\_\allowbreak{}unique=true}). & Code/\allowbreak{}Baseline \\
\texttt{clang} & LLVM compiler (\texttt{Clang\allowbreak{}Compiler\allowbreak{}Axis}, driver \texttt{clang++-22}, \texttt{supports\_\allowbreak{}fno\_\allowbreak{}gnu\_\allowbreak{}unique=false}). & Code/\allowbreak{}Baseline \\
\bottomrule
\end{longtable}
\end{scriptsize}

\emph{Offering and build minimum:} the driver identifiers \texttt{g++-16} and \texttt{clang++-22} are the default drivers of the offering (\texttt{driver\_\allowbreak{}default} of the compiler building blocks); the build itself requires \texttt{g++} $\geq$ 15.3 as the minimum and standard of the measurement machines and deliberately pins no version (newer \texttt{g++} admitted); the CI test cells prefer \texttt{g++-15} where available, and the clang full build runs on clang 22. Clang and MSVC have no version minimum.

The five options of the sub-axis \texttt{opt\_\allowbreak{}level} with their flags per dialect; the column \texttt{ieee754\_\allowbreak{}deterministic} states whether the level leaves the floating-point semantics untouched:

\begin{scriptsize}
\begin{longtable}{@{}>{\raggedright\arraybackslash}p{2.0cm} >{\raggedright\arraybackslash}p{3.4cm} >{\raggedright\arraybackslash}p{2.6cm} >{\raggedright\arraybackslash}p{5.5cm}@{}}
\toprule
Option & Flag (\texttt{gpp} and \texttt{clang}) & Flag (\texttt{msvc}) & \texttt{ieee754\_\allowbreak{}deterministic} \\
\midrule
\endhead%
\texttt{O0} & \texttt{-O0} & \texttt{/Od} & \texttt{true} \\
\texttt{O1} & \texttt{-O1} & \texttt{/O1} & \texttt{true} \\
\texttt{O2} & \texttt{-O2} & \texttt{/O2} & \texttt{true} \\
\texttt{O3} & \texttt{-O3} & \texttt{/O2} & \texttt{true} \\
\texttt{Ofast} & \texttt{-Ofast} & \texttt{/O2} & \texttt{false} \\
\bottomrule
\end{longtable}
\end{scriptsize}

The default of the build is \texttt{O2}; \texttt{O3} is selectable via the XML, and the choice enters the build-version suffix as a field.

The two options of the sub-axis \texttt{atomic128}:

\begin{scriptsize}
\begin{longtable}{@{}>{\raggedright\arraybackslash}p{3.0cm} >{\raggedright\arraybackslash}p{5.2cm} >{\raggedright\arraybackslash}p{5.3cm}@{}}
\toprule
Option & Flag (\texttt{gpp} and \texttt{clang}) & Flag (\texttt{msvc}) \\
\midrule
\endhead%
\texttt{no\_\allowbreak{}cx16} & --- (empty) & --- (empty) \\
\texttt{cx16} & \texttt{-mcx16} & --- (empty) \\
\bottomrule
\end{longtable}
\end{scriptsize}

\emph{Build version identifier:} \texttt{opt\_\allowbreak{}level} and \texttt{simd} materialize as flags of the compilation seam and appear as segments (\texttt{+opt=}, \texttt{+ext=}) in the build-version suffix of the binary --- an ordered eight-segment sequence in which an empty member simply produces no segment (cf.\ Section~\ref{sec:stamp-model}); \texttt{atomic128} likewise materializes as a flag of the compilation seam, but stands not in the suffix but as a member-owned field of the toolchain member in the fingerprint preimage (beside \texttt{vendoropt}, ADR-17 in Appendix~\ref{app:adr}). They never appear in the \texttt{binary\_\allowbreak{}id}.

\emph{Version flag grammar:} the version of every axis algorithm is uniformly three-part; a \texttt{v} prefix is excluded, short forms are forbidden. Compile-time checks enforce the form:

\begin{quote}
\verb+version := UINT '.' UINT '.' UINT [ '.' flag ]*+\\
\verb+flag     := basis [ '{' sub [ '.' sub ]* '}' ] | companion+\\
\verb+basis    := 'c' | 'g' | 'f' | 'n' | 'x128' | 'x256' | 'x512'+\\
\verb+            | 'm64'+\\
\verb+sub      := token [ '{' sub [ '.' sub ]* '}' ]  (recursive)+
\end{quote}

A dot precedes every flag, including the first (\texttt{1.0.0.c}, never \texttt{1.0.0c}); the basis attaches directly to its opening brace, with no dot of its own in between (\texttt{c\{p.e\}}, never \texttt{c.\{p.e\}}). The suffix \texttt{e} denotes an \emph{efficiency core} and is a sub-flag like any other under the basis \texttt{c}; the basis \texttt{c} without a brace is semantically \texttt{c\{p\}}; a sub-flag for experimentally introduced candidate algorithms does not exist. The bases \texttt{c}, \texttt{g}, \texttt{f} and \texttt{n} name the target hardware class (CPU, GPU, FPGA and NPU, of which only \texttt{c} is produced and the other three are reserved), \texttt{x128}, \texttt{x256} and \texttt{x512} the three views onto the same vector register file (XMM, YMM, ZMM), and \texttt{m64} the media basis of the MMX/3DNow! family on the x87-aliased MM registers, whose sub-flags the catalogue check passes only under this basis (\texttt{1.0.0.c.m64\{mmx.mmxext.3dnow\}}, never as bare tokens without a basis). Example form: \texttt{memory\_\allowbreak{}layout=\allowbreak{}SoaMemoryLayout@\allowbreak{}1.0.0.c\{p.e\}.\allowbreak{}x512\{f.vl.bw.dq\}.\allowbreak{}gfni}.

\section[head={Measurement Sphere: Measurement Categories}]{Measurement Sphere: Measurement Categories (\texttt{measurement\_category})}

This quantity belongs to the third realm --- the measurement axes, the equipment of the \emph{laboratory} (cf.\ Section~\ref{sec:axis-triad}).

It carries \texttt{stage=ct}, \texttt{binary\_\allowbreak{}id=never} and \texttt{baustein\_\allowbreak{}count=16}. Its values manifest as columns of the result tables and as the setting label of a run, never in the binary identity. In the chapter model the sixteen categories are the \emph{sub}-axis of the measurement tooling; the measurement registry additionally carries them as an offering block of their own, because the tooling does not invent its categories but selects them from this catalogue. The column \texttt{regime\_\allowbreak{}ordinal} separates the time- and observer-based regime (0) from the counter-based PMC regime (1). The sixteen categories are the offering catalogue of the category axis; the number of measurement sources acquired per axis interface and parameter set is different from that and not a subject of this catalogue (Section~\ref{sec:measurement-system}).

\subsection*{Building blocks}

\begin{scriptsize}
\begin{longtable}{@{}>{\raggedright\arraybackslash}p{5.0cm} >{\raggedright\arraybackslash}p{2.6cm} >{\raggedright\arraybackslash}p{5.9cm}@{}}
\toprule
Category & \texttt{category\_\allowbreak{}ordinal} & Regime (\texttt{regime\_\allowbreak{}ordinal}) \\
\midrule
\endhead%
\texttt{CLU} & 0 & 0 --- time/\allowbreak{}observer \\
\texttt{CACHE\_\allowbreak{}MISS\_\allowbreak{}L1} & 1 & 1 --- PMC \\
\texttt{CACHE\_\allowbreak{}MISS\_\allowbreak{}L2} & 2 & 1 --- PMC \\
\texttt{CACHE\_\allowbreak{}MISS\_\allowbreak{}L3} & 3 & 1 --- PMC \\
\texttt{DTLB\_\allowbreak{}MISS} & 4 & 1 --- PMC \\
\texttt{MEMORY\_\allowbreak{}FOOTPRINT} & 5 & 0 --- time/\allowbreak{}observer \\
\texttt{BRANCH\_\allowbreak{}MISS} & 6 & 1 --- PMC \\
\texttt{IPC\_\allowbreak{}CPI} & 7 & 1 --- PMC \\
\texttt{LATENCY\_\allowbreak{}MEAN} & 8 & 0 --- time/\allowbreak{}observer \\
\texttt{LATENCY\_\allowbreak{}P50} & 9 & 0 --- time/\allowbreak{}observer \\
\texttt{LATENCY\_\allowbreak{}P95} & 10 & 0 --- time/\allowbreak{}observer \\
\texttt{LATENCY\_\allowbreak{}P99} & 11 & 0 --- time/\allowbreak{}observer \\
\texttt{LATENCY\_\allowbreak{}P999} & 12 & 0 --- time/\allowbreak{}observer \\
\texttt{THROUGHPUT} & 13 & 0 --- time/\allowbreak{}observer \\
\texttt{ENERGY\_\allowbreak{}J} & 14 & 1 --- PMC \\
\texttt{FILL\_\allowbreak{}BUFFER\_\allowbreak{}OCCUPANCY} & 15 & 0 --- time/\allowbreak{}observer \\
\bottomrule
\end{longtable}
\end{scriptsize}

\emph{Regime split:} nine categories lie in the time and observer regime (\texttt{CLU}, \texttt{MEMORY\_\allowbreak{}FOOTPRINT}, the five latency categories, \texttt{THROUGHPUT}, \texttt{FILL\_\allowbreak{}BUFFER\_\allowbreak{}OCCUPANCY}), seven in the PMC regime (the three cache-miss levels, \texttt{DTLB\_\allowbreak{}MISS}, \texttt{BRANCH\_\allowbreak{}MISS}, \texttt{IPC\_\allowbreak{}CPI}, \texttt{ENERGY\_\allowbreak{}J}). The split coincides with the \enquote{nine/\allowbreak{}seven} division of Section~\ref{sec:axis-triad}.

\emph{Contributing sources:} --- (CE baseline, own building blocks).

\section[head={Measurement Sphere: Measurement Tooling/Collector}]{Measurement Sphere: Measurement Tooling/\allowbreak{}Collector (\texttt{collector})}

This quantity is the instrument axis of the measurement realm: it records with which tool a category is acquired (cf.\ Section~\ref{sec:axis-triad}).

It carries \texttt{stage=ct}, \texttt{binary\_\allowbreak{}id=never} and \texttt{baustein\_\allowbreak{}count=3}. Every building block carries its regime and the set of categories it serves; the union of the three sets covers ten of the sixteen categories. Of the remaining six, the four latency-percentile categories (\texttt{LATENCY\_\allowbreak{}P50} through \texttt{LATENCY\_\allowbreak{}P999}) are served not by an acquisition tool but by the downstream evaluation (nearest-rank over the raw values, complemented by the HDR histogram, Section~\ref{sec:impl-system-axes}); \texttt{MEMORY\_\allowbreak{}FOOTPRINT} and \texttt{FILL\_\allowbreak{}BUFFER\_\allowbreak{}OCCUPANCY} are described in the offering catalogue but assigned to no acquisition tool.

\subsection*{Building blocks}

\begin{scriptsize}
\begin{longtable}{@{}>{\raggedright\arraybackslash}p{3.6cm} >{\raggedright\arraybackslash}p{1.6cm} >{\raggedright\arraybackslash}p{8.3cm}@{}}
\toprule
Variant & Regime & Categories served \\
\midrule
\endhead%
\texttt{Wall\allowbreak{}Clock\allowbreak{}System\allowbreak{}Axis} & 0 & \texttt{LATENCY\_\allowbreak{}MEAN}, \texttt{THROUGHPUT} \\
\texttt{Observer\allowbreak{}Snapshot\allowbreak{}System\allowbreak{}Axis} & 0 & \texttt{CLU} \\
\texttt{Pmc\allowbreak{}System\allowbreak{}Axis} & 1 & \texttt{CACHE\_\allowbreak{}MISS\_\allowbreak{}L1}, \texttt{CACHE\_\allowbreak{}MISS\_\allowbreak{}L2}, \texttt{CACHE\_\allowbreak{}MISS\_\allowbreak{}L3}, \texttt{DTLB\_\allowbreak{}MISS}, \texttt{BRANCH\_\allowbreak{}MISS}, \texttt{IPC\_\allowbreak{}CPI}, \texttt{ENERGY\_\allowbreak{}J} \\
\bottomrule
\end{longtable}
\end{scriptsize}

\emph{Contributing sources:} --- (CE baseline, own building blocks).

\section[head={Measurement Sphere: Telemetry Building Blocks}]{Measurement Sphere: Telemetry Building Blocks (Conserved Catalogue)}

This quantity is a \emph{two-way split} main quantity --- sweep sub-axis of the measurement tooling in the planner \emph{and} compile-time main quantity in the builder (cf.\ Section~\ref{sec:axis-triad}). Its building-block catalogue is retained here as a conserved catalogue.

Defines strategies for measuring index operations (insert, lookup, erase) with different overhead profiles. Supports density tracking for adaptive node-type selection, global insert counters for throughput measurements, HDR histograms for latency quantiles, and leaf-only counters to avoid cache-line false sharing in the measurement cycles.

\subsection*{Static sub-axes}

\begin{scriptsize}
\begin{longtable}{@{}>{\raggedright\arraybackslash}p{1.8cm} >{\raggedright\arraybackslash}p{3.2cm} >{\raggedright\arraybackslash}p{8.5cm}@{}}
\toprule
ID & Name & Description \\
\midrule
\endhead%
TM1 & Scope & Validity scope of the counter management: leaf-only (leaf nodes only), per-node (all nodes), or global (single accumulator for the entire structure). Compile-time choice, impacts overhead through contention on shared cache lines. \\
TM2 & Metric Type & Type of the measured metric: counter (event counter), histogram (distribution with quantiles), or density (slot occupancy per node). Impacts the memory footprint and the measurement granularity. \\
TM3 & Overhead Level & Magnitude of the measurement overhead: low (atomic operations), medium (per-node data structures), high (histogram logging). Determines the suitability for different workload scenarios (benchmarking vs.\ production). \\
\bottomrule
\end{longtable}
\end{scriptsize}

\subsection*{Building blocks}

\begin{scriptsize}
\begin{longtable}{@{}>{\raggedright\arraybackslash}p{2.5cm} >{\raggedright\arraybackslash}p{7.5cm} >{\raggedright\arraybackslash}p{3.5cm}@{}}
\toprule
Variant & Description & Source \\
\midrule
\endhead%
\texttt{Leaf\allowbreak{}Only\allowbreak{}Counter} & Counter \emph{only} in leaf nodes (avoids cache-line ping-pong by dispensing with updates of inner nodes). Low overhead with leaf-only scope. Corresponds to the coherence-friendly counter concept from Kuehn Da\allowbreak{}Mo\allowbreak{}N 2023, although the original paper is not a leaf-counter algorithm --- CE baseline for false-sharing avoidance. & P28 (Kuehn Da\allowbreak{}Mo\allowbreak{}N 2023, context; Code/\allowbreak{}Baseline) \\
\texttt{Density\allowbreak{}Tracker} & Density/\allowbreak{}slot-occupancy tracking per node for adaptive node-type selection (4→16→48→256). Contextual telemetry tied to the ART structure (Leis ICDE 2013), measures the slot occupancy of all nodes. Higher overhead than Leaf\allowbreak{}Only, but provides the adaptation signal for dynamic rebalancing decisions. & P01 (ART/\allowbreak{}Leis ICDE 2013, context; Code/\allowbreak{}Baseline) \\
\texttt{Insert\allowbreak{}Counter} & Global atomic insert counter with minimal overhead (1 atomic increment per insert). Contextual telemetry tied to the HOT structure (Binna SIGMOD 2018), global scope. Sufficient for throughput measurement series, no latency granularity. & P02 (HOT/\allowbreak{}Binna SIGMOD 2018, context; Code/\allowbreak{}Baseline) \\
\texttt{Latency\allowbreak{}Histogram} & HDR histogram for the lookup latency distribution (p50/\allowbreak{}p95/\allowbreak{}p99 quantiles). Based on Hdr\allowbreak{}Histogram (Gil Tene), CC0-1.0 licence. Highest overhead (per lookup), but provides complete latency statistics for SLA-oriented measurement series and measurement series with scientific ambition. Only is\_\allowbreak{}original candidate of the axis. & Hdr\allowbreak{}Histogram (github.com/\allowbreak{}Hdr\allowbreak{}Histogram/\allowbreak{}Hdr\allowbreak{}Histogram\_\allowbreak{}c, CC0-1.0, is\_\allowbreak{}original) \\
Probability-Hints & Probability hints of the cache-oblivious Bender line (P16/\allowbreak{}P17): distribution- and probability-related occupancy hints that serve as input to the layout decisions. Paper building block of the axis dissection (reconstruction target of this thesis, no CE baseline). & P16/\allowbreak{}P17 (Bender 2000/\allowbreak{}2005)~\cite{bender2000cobtree,bender2005coboblivious} \\
\bottomrule
\end{longtable}
\end{scriptsize}

The CRTP organ registry of the telemetry (\texttt{All\allowbreak{}Telemetries}) carries exactly the four strategies of the first four rows; beside it the code contains the vtable family \texttt{ITelemetry\allowbreak{}Strategy} with six kinds (per-node counter, leaf-only counter, sampled leaf-only counter, retroactive aggregation, path-read counter P26 and the probability hints of this thesis), which is no organ registry and is consumed by the telemetry engine (\texttt{c05\_\allowbreak{}telemetry\_\allowbreak{}engine}) and its tests. Two further telemetry targets with ZIH relation --- an OTF2 trace connection to the performance analysis tool Vampir~\cite{knupfer2008vampir} and a connection to the centre-wide job monitor PIKA~\cite{dietrich2020pika} --- are foreseen as marker types of a variant catalogue consumed only by tests and isolated from the build, and are not part of the build-driving offering; the reconstruction of the VAMPIR memory-type design (A24) is independent of that.

\emph{Contributing sources:} Hdr\allowbreak{}Histogram (Gil Tene)~\cite{tene_hdrhistogram}, P16/\allowbreak{}P17 (Bender)~\cite{bender2000cobtree,bender2005coboblivious}.

\emph{Related literature:} P01, P02, P28~\cite{leis2013art,binna2018hot,kuehn2023bplustree}.

\section[head={Measurement Sphere: Load Framework --- Meta-Meta Axis}]{Measurement Sphere: Load Framework (\texttt{load\_framework}) --- Meta-Meta Axis}

This quantity is the meta-meta main axis of the measurement realm and therefore stands at its end (cf.\ Section~\ref{sec:axis-triad}).

It carries \texttt{category=measurement\_\allowbreak{}meta\_\allowbreak{}meta}, \texttt{stage=ct}, \texttt{binary\_\allowbreak{}id=never} and \texttt{baustein\_\allowbreak{}count=1}. It is the meta-meta main axis of the planner in the measurement realm; it explicitly does not belong under the hub of the system meta-metas. In the stamp it receives \emph{no} meta-meta line of its own: its identifier appears as a dynamic bracket appendage at the \emph{end} of the existing measurement realm line. The feeding path here is entry-driven --- the chosen identifier together with its version comes from the framework registry of the measurement realm and therefore does not reside in the type, unlike the type-driven system hub.

\subsection*{Dynamic sub-axes}

The registry carries seven sweep dimensions (\texttt{resource\_\allowbreak{}control\_\allowbreak{}version=1}); six of them are \texttt{stage=runtime}, the seventh (\texttt{work\_\allowbreak{}mode}) is \texttt{stage=plan}.

\begin{scriptsize}
\begin{longtable}{@{}>{\raggedright\arraybackslash}p{3.6cm} >{\raggedright\arraybackslash}p{4.4cm} >{\raggedright\arraybackslash}p{5.5cm}@{}}
\toprule
Dimension & Origin (\texttt{source}) & Value type and meaning \\
\midrule
\endhead%
\texttt{workload} & \texttt{measurement:\allowbreak{}load\_\allowbreak{}framework} & \texttt{token} --- the load profile; the single source is the sub-axis label of the load framework itself. \\
\texttt{thread\_\allowbreak{}count} & \texttt{resource\_\allowbreak{}control\_\allowbreak{}pod} & \texttt{uint} --- number of worker threads. \\
\texttt{prefetch\_\allowbreak{}distance} & \texttt{resource\_\allowbreak{}control\_\allowbreak{}pod} & \texttt{uint} --- prefetch distance. \\
\texttt{pool\_\allowbreak{}budget\_\allowbreak{}bytes} & \texttt{resource\_\allowbreak{}control\_\allowbreak{}pod} & \texttt{uint} --- memory budget of the pool. \\
\texttt{batch\_\allowbreak{}size} & \texttt{resource\_\allowbreak{}control\_\allowbreak{}pod} & \texttt{uint} --- batch size per operation. \\
\texttt{inline\_\allowbreak{}threshold\_\allowbreak{}bytes} & \texttt{resource\_\allowbreak{}control\_\allowbreak{}pod} & \texttt{uint} --- threshold of inline storage. \\
\texttt{work\_\allowbreak{}mode} & \texttt{measurement:\allowbreak{}run\_\allowbreak{}methodology} & \texttt{token} --- the run method of a run (\texttt{build}, \texttt{measure}, \texttt{compare}, \texttt{release}), \texttt{stage=plan}; the single source is the compile-time registry of the run methodology, no hand-written value. \\
\bottomrule
\end{longtable}
\end{scriptsize}

\emph{Note on non-observation:} the repetition count of a run is deliberately \emph{not} in this list: it is a pure label dimension without a compile-time counterpart and becomes reflectable only after its typing; a manual entry would be an unsubstantiated value --- the registry instead carries an explicitly coded non-observation. The run methods, by contrast, are typed: there are \emph{four} --- build, measurement, comparison and release execution (\texttt{build}, \texttt{measure}, \texttt{compare}, \texttt{release}) ---, backed by a compile-time registry without run-time branching; the comparison execution possesses the selectability and reads the measurement store without acquiring reference values itself. They stand in the list as the dimension \texttt{work\_\allowbreak{}mode}.

\subsection*{Building blocks}

\begin{scriptsize}
\begin{longtable}{@{}>{\raggedright\arraybackslash}p{2.4cm} >{\raggedright\arraybackslash}p{7.6cm} >{\raggedright\arraybackslash}p{3.5cm}@{}}
\toprule
Variant & Description & Source \\
\midrule
\endhead%
\texttt{ycsb} & The YCSB load-profile framework; enabled, with the sub-axis label \texttt{workload}. The only enabled building block of the axis. & Code/\allowbreak{}Baseline \\
\bottomrule
\end{longtable}
\end{scriptsize}

\emph{Contributing sources:} --- (CE baseline, own building blocks).

\section{A-Korpus Allocator Corpus (A01--A23 and A24)}

The allocator axis documents 26 memory management strategies (identifiers A01--A21, the three standard-library baselines A22a--A22c, A23 and A24) that affect dynamic memory management in search workloads. It covers classic textbook allocators (Buddy, Slab), industrial standards (dlmalloc, jemalloc, TCMalloc, mimalloc, snmalloc), specialized variants (NUMA-aware, PIM-Malloc, cache-aware), modern hardened implementations (Star\allowbreak{}Malloc) and a memory-type virtualization following the VAMPIR design (A24). Here A01--A23 form the literature corpus of the 23 works of Section~\ref{sec:sota-instances} with one allocator profile each (Table~\ref{tab:allocator-profiles}), whereas A24 is the allocator front reconstructed from the SOTA profile P33 without an allocator profile of its own. All 26 wrappers satisfy the \texttt{Allocator\allowbreak{}Strategy} concept and enable systematic measurement of the influence of memory management on search-tree performance (F15 operativity). The enabled inventory of this build configuration comprises three wrappers (see the registry-state note in section T6).

\subsection*{Static sub-axes}

The sub-axes AA1--AA7 are defined in section T6; the following table assigns the corpus examples to them.

\begin{scriptsize}
\begin{longtable}{@{}>{\raggedright\arraybackslash}p{1.8cm} >{\raggedright\arraybackslash}p{3.2cm} >{\raggedright\arraybackslash}p{8.5cm}@{}}
\toprule
ID & Name & Description \\
\midrule
\endhead%
AA1 & Free-List Topology & Structure of the free memory blocks: per-page sharding (mimalloc), per-thread caching (snmalloc, rpmalloc, TCMalloc), buddy tree (Buddy allocator), single global free-list (dlmalloc), per-heap locks (Hoard). \\
AA2 & Size-Class Schema & Classification of block sizes: slab-per-object-size (Slab), dlmalloc bins, jemalloc size-classes, scalloc spans, Buddy power-of-2. \\
AA3 & Thread Locality & Degree of thread isolation: tcmalloc thread-local caches, rpmalloc per-thread spanning, snmalloc message-passing, Hoard per-heap locks. \\
AA4 & Synchronization & Concurrency mechanism: lock-free CAS (Michael, LRMalloc), per-heap locks (Hoard), wait-free (Crystalline), mutex-based. \\
AA5 & Allocation Policy & Routing of allocation requests: NUMA-origin-aware (NUMAlloc), cache-set-aware (CAMA), NUCA-aware (TCMalloc-Warehouse), PIM-aware (PIM-Malloc), NFP memory-type virtualization (Vampir\allowbreak{}Nfp). \\
AA6 & Reclamation & Release of deleted blocks: magazine cache (Slab), page-heap release (TCMalloc), lock-free reclamation (LRMalloc), deferred free (mimalloc). \\
AA7 & Fragmentation Strategy & Avoidance of fragmentation: dlmalloc coalescing, object colouring (Slab), buddy splitting, page-local sharding (mimalloc), span pooling (scalloc). \\
\bottomrule
\end{longtable}
\end{scriptsize}

\subsection*{Building blocks}

\begin{scriptsize}
\begin{longtable}{@{}>{\raggedright\arraybackslash}p{3.0cm} >{\raggedright\arraybackslash}p{6.0cm} >{\raggedright\arraybackslash}p{4.5cm}@{}}
\toprule
Variant & Description & Source \\
\midrule
\endhead%
\texttt{Hoard\allowbreak{}Allocator} (A01) & Per-heap free-lists with SMP scaling. Three-tier structure: thread-local heaps, global heap for cross-thread migration, false-sharing reduction via segregated allocation requests. & A01~\cite{berger2000hoard} (Berger/\allowbreak{}Mc\allowbreak{}Kinley/\allowbreak{}Blumofe/\allowbreak{}Wilson, ASPLOS 2000, 10.1145/\allowbreak{}378993.379232) + Code/\allowbreak{}Baseline \\
\texttt{Slab\allowbreak{}Allocator} (A02) & Object-caching kernel allocator. Slab-per-object-size with per-CPU magazine cache (Bonwick 1994/\allowbreak{}2001). Classic Solaris kernel implementation for pre-allocated object pools. & A02~\cite{bonwick1994slab} (Bonwick, USENIX 1994) + Code/\allowbreak{}Baseline \\
\texttt{Michael\allowbreak{}Lock\allowbreak{}Free\allowbreak{}Allocator} (A03) & Lock-free CAS allocator per Michael 2004. Hazard-pointer-protected lock-free data structures, LGPL re-implementation with IBM patent background. & A03~\cite{michael2004lockfree} (Michael, PLDI 2004, 10.1145/\allowbreak{}996841.996848) + Code/\allowbreak{}Baseline \\
\texttt{Mimalloc\allowbreak{}Allocator} (A04) & Free-list sharding (3 free-lists per page). Microsoft Research 2019: deferred\_\allowbreak{}free hooks, temporal cadence, page-local sharding.\ is\_\allowbreak{}original: MIT direct binding. & A04~\cite{leijen2019mimalloc} (Leijen/\allowbreak{}Zorn/\allowbreak{}de Moura, APLAS 2019, 10.1007/\allowbreak{}978-3-030-34175-6\_\allowbreak{}13) + is\_\allowbreak{}original (github.com/\allowbreak{}microsoft/\allowbreak{}mimalloc, MIT) \\
\texttt{Jemalloc\allowbreak{}Allocator} (A05) & Arena-based size-class allocator. Evans 2006 /\allowbreak{} Free\allowbreak{}BSD\@. Per-CPU arenas with tcache layer for thread-local caches. Facebook standard\@. is\_\allowbreak{}original: BSD-2-Clause direct binding. & A05~\cite{evans2006jemalloc} (Evans, BSDCan 2006, papers.freebsd.org/\allowbreak{}\ldots/\allowbreak{}jemalloc) + is\_\allowbreak{}original (github.com/\allowbreak{}jemalloc/\allowbreak{}jemalloc, BSD-2-Clause) \\
\texttt{TCMalloc\allowbreak{}Allocator} (A06) & Thread-cache + central-cache + page-heap. Google gperftools 2005. Later TEMERAIRE (OSDI 2021) with warehouse-scale optimizations. & A06~\cite{hunter2021temeraire} (TCMalloc 2005; TEMERAIRE, OSDI 2021) + Code/\allowbreak{}Baseline \\
\texttt{Snmalloc\allowbreak{}Allocator} (A07) & Message-passing allocator. Paul Liétar et al. ISMM 2019: free operations via lock-free queue to the allocating thread, cross-thread lock-contention avoidance\@. is\_\allowbreak{}original: MIT direct binding. & A07~\cite{lietar2019snmalloc} (Liétar, ISMM 2019, 10.1145/\allowbreak{}3315573.3329980) + is\_\allowbreak{}original (github.com/\allowbreak{}microsoft/\allowbreak{}snmalloc, MIT) \\
\texttt{Scalloc\allowbreak{}Allocator} (A08) & Span-based, virtual-spans technique. Aigner/\allowbreak{}Kirsch/\allowbreak{}Lippautz/\allowbreak{}Sokolova OOPSLA 2015. Multicore-scalable with low fragmentation. & A08~\cite{aigner2015scalloc} (Aigner et al., OOPSLA 2015) + Code/\allowbreak{}Baseline \\
\texttt{NUMAlloc\allowbreak{}Allocator} (A09) & NUMA-origin-aware allocator. UTSASRG ISMM 2023. Node hint for local NUMA allocation, multi-node systems optimized. & A09~\cite{yang2023numalloc} (Yang et al., ISMM 2023) + Code/\allowbreak{}Baseline \\
\texttt{RPMalloc\allowbreak{}Allocator} (A10) & Per-thread span caching. Mattias Jansson 2017, open-source implementation without a publication. Special case: mandatory init (rpmalloc\_\allowbreak{}initialize, rpmalloc\_\allowbreak{}thread\_\allowbreak{}initialize)\@. is\_\allowbreak{}original: Public-Domain/\allowbreak{}MIT direct binding. & Code/\allowbreak{}Baseline (github.com/\allowbreak{}mjansson/\allowbreak{}rpmalloc, Public-Domain/\allowbreak{}MIT) + is\_\allowbreak{}original \\
\texttt{LRMalloc\allowbreak{}Allocator} (A11) & Lock-free + hazard pointers. Leite/\allowbreak{}Rocha VECPAR 2018 (LNCS 11333). Modern lock-free allocator with ABA protection.\ is\_\allowbreak{}original: MIT direct binding. & A11~\cite{leite2018lrmalloc} (Leite/\allowbreak{}Rocha, VECPAR 2018, 10.1007/\allowbreak{}978-3-030-15996-2\_\allowbreak{}17) + is\_\allowbreak{}original (github.com/\allowbreak{}ricleite/\allowbreak{}lrmalloc, MIT) \\
\texttt{CAMAAllocator} (A12) & Cache-aware memory allocator. Herter/\allowbreak{}Backes/\allowbreak{}Haupenthal/\allowbreak{}Reineke ECRTS 2011 (Saarland). Predictable cache behaviour, no public code reference. & A12~\cite{herter2011cama} (Herter et al., ECRTS 2011) + Code/\allowbreak{}Baseline \\
\texttt{Star\allowbreak{}Malloc\allowbreak{}Allocator} (A13) & Formally verified hardened allocator. Inria-Prosecco (Reitz/\allowbreak{}Fromherz/\allowbreak{}Protzenko et al.) OOPSLA 2024. F*/\allowbreak{}Steel formal verification, Apache-2.0 licence. & A13~\cite{reitz2024starmalloc} (Reitz et al., OOPSLA 2024) + Code/\allowbreak{}Baseline \\
\texttt{TCMalloc\allowbreak{}Warehouse\allowbreak{}Allocator} (A14) & TCMalloc-Warehouse /\allowbreak{} Hyperscale. Google Cloud OSDI 2021 /\allowbreak{} ASPLOS 2024. Hugepage-aware variant, product variant without a separate artefact. & A14~\cite{hunter2021temeraire,zhou2024tcmallocwarehouse} (Hunter OSDI 2021 / Zhou ASPLOS 2024) + Code/\allowbreak{}Baseline \\
\texttt{HMalloc\allowbreak{}Allocator} (A15) & Hybrid, scalable, lock-free memory allocator. Li/\allowbreak{}Yao/\allowbreak{}Tang/\allowbreak{}Lin ICPADS 2019. Hybrid approach for scalability, no public code reference. & A15~\cite{li2019hmalloc} (Li/\allowbreak{}Yao/\allowbreak{}Tang/\allowbreak{}Lin, ICPADS 2019, 10.1109/\allowbreak{}ICPADS47876.2019.00064) + Code/\allowbreak{}Baseline \\
\texttt{PIMMalloc\allowbreak{}Allocator} (A16) & Processing-in-memory allocator. VIA-Research HPCA 2026 (ar\allowbreak{}Xiv:2505.13002). Memory distribution across compute nodes, PIM-hardware-optimized. & A16~\cite{pimmalloc2026} (VIA-Research, HPCA 2026, arxiv.org/\allowbreak{}abs/\allowbreak{}2505.13002) + Code/\allowbreak{}Baseline \\
\texttt{Crystalline\allowbreak{}Reclamation} (A17) & Wait-free or lock-/\allowbreak{}wait-free safe-memory-reclamation family for non-blocking data structures. Not a standalone malloc/\allowbreak{}free allocator, but a reclamation policy. & A17~\cite{nikolaev2021crystalline,nikolaev2024crystalline} (Nikolaev/\allowbreak{}Ravindran, DISC 2021 / PLDI 2024) + Code/\allowbreak{}Baseline \\
\texttt{Exgen\allowbreak{}Allocator} (A18) & Single-threaded specialized allocator. UT Austin IEEE CAL 2025 (ar\allowbreak{}Xiv:2510.10219)\@. \enquote{Old is Gold} optimizations for single-thread workloads, no public code reference. & A18~\cite{li2025exgen} (UT Austin, IEEE CAL 2025, arxiv.org/\allowbreak{}abs/\allowbreak{}2510.10219) + Code/\allowbreak{}Baseline \\
\texttt{Buddy\allowbreak{}Allocator} (A19) & Buddy system (power-of-2 splitting). Knuth TAo\allowbreak{}CP Vol.1 (1968), CACM 1965. Classic textbook allocator, binary-tree splitting without OSS code reference. & A19~\cite{knowlton1965buddy} (Knowlton, CACM 1965, 10.1145/\allowbreak{}365628.365655) + Code/\allowbreak{}Baseline \\
\texttt{Dlmalloc\allowbreak{}Allocator} (A20) & Doug Lea malloc (bins-based). Lea 1987--2012, public domain (CC0). Classic since 1987, bins size-class schema, oldest production implementation\@. is\_\allowbreak{}original: public-domain direct binding. & A20~\cite{lea1996dlmalloc} (Lea, gee.cs.oswego.edu/\allowbreak{}\ldots/\allowbreak{}malloc.html, 1996) + is\_\allowbreak{}original (github.com/\allowbreak{}ennorehling/\allowbreak{}dlmalloc, Public-Domain/\allowbreak{}CC0) \\
\texttt{Pt\allowbreak{}Malloc2\allowbreak{}Allocator} (A21) & ptmalloc2 (glibc malloc). Wolfram Gloger 1990s-2000s. LGPL-2.1+ copyleft, linking blocked (no is\_\allowbreak{}original). Standard basis for glibc-based systems. & Code/\allowbreak{}Baseline (malloc.de/\allowbreak{}en/\allowbreak{}) + License restriction \\
\texttt{Std\allowbreak{}Malloc} (A22a) & Standard libc malloc (proof of concept), portable fallback allocator. Baseline for all standard-library implementations without vendor linking. & Code/\allowbreak{}Baseline (C++ Standard Library) + Baseline \\
\texttt{Pmr\allowbreak{}Resource\allowbreak{}Allocator} (A22b) & std::pmr::memory\_\allowbreak{}resource. WG21 N3916 (2014 $\rightarrow$ C++17). Polymorphic Memory Resources standard library, fallback to the new/\allowbreak{}delete allocator. & Code/\allowbreak{}Baseline (C++17 Standard Library, open-std.org/\allowbreak{}\ldots/\allowbreak{}n3916.pdf) + Baseline \\
\texttt{Pool\allowbreak{}Resource\allowbreak{}Allocator} (A22c) & std::pmr::unsynchronized\_\allowbreak{}pool\_\allowbreak{}resource. WG21 N3916 (C++17). Own size-class pool, F15-operative (first wrapper with its own memory-retention behaviour without vendor linking). & Code/\allowbreak{}Baseline (C++17 Standard Library, open-std.org/\allowbreak{}\ldots/\allowbreak{}n3916.pdf) + is\_\allowbreak{}operative \\
\texttt{Vmem\allowbreak{}Magazines\allowbreak{}Allocator} (A23) & Vmem + magazines (slab extension). Bonwick 2001 USENIX ATC\@. Per-CPU magazines for scaling, kernel classic without public code reference. & A23~\cite{bonwick2001vmem} (Bonwick, USENIX ATC 2001, usenix.org/\allowbreak{}\ldots/\allowbreak{}bonwick.html) + Code/\allowbreak{}Baseline \\
\texttt{Vampir\allowbreak{}Nfp\allowbreak{}Allocator} (A24) & Memory-type virtualization following the VAMPIR design (Virtualized Non-Functional Memory Properties for Data-Pipeline Scheduling): allocator front over an NFP-labelled memory tier (target DRAM) with an exposed NFP descriptor (\texttt{bandwidth\_\allowbreak{}class}, \texttt{latency\_\allowbreak{}class}, \texttt{capacity}); structural reconstruction of the buildable part, no cost formula, no scheduling, no energy metric; disabled by default. & A24~\cite{vampir2023sosp} (VAMPIR poster, SOSP 2023; TU Dresden) + Code/\allowbreak{}Baseline (is\_\allowbreak{}original=false) \\
\bottomrule
\end{longtable}
\end{scriptsize}

\emph{Contributing sources:} cited per allocator A01--A24 individually in the source column of the building-block table.

\setlength{\tabcolsep}{6pt}

\section{Axis Inventory from the Generated Registries}

The following table is not maintained by hand: the appendix generator
(\texttt{Code/\allowbreak{}08\_\allowbreak{}appendix\_\allowbreak{}generator}) derives it from the three generated axis registries of the
code --- organ, system and measurement realm --- and writes it, provided the run is invoked with at least one of the three
registry options (\texttt{-{}-organ-registry}, \texttt{-{}-system-registry}, \texttt{-{}-measurement-registry}),
per language to \texttt{anhang/<lang>/tabellen/}. It is therefore a derivation from the
compile-time reflection and not a second, hand-kept state beside the registries. If the
registries contain no axis, no file is produced on purpose; an empty table is not
synthesized.\IfFileExists{anhang/en/tabellen/axis_inventory.tex}{ Table~\ref{tab:axis-inventory} shows the state of the
last generator run.}{}

\InputIfFileExists{anhang/en/tabellen/axis_inventory.tex}{}{%
  \emph{(The axis inventory is not available for this state; a generator run with registry options produces it.)}}

\section{Conserved Designs}\label{sec:app-blocks-konserviert}
The three following offering catalogues are conserved designs: living code with registries of its own that
stands outside the build-driving offering of the three axis registries --- the SIMD catalogue Build-SIMD under the
extension hardware, the page-type catalogue Build-PG outside the eighteen composition slots and the
platform hardware strategy Build-HW, which the system registry explicitly excludes. They remain listed here
in full --- with purpose, sub-axes, building blocks and sources --- and are to be distinguished from the offering of the
registries, which the preceding sections document (cf.\ Section~\ref{sec:axis-triad}).

\subsection[Build-SIMD: Vector/Accelerator Extension Axis]{Build-SIMD: Vector/\allowbreak{}Accelerator Extension Axis (Hardware Accelerators)} % chktex 13 (Titel-Doppelpunkt nach Grossbuchstabe, kein Satzende)

\emph{Conserved offering catalogue outside the build-driving offering:} the axis \texttt{axis\_\allowbreak{}09b\_\allowbreak{}simd\_\allowbreak{}extension} is an offering catalogue under the main axis \texttt{external\_\allowbreak{}utils}, whose sub-axis \texttt{simd} carries the build-driving choice; its eight building blocks and the sub-axes SE1--SE3 remain listed in full.

The axis \texttt{axis\_\allowbreak{}09b\_\allowbreak{}simd\_\allowbreak{}extension} models the available SIMD/\allowbreak{}accelerator extensions of the CPU instruction set architecture as a sub-axis to \texttt{axis\_\allowbreak{}09} (main ISA). It enables compile-time selection among various vector extensions (SSE2/\allowbreak{}AVX2/\allowbreak{}AVX-512 for x86; NEON/\allowbreak{}SVE2 for ARM\@; RVV for RISC-V) plus baseline (scalar) and GPU acceleration (CUDA), with backward-cumulative layer hierarchy and hardware-specific topology (SIMD units per socket, P/\allowbreak{}E core access).

\subsubsection*{Static sub-axes}

\begin{scriptsize}
\begin{longtable}{@{}>{\raggedright\arraybackslash}p{1.8cm} >{\raggedright\arraybackslash}p{3.2cm} >{\raggedright\arraybackslash}p{8.5cm}@{}}
\toprule
ID & Name & Description \\
\midrule
\endhead%
SE1 & Vector Width & Compile-time-known maximum vector width of the SIMD extension in bits: 0 (none), 128, 256, 512, or -1 (scalable). Determines lane width and parallelization degree. \\
SE2 & Accelerator Type & Categorizes the accelerator class: none (baseline), simd (vector SIMD), matrix (matrix engine), gpu (GPU), npu (neural processing unit). Influences the composition with the hardware topology (\texttt{axis\_\allowbreak{}12}). \\
SE3 & Compatibility Family & Constraint to the main ISA (\texttt{axis\_\allowbreak{}09}): none (universal), x86, arm, riscv, cuda. Ensures compile-time validation: e.g., \texttt{Sse2\allowbreak{}Simd\allowbreak{}Extension} compatible only with \texttt{Amd64\allowbreak{}Isa}. \\
\bottomrule
\end{longtable}
\end{scriptsize}

\subsubsection*{Building blocks}

\begin{scriptsize}
\begin{longtable}{@{}>{\raggedright\arraybackslash}p{3.0cm} >{\raggedright\arraybackslash}p{7.5cm} >{\raggedright\arraybackslash}p{3.0cm}@{}}
\toprule
Variant & Description & Source \\
\midrule
\endhead%
\texttt{No\allowbreak{}Simd\allowbreak{}Extension} & Baseline variant without SIMD acceleration. Scalar fallback strategy for all permutations without accelerator. Universally compatible with all main ISAs (x86/\allowbreak{}ARM/\allowbreak{}RISC-V/\allowbreak{}Power\allowbreak{}PC). Mandatory comparison point. & Code/\allowbreak{}Baseline \\
\texttt{Sse2\allowbreak{}Simd\allowbreak{}Extension} & Intel SSE2 128-bit vector extension (x86\_\allowbreak{}64 ABI mandatory since 2003). Compatible only with \texttt{Amd64\allowbreak{}Isa}. Baseline on x86\_\allowbreak{}64; provides 4× uint32 lanes/\allowbreak{}iteration. Standard on the cluster nodes Barnard/\allowbreak{}Capella. & Intel ISA extensions reference 319433~\cite{intel_isa_ext_ref} \\
\texttt{Avx2\allowbreak{}Simd\allowbreak{}Extension} & Intel AVX2 256-bit vector extension (Haswell+ 2013). Only x86\_\allowbreak{}64. Backward-cumulative: includes all SSE+AVX\@. Default optimization level for server DBMS\@. 2 AVX2 units/\allowbreak{}socket. ZIH standard on Barnard (Intel Sapphire Rapids) and Capella (AMD EPYC 9334). & Intel ISA extensions reference 319433~\cite{intel_isa_ext_ref} \\
\texttt{Avx512\allowbreak{}Simd\allowbreak{}Extension} & Intel AVX-512 512-bit vector extension (Skylake-X+, Ice Lake Server, Zen 4+). Backward-cumulative: all SSE/\allowbreak{}AVX/\allowbreak{}AVX2 + 15 separate flags (VNNI for vector indexes, VPOPCNTDQ for bitmap indexes). Not on all x86\_\allowbreak{}64 CPUs (for instance Haswell without AVX-512; on Intel's hybrid Alder Lake, Raptor Lake and Meteor Lake processors AVX-512 is switched off in the BIOS because of the E cores --- so also on prod2, whose signature contains no avx512 feature). 1 AVX-512 unit/\allowbreak{}socket; E cores no access (Intel Hybrid). & Intel ISA extensions reference 319433~\cite{intel_isa_ext_ref} \\
\texttt{Neon\allowbreak{}Simd\allowbreak{}Extension} & ARM NEON /\allowbreak{} Advanced SIMD 128-bit (AArch64 ABI mandatory ARMv8+). Compatible only with \texttt{Aarch64\allowbreak{}Isa}. Baseline on all AArch64 CPUs (Apple M-Series, AWS Graviton 2/\allowbreak{}3/\allowbreak{}4, NVIDIA Grace). 1 NEON unit/\allowbreak{}socket; all cores accessible (big.LITTLE architecture). & ARM ARM DDI 0487 /\allowbreak{} Cortex-A8 TRM \\
\texttt{Sve2\allowbreak{}Simd\allowbreak{}Extension} & ARM SVE2 (Scalable Vector Extension 2) ARMv9 mandatory (scalable 128--2048 bit). Compatible only with \texttt{Aarch64\allowbreak{}Isa}. ZIH test system NVIDIA Grace Hopper Superchip (Grace: 72 Arm Neoverse V2 cores with SVE2; test system with limited availability~\cite{zih_hpc_compendium_gh200}). Scalable lane width; 1 unit/\allowbreak{}socket; all cores. & SVE~\cite{stephens2017sve} (IEEE Micro 2017); SVE2/\allowbreak{}Armv9 per ARM ARM~\cite{arm_arm} \\
\texttt{Rvv\allowbreak{}Simd\allowbreak{}Extension} & RISC-V Vector Extension (RVV v1.0, scalable VLEN 128--65536 bit). Compatible only with \texttt{Risc\allowbreak{}VIsa}. Optional (not baseline like SSE2/\allowbreak{}NEON). Growing relevance: Si\allowbreak{}Five Performance P870 (128-bit), T-Head Xuantie C908 (128-bit). ZIH presence: not listed in the ZIH compendium. & RVV-v1.0-Spec~\cite{riscv_vector_spec_1_0}; Vitruvius+ (TACO 2023)~\cite{minervini2023vitruvius} \\
\texttt{Cuda\allowbreak{}Gh200\allowbreak{}Simd\allowbreak{}Extension} & NVIDIA CUDA on Hopper GH200 (Grace Hopper Superchip GPU accelerator). Hopper GPUs at the ZIH are located in Capella (4 $\times$ H100-SXM5 per node on AMD EPYC 9334 hosts~\cite{zih_hpc_compendium_hardware}) and in the Grace Hopper Superchip test system (Hopper GPU with 96 GB HBM3, NVLink-C2C 900 GB/\allowbreak{}s~\cite{zih_hpc_compendium_gh200}). Compatible with \texttt{Aarch64\allowbreak{}Isa} (Grace CPU) \emph{and} \texttt{Amd64\allowbreak{}Isa} (Hopper host). Massively parallel: 16\,896 FP32 CUDA cores per H100-SXM5 GPU~\cite{andersch2022hopper}. \texttt{units\_\allowbreak{}per\_\allowbreak{}socket()} = -1 (GPU is a separate device). & NVIDIA Tesla~\cite{lindholm2008tesla} (IEEE Micro 2008); NVIDIA Hopper Architecture In-Depth~\cite{andersch2022hopper} \\
\bottomrule
\end{longtable}
\end{scriptsize}

\emph{Sources:} ISA specifications and whitepapers (Intel, ARM, IEEE Micro, ACM TACO, NVIDIA); individual references in the source column of the building blocks.

\subsection{Build-PG Page Type Axis (PAGE-TYPE)}

\emph{Delimitation and conservation:} this section describes the \texttt{Page\allowbreak{}Kind} of the organ world (\texttt{axis\_\allowbreak{}01\_\allowbreak{}page\_\allowbreak{}type}) and is distinct from the system sub-axis \texttt{page} of the target ISA (page-size release 4k/\allowbreak{}2m/\allowbreak{}1g) --- both quantities carry similar names but share neither building blocks nor semantics. The catalogue is a conserved offering catalogue outside the eighteen composition slots and remains listed in full.

The page type axis defines the structural and indexing strategies for branch pages in the trie index: from dense direct-addressed byte indexing (Dense\allowbreak{}Byte), through multi-level density variants (Extended\allowbreak{}Dense, Sparse\allowbreak{}Patricia), to path-collapse optimizations (Redirect, Custom\allowbreak{}Cache) and B+-style sorted pages (BPlus). All page types are cache-line aware and implement the \texttt{Page\allowbreak{}Type\allowbreak{}Strategy} concept interface with guaranteed interchangeability (F15 permutation).

\subsubsection*{Static sub-axes}

\begin{scriptsize}
\begin{longtable}{@{}>{\raggedright\arraybackslash}p{1.8cm} >{\raggedright\arraybackslash}p{3.2cm} >{\raggedright\arraybackslash}p{8.5cm}@{}}
\toprule
ID & Name & Description \\
\midrule
\endhead%
PG1 & Structure Role (Branching) & Classification of the branch page architecture: branch page with full branching, collapsed single-path node (Redirect), or special cache-optimized variant. \\
PG2 & Density Class (Density) & Degree of slot utilization and the resulting indexing strategy: densely populated (0-256 slots), extended (>256 slots), or sparse with Patricia compression. \\
PG3 & Path Collapse Strategy & Handling of unique remainder paths: no compression (standard), or collapse via Redirect/\allowbreak{}Co\allowbreak{}Co technique for space minimization. \\
\bottomrule
\end{longtable}
\end{scriptsize}

\subsubsection*{Building blocks}

\begin{scriptsize}
\begin{longtable}{@{}>{\raggedright\arraybackslash}p{3.0cm} >{\raggedright\arraybackslash}p{8.0cm} >{\raggedright\arraybackslash}p{2.5cm}@{}}
\toprule
Variant & Description & Source \\
\midrule
\endhead%
\texttt{Dense\allowbreak{}Byte\allowbreak{}Page\allowbreak{}Type} & Dense byte-indexed branch page with direct 256-slot addressing (direct-addressed). Page type 1: direct mapping from key byte to slot index without redirection. Source: Adaptive Radix Tree (ART), Leis et al. 2013. & P01 \\
\texttt{Extended\allowbreak{}Dense\allowbreak{}Page\allowbreak{}Type} & Extended dense page for >256 entries with multi-level indexing scheme (e.g., 2-level lookup or variable k-byte span). Page type 2: density class for more complex slot arrangements. Source: Adaptive Radix Tree extensions, PRT-ART\@. & P01 (extended) \\
\texttt{Sparse\allowbreak{}Patricia\allowbreak{}Page\allowbreak{}Type} & Sparsely populated page with Patricia path compression (single-bit split). Page type 3: density class with bit-vector traversal for low branchingness. Source: HOT (Height Optimized Trie, Binna et al. 2018, based on the Patricia trie, Morrison 1968). & P02 \\
\texttt{Redirect\allowbreak{}Page\allowbreak{}Type} & Collapsed unique remainder path (single-path node). Page type 4: path-collapse strategy for memory optimization under monopoly-unique successor sequences. Neither classical branch nor leaf, stores only the suffix up to the next actual branching or value node. Source: Co\allowbreak{}Co-Trie (Boffa et al. 2024). & P04 \\
\texttt{Custom\allowbreak{}Cache\allowbreak{}Page\allowbreak{}Type} & Cache-line-oriented custom page with explicit layout design for 64-byte alignment. Page type 5: structure role branch with hardware-aware geometry. Source: PRT-ART own development with TLB offset + cache-line-aligned memory layout (axis 5). & PRT-ART \\
\texttt{BPlus\allowbreak{}Page\allowbreak{}Type} & B+-style sorted branch page with range-walk traversal. Page type 6: structure role branch, but with sorted key-array organization instead of byte indexing. Source: Masstree (Mao et al. 2012), B+-tree family. & P03 \\
\bottomrule
\end{longtable}
\end{scriptsize}

\emph{Contributing sources:} P01, P02, P03, P04, Morrison 1968, PRT-ART~\cite{leis2013art,binna2018hot,mao2012masstree,boffa2024coco,morrison1968patricia}.

\emph{Related literature:} P05, P06, P10, P11, P12, P13, P20~\cite{fent2020start,schmeisser2022b2tree,zhang2018surf,rao1999css,rao2000csb,hankins2003nodesize,mueller2025btreesback}.

\subsection{Build-HW — Platform Hardware Strategy}

\emph{Status (conserved offering catalogue):} this axis is excluded from the build-driving offering of the system registry. The registry records the exclusion verbatim: \texttt{extension\_\allowbreak{}hardware\_\allowbreak{}system\_\allowbreak{}axis.hpp} is marked as deprecated, \texttt{hardware\_\allowbreak{}isa\_\allowbreak{}system\_\allowbreak{}axis.hpp} a pure host descriptor that does not drive the build --- neither belongs to the build-driving offering. The static platform profiles with their fixed constants for cache-line size, page size and SIMD bit width are conceptually superseded by the run-time hardware detection of the \texttt{Cache\allowbreak{}Engine\allowbreak{}Builder} --- a factory over ISA and operating system that acquires hardware values instead of setting them statically (cf.\ Section~\ref{sec:hw-probe}). The four sub-axes HW1--HW4 and the three hardware profiles remain listed in full as an offering catalogue.

Defines the platform hardware family (CPU architecture, SIMD capabilities, cache topology, memory paging) as compile-time constants for all cache-engine algorithms. The axis captures which platforms are configurable at build time and provides static properties (cache-line size, page size, SIMD bit width, NUMA capability, Huge\allowbreak{}Page capability) for algorithm-appropriate data structures.

\subsubsection*{Static sub-axes}

The four sub-axes are orthogonal to one another; a platform concept assigns its primary characteristic to exactly one tag, and the \texttt{Cache\allowbreak{}Engine\allowbreak{}Builder} can group platforms by sub-axes (for instance HW1 $\times$ HW2 across all platforms).

\begin{scriptsize}
\begin{longtable}{@{}>{\raggedright\arraybackslash}p{1.8cm} >{\raggedright\arraybackslash}p{3.2cm} >{\raggedright\arraybackslash}p{8.5cm}@{}}
\toprule
ID & Name & Description \\
\midrule
\endhead%
HW1 & CPU Architecture Family (\texttt{cpu\_\allowbreak{}family}) & x86\_\allowbreak{}64 (Intel/\allowbreak{}AMD), aarch64 (ARM/\allowbreak{}Apple), riscv64 (RISC-V), generic (platform-independent). \\
HW2 & SIMD Instruction Set Concept (\texttt{simd\_\allowbreak{}capability}) & SSE/\allowbreak{}AVX/\allowbreak{}AVX2/\allowbreak{}AVX-512 (x86), NEON/\allowbreak{}SVE (ARM), RVV (RISC-V), scalar (no SIMD). \\
HW3 & NUMA and Cache Hierarchy (\texttt{memory\_\allowbreak{}topology}) & UMA-single-socket, NUMA-multi-socket, Apple-unified-memory, HBM-on-package (Sapphire Rapids HBM, Grace Hopper), CXL-attached. \\
HW4 & Page Configuration (\texttt{page\_\allowbreak{}topology}) & Standard and HugePage configuration: 4K-standard (POSIX), 16K-apple-silicon, 2M-hugepage (x86), 64K-arm-server, 1G-gigantic-hugepage. \\
\bottomrule
\end{longtable}
\end{scriptsize}

\subsubsection*{Building blocks}

\begin{scriptsize}
\begin{longtable}{@{}>{\raggedright\arraybackslash}p{3.2cm} >{\raggedright\arraybackslash}p{8.0cm} >{\raggedright\arraybackslash}p{2.3cm}@{}}
\toprule
Variant & Description & Source \\
\midrule
\endhead%
\texttt{X86\_\allowbreak{}64\allowbreak{}Hardware\allowbreak{}Profile} & Modern x86-64 desktop/\allowbreak{}server platform (Intel/\allowbreak{}AMD). Build-time constants: \texttt{cache\_\allowbreak{}line\_\allowbreak{}size}=64, \texttt{memory\_\allowbreak{}page\_\allowbreak{}size}=4096, \texttt{simd\_\allowbreak{}width\_\allowbreak{}bits}=256 (AVX2 conservative), \texttt{numa\_\allowbreak{}capable}=true, \texttt{huge\_\allowbreak{}page\_\allowbreak{}capable}=true. Basis for desktop- and server-side experiments. & Code/\allowbreak{}Baseline \\
\texttt{Aarch64\allowbreak{}Hardware\allowbreak{}Profile} & ARM64/\allowbreak{}AArch64 server and mobile platforms (Ampere Altra, Apple Silicon M-series, Linux ARM). Build-time constants: \texttt{cache\_\allowbreak{}line\_\allowbreak{}size}=64, \texttt{memory\_\allowbreak{}page\_\allowbreak{}size}=4096, \texttt{simd\_\allowbreak{}width\_\allowbreak{}bits}=128 (NEON standard), \texttt{numa\_\allowbreak{}capable}=true, \texttt{huge\_\allowbreak{}page\_\allowbreak{}capable}=true. Enables cross-platform comparisons. & Code/\allowbreak{}Baseline \\
\texttt{Generic\allowbreak{}Hardware\allowbreak{}Profile} & Conservative fallback profile for platform-independent configurations. Build-time constants: \texttt{cache\_\allowbreak{}line\_\allowbreak{}size}=64, \texttt{memory\_\allowbreak{}page\_\allowbreak{}size}=4096, \texttt{simd\_\allowbreak{}width\_\allowbreak{}bits}=0 (scalar-only), \texttt{numa\_\allowbreak{}capable}=false, \texttt{huge\_\allowbreak{}page\_\allowbreak{}capable}=false. Lowest common denominator when no specialized adapter matches. & Code/\allowbreak{}Baseline \\
\bottomrule
\end{longtable}
\end{scriptsize}

\emph{Sources:} platform profiles as Code/\allowbreak{}Baseline (no external publications).

